\documentclass[10pt,a4paper]{article}

% ----------------- Paquetes -------------------%
\usepackage[T1]{fontenc}
\usepackage[utf8]{inputenc}
\usepackage[a4paper,margin=2.5cm]{geometry}
\usepackage{mathptmx}             
\usepackage{changepage} 
\usepackage{setspace}
\usepackage[spanish]{babel}
\usepackage[labelfont=bf,labelsep=space]{caption}
\usepackage{amssymb}
\usepackage{amsmath}
\usepackage{ebgaramond}
\usepackage{placeins}
\usepackage{etoolbox}
\usepackage{setspace}
\usepackage{parskip} 
\usepackage{tcolorbox}
\usepackage{needspace}
\usepackage{xparse}
\usepackage{ocgx2}
\usepackage{hyperref}
\usepackage{hypcap}


%%%%%%%%%%%%%%%%%%%%%%%%%%%%%%%%%%%%%%%%%%%%%%%%%%%%%%%%%%%%%%%%
% --- Fija primero tus valores globales "buenos" ---
\setstretch{1.2}                 % Interlineado global
\setlength{\parskip}{0.7\baselineskip}
\setlength{\parindent}{0pt}
\newlength{\GlobalParskip}
\setlength{\GlobalParskip}{\parskip}

\newcounter{eqnum}[subsection]
\renewcommand{\theeqnum}{\arabic{eqnum}}
\newcounter{alphaitem}
\newcounter{numitem}

%% ------------------- Recuadros----------------------%%
\newcommand{\puntosclave}[1]{%
  \begin{tcolorbox}[
    colframe=black,
    title={Puntos clave},
    before upper={\setlength{\parindent}{0.5cm}\setlength{\parskip}{0.5\baselineskip}}
  ]
  #1
  \end{tcolorbox}
}

\newcommand{\modificaciones}[1]{
  \begin{tcolorbox}[
    colback=white,
    colframe=red,
    title={Modificaciones a realizar},
    before upper={\setlength{\parindent}{0.5cm}\setlength{\parskip}{0.5\baselineskip}}
  ]
  #1
  \end{tcolorbox}
}

%% ------------------- Preguntas TEST----------------------%%
% Opciones: radio (single-choice)
\NewDocumentCommand{\OptRadio}{m m m}{%
  \noindent\CheckBox[radio,name=#1,value=#2,width=1.2ex,height=1.2ex]{}%
  \;\textbf{#2.}~#3\par
}

% Opciones: checkbox (multi-choice)
\NewDocumentCommand{\OptCheck}{m m}{%
  \noindent\CheckBox[name=#1,width=1.2ex,height=1.2ex,bordercolor={0 0 0}]{}%
  \;#2\par
}

% Pregunta single-choice (radio)
% Args: 1=id, 2=enunciado, 3=opciones, 4=solución+justificación, 5=imagen (o {}), 6=origen
\NewDocumentCommand{\PreguntaTestSingle}{m m m m m m}{%
  \par\medskip
  \noindent\textbf{#1.}~#2\par\smallskip

    % Imagen (si no es {})
    \IfBlankTF{#5}{}{%
      \begin{center}
        \includegraphics[width=0.75\linewidth]{#5}
      \end{center}
    }

    \begin{Form}
      #3
    \end{Form}

    \par\smallskip
    \switchocg{sol-#1}{\fbox{Mostrar / ocultar solución}}%
    \hfill{\footnotesize #6}\par\smallskip

    \begin{ocg}{Solución}{sol-#1}{0}
      \noindent #4
    \end{ocg}
  \par\medskip
}

% Pregunta multi-choice (checkboxes)
\NewDocumentCommand{\PreguntaTestMulti}{m m m m m m}{%
  \par\medskip
  \noindent\textbf{#1.}~#2\par\smallskip

    % Imagen (si no es {})
    \IfBlankTF{#5}{}{%
      \begin{center}
        \includegraphics[width=0.75\linewidth]{#5}
      \end{center}
    }
    \begin{Form}
      #3
    \end{Form}

    \par\smallskip
    \switchocg{sol-#1}{\fbox{Mostrar / ocultar solución}}%
    \hfill{\footnotesize #6}\par\smallskip

    \begin{ocg}{Solución}{sol-#1}{0}
      \noindent #4
    \end{ocg}
  \par\medskip
}

%% ------------------- Listas----------------------%%
    % --- Lista alfabética ---
    \newenvironment{la}
      {\begin{list}{\alph{alphaitem})}{%
          \usecounter{alphaitem}%
          \setlength{\leftmargin}{1cm}%
          \setlength{\parskip}{3pt}% 
          \setlength{\itemsep}{0pt}% ← espacio entre ítems
          \setlength{\parsep}{0pt}%  ← párrafos dentro de un ítem
      }}
      {\end{list}}
      
    % --- Lista numérica ---
    \newenvironment{lnum}
      {\begin{list}{\arabic{numitem}.}{%
          \usecounter{numitem}%
          \setlength{\leftmargin}{1cm}%
          \setlength{\parskip}{3pt}% 
          \setlength{\itemsep}{0pt}% ← espacio entre ítems
          \setlength{\parsep}{0pt}%  ← párrafos dentro de un ítem
      }}
      {\end{list}}
      
% ------------------- Imágenes----------------------%
    \graphicspath{{Img/}}% Rutas por defecto 
    % --- Imagen adaptativa ---
    % Registros de dimensión definidos una sola vez
    \newdimen\imgcap
    \newdimen\imgmin
    \newdimen\imgmax
    \newdimen\imgremain
    \newdimen\imgh
    \newdimen\imgneed
    
    \newcommand{\imagenfit}[3]{%
      \addvspace{10pt}% separación antes
      \par\begingroup
        % Parámetros de composición (asignaciones, no \newdimen)
        \imgcap=1\baselineskip            % alto estimado de título+fuente
        \imgmin=0.15\textheight
        \imgmax=0.15\textheight
        \imgremain=\pagegoal
        \ifdim\pagegoal=\maxdimen \imgremain=0pt\fi
        \advance\imgremain by -\pagetotal
        \advance\imgremain by -\imgcap
        % Decisión de altura destino
        \ifdim\imgremain<\imgmin
          \newpage
          \imgh=0.20\textheight
        \else
          \imgh=\imgremain
          \ifdim\imgh>\imgmax \imgh=\imgmax \fi
        \fi
        % Reservar espacio para evitar cortes del bloque completo
        \imgneed=\imgh \advance\imgneed by \imgcap
        \Needspace{\imgneed}
        % Cuerpo
        \noindent\begin{minipage}{\textwidth}
          \centering
          \captionof{figure}{\textemdash\ #2}\par\vspace{0.3em}% título arriba
          \includegraphics[width=\linewidth,height=\imgh,keepaspectratio]{\detokenize{#1}}\par
          \vspace{0.3em}\small\textit{Fuente: }#3% fuente abajo
        \end{minipage}%
      \endgroup
      \par\addvspace{10pt}%
    }

%------------ Bloque de RECUADRO ESPECIAL -------------%
    \newenvironment{specialblock}[1]{%
      \begin{quote} % crea sangría a izquierda y derecha
      \small % texto más pequeño
      \noindent\textbf{#1}\\[-0.3em] % título en negrita
      \rule{\linewidth}{0.4pt}\\[0.6em]}{
      \end{quote}}
      
%-------------------------- Portada -----------------------%
\newcommand{\oldtitle}[1]{\def\OldTitle{#1}}
\newcommand{\changerationale}[1]{\def\ChangeWhy{#1}}

\newcommand{\changebox}{%
\begin{tcolorbox}[title={Historial de cambios del tema}]
\textbf{Título anterior:} \OldTitle\par
\textbf{Motivo del cambio:} \ChangeWhy
\end{tcolorbox}
}

%-------------------------- Author Font -----------------------%
    \newcommand{\authorfont}[1]{{\fontfamily{EBGaramond-TLF}\selectfont #1}}

%-------------------------- Anexos -----------------------%
    \makeatletter
    \newcounter{anexo}
    \renewcommand{\theanexo}{\Roman{anexo}}
    \newcommand{\anexo}[1]{%
      \clearpage
      \phantomsection
      \refstepcounter{anexo}%
      \noindent\textbf{Anexo \theanexo.- }#1\par\medskip
      \addcontentsline{stc}{section}{Anexo \theanexo.- #1}%
    }
    \makeatother

    %-------------------------- Lista de figuras (personal) -----------------------%
\makeatletter
\newcommand{\MyListOfFigures}{%
  \clearpage
  \phantomsection
  {\centering\bfseries\Large Índice de figuras\par}%
  \medskip
  % Entrada en nuestro índice de contenido (stc)
  \addcontentsline{stc}{section}{Índice de figuras}%
  % Recuperación del listado (sin cabecera propia)
  \begingroup
    \parindent=0pt\parskip=2pt
    \@starttoc{lof}%
  \endgroup
}
\makeatother

%-------------------------- Bibliografía -----------------------%
\makeatletter
% Contador para las referencias
\newcounter{myref}
% Cabecera de la bibliografía, entra en el índice 'stc'
\newcommand{\MyBibliography}{%
  \clearpage
  \phantomsection
  {\centering\bfseries\Large Bibliografía y referencias\par}%
  \medskip
  \addcontentsline{stc}{section}{Bibliografía y referencias}%
  \setcounter{myref}{0}%
}
\newcommand{\MyBibItem}[4]{%
  \refstepcounter{myref}%
  \noindent[\themyref]\ #1.\ \emph{#2}. #3. #4\par
}
\makeatother

%--------- Establecer cuatro niveles de índice --------%
    \setcounter{tocdepth}{4}   
    \setcounter{secnumdepth}{4}% numera hasta \paragraph

%%%%%%%%%%%%%%%%%%%%%%%%%%%%%%%%

\makeatletter

    %--------- Interlineado Global --------%
    \edef\GlobalBaseLineStretch{\baselinestretch}
    
    %--------- Espaciado de seccinones --------%
    \renewcommand\section{\@startsection{section}
      {1}{\z@}
      {2.5ex \@plus 1ex \@minus .2ex} % espacio antes
      {1.5ex \@plus .2ex}             % espacio después
      {\normalfont\Large\bfseries}    % formato del título
    }
    \renewcommand\subsection{\@startsection{subsection}{2}{\z@}
      {3ex \@plus 0.8ex \@minus .2ex}
      {1.5ex \@plus .2ex}
      {\normalfont\large\bfseries}
    }
    \renewcommand\subsubsection{\@startsection{subsubsection}{3}{\z@}
      {3ex \@plus 0.5ex \@minus .2ex}
      {1ex \@plus .2ex}
      {\normalfont\normalsize\bfseries}
    }
    \renewcommand\paragraph{\@startsection{paragraph}{4}{\z@}%
    {-2ex \@plus -0.5ex \@minus -.2ex}% espacio antes
    {0.8ex \@plus .2ex}% espacio después
    {\normalfont\normalsize\itshape}}% ← cursiva en lugar de negrita

    %--------- Ecuaciones --------%   
    % Variables globales para almacenar títulos y notas
    \gdef\leq@current@section{}%
    \gdef\leq@current@subsection{}%
    \gdef\leq@last@printed@section{}%
    \gdef\leq@last@printed@subsection{}%
    
    % Comando para añadir nota (invisible en el cuerpo)
    \newcommand{\eqnote}[1]{%
      \addcontentsline{leq}{leqnote}{#1}%
    }
    
    % Comando para ecuaciones
    \newcommand{\eqblock}[2]{%
      % Verificar si necesitamos imprimir el título de sección
      \ifx\leq@current@section\leq@last@printed@section\else
        \addcontentsline{leq}{leqsection}{\leq@current@section}%
        \global\let\leq@last@printed@section\leq@current@section
        \gdef\leq@last@printed@subsection{}% Reset subsección al cambiar sección
      \fi
      % Verificar si necesitamos imprimir el título de subsección
      \ifx\leq@current@subsection\leq@last@printed@subsection\else
        \ifx\leq@current@subsection\@empty\else
          \addcontentsline{leq}{leqsubsection}{\leq@current@subsection}%
          \global\let\leq@last@printed@subsection\leq@current@subsection
        \fi
      \fi
      % Ecuación
      \refstepcounter{eqnum}%
      \begin{equation*}
        #1\tag{E.\theeqnum}
      \end{equation*}%
      \addcontentsline{leq}{leqt}{[E.\theeqnum] -- #2}%
      \addcontentsline{leq}{leqm}{#1}%
    }
    
    %%% ----- Índice de ecuaciones ----------%%%
    % Tamaño para []
    \newcommand{\leqIndexFont}{\normalsize}
    \newcommand{\leqTitleFont}{\tiny}
    \newcommand{\leqNoteFont}{\tiny}
    % Cabecera de sección 
    \newcommand*\l@leqsection[2]{%
      \addvspace{25pt}% Mayor separación antes de nueva sección
      \noindent\leqIndexFont
      \makebox[\linewidth]{\rule{.38\linewidth}{0.4pt}\ \ \textbf{  \MakeUppercase{#1}}\ \ \rule{.38\linewidth}{0.4pt}}%
      \par\vspace{-10pt}
    }
    % Cabecera de subsección 
    \newcommand*\l@leqsubsection[2]{%
      \par\addvspace{15pt}%
      \noindent\leqIndexFont #1
      \par\vspace{-7pt}
    }
    % Título de ecuación 
    \newcommand*\l@leqt[2]{%
    \par\addvspace{10pt}%
      \par\noindent\leqTitleFont #1\par
    }
    % Miniatura de ecuación
    \newcommand*\l@leqm[2]{%
      \vspace{0.3\baselineskip}%
      \par\scriptsize\noindent\hspace*{1.5em}\(#1\)\par%
    }
    % Nota de ecuación 
    \newcommand*\l@leqnote[2]{%
      \vspace{0.2\baselineskip}%
      \par\noindent\leqNoteFont\hspace*{1.5em}\textbf{Nota.--} #1\par\vspace{0.5\baselineskip}%
    }
    % Generación del índice
    \newcommand{\mylistofequations}{%
      \clearpage
      \phantomsection
      % Título visual de la lista (ajústalo a tu gusto):
      {\centering\bfseries\Large Índice de ecuaciones\par}
      % Entrada en nuestro índice 'stc':
      \addcontentsline{stc}{section}{Índice de ecuaciones}%
      % Llamada a tu lista real (usa el nombre exacto de tu macro):
      \@starttoc{leq}%
    }
    
    %--------- Captura de títulos de sección --------%
    \let\leq@original@section\section
    \let\leq@original@subsection\subsection
    % Flag para saber si estamos en una sección asterisco
    \newif\ifLeqStarSection
    \LeqStarSectionfalse
    
    % Redefinir \section CON CONTROL DE ASTERISCO
    \renewcommand{\section}{%
      \@ifstar{\leq@section@star}{\@dblarg\leq@section@nostar}%
    }
    \def\leq@section@star#1{%
      \global\LeqStarSectiontrue
      \gdef\leq@current@section{#1}%
      \gdef\leq@current@subsection{}%
      \leq@original@section*{#1}%
      \global\LeqStarSectionfalse
    }
    \def\leq@section@nostar[#1]#2{%
      \global\LeqStarSectionfalse
      \gdef\leq@current@section{#2}%
      \gdef\leq@current@subsection{}%
      \leq@original@section[#1]{#2}%
    }
    % Redefinir \subsection
    \renewcommand{\subsection}{%
      \@ifstar{\leq@subsection@star}{\@dblarg\leq@subsection@nostar}%
    }
    \def\leq@subsection@star#1{%
      \global\LeqStarSectiontrue
      \gdef\leq@current@subsection{#1}%
      \leq@original@subsection*{#1}%
      \global\LeqStarSectionfalse
    }
    \def\leq@subsection@nostar[#1]#2{%
      \global\LeqStarSectionfalse
      \gdef\leq@current@subsection{#2}%
      \leq@original@subsection[#1]{#2}%
    }

    % ----- Restauración de valores Globales de estilo ----------%
    % Flags
    \newif\ifInFootnote
    \InFootnotefalse
    \pretocmd{\@footnotetext}{\InFootnotetrue}{}{}
    \apptocmd{\@footnotetext}{\InFootnotefalse}{}{}
    % Profundidad de listas (soporta anidadas)
    \newcount\MyListDepth \MyListDepth=0
    \AtBeginEnvironment{la}{\advance\MyListDepth by 1}
    \AfterEndEnvironment{la}{\advance\MyListDepth by -1}
    \AtBeginEnvironment{lnum}{\advance\MyListDepth by 1}
    \AfterEndEnvironment{lnum}{\advance\MyListDepth by -1}
    \AtBeginEnvironment{itemize}{\advance\MyListDepth by 1}
    \AfterEndEnvironment{itemize}{\advance\MyListDepth by -1}
    \AtBeginEnvironment{enumerate}{\advance\MyListDepth by 1}
    \AfterEndEnvironment{enumerate}{\advance\MyListDepth by -1}
    % Restauración diferida: hazlo al INICIO del próximo párrafo real
    \newif\if@pendingRestore
    \@pendingRestorefalse
    \newcommand{\RequestRestore}{\global\@pendingRestoretrue}
    \everypar{%
      \if@pendingRestore
        \global\@pendingRestorefalse
        \ifInFootnote\else
          \ifnum\MyListDepth=0
            \setlength{\parskip}{\GlobalParskip}%
            \setstretch{\GlobalBaseLineStretch}%
          \fi
        \fi
      \fi
    }
    % Al final de bloques:
    \AfterEndEnvironment{equation}{\RequestRestore}
    \AfterEndEnvironment{equation*}{\RequestRestore}
    \AfterEndEnvironment{la}{\RequestRestore}
    \AfterEndEnvironment{lnum}{\RequestRestore}
    % Al final de macros:
    \apptocmd{\eqblock}{\RequestRestore}{}{}
    \apptocmd{\imagenfit}{\RequestRestore}{}{}
    
\makeatother

% ===================== ToC propio con 4 niveles (único bloque) =====================
\makeatletter

% --- Escritores: añadimos una línea al .stc en cada cabecera numerada ---
% Guardamos las versiones actuales para llamar al comportamiento normal
\let\MyTOC@orig@section\section
\let\MyTOC@orig@subsection\subsection
\let\MyTOC@orig@subsubsection\subsubsection
\let\MyTOC@orig@paragraph\paragraph

% SECTION
\renewcommand{\section}{%
  \@ifstar{\MyTOC@section@star}{\@dblarg\MyTOC@section@nostar}%
}
\def\MyTOC@section@star#1{%
  \MyTOC@orig@section*{#1}% sin entrada al .stc
}
\def\MyTOC@section@nostar[#1]#2{%
  \MyTOC@orig@section[{#1}]{#2}% comportamiento normal (escribe al .toc)
  % Escribimos también nuestra línea al .stc (texto con \numberline)
  \addcontentsline{stc}{section}{\protect\numberline{\thesection}#2}%
}

% SUBSECTION
\renewcommand{\subsection}{%
  \@ifstar{\MyTOC@subsection@star}{\@dblarg\MyTOC@subsection@nostar}%
}
\def\MyTOC@subsection@star#1{%
  \MyTOC@orig@subsection*{#1}%
}
\def\MyTOC@subsection@nostar[#1]#2{%
  \MyTOC@orig@subsection[{#1}]{#2}%
  \addcontentsline{stc}{subsection}{\protect\numberline{\thesubsection}#2}%
}

% SUBSUBSECTION
\renewcommand{\subsubsection}{%
  \@ifstar{\MyTOC@subsubsection@star}{\@dblarg\MyTOC@subsubsection@nostar}%
}
\def\MyTOC@subsubsection@star#1{%
  \MyTOC@orig@subsubsection*{#1}%
}
\def\MyTOC@subsubsection@nostar[#1]#2{%
  \MyTOC@orig@subsubsection[{#1}]{#2}%
  \addcontentsline{stc}{subsubsection}{\protect\numberline{\thesubsubsection}#2}%
}

% PARAGRAPH
\renewcommand{\paragraph}{%
  \@ifstar{\MyTOC@paragraph@star}{\@dblarg\MyTOC@paragraph@nostar}%
}
\def\MyTOC@paragraph@star#1{%
  \MyTOC@orig@paragraph*{#1}%
}
\def\MyTOC@paragraph@nostar[#1]#2{%
  \MyTOC@orig@paragraph[{#1}]{#2}%
  % Si no numeras los \paragraph, cambia \theparagraph por texto plano sin \numberline
  \addcontentsline{stc}{paragraph}{\protect\numberline{\theparagraph}#2}%
}

% --- Impresor del índice propio (lee el .stc) ---
\newcommand{\MyTOC}{%
  \par\bigskip
  {\centering\bfseries\Large Índice de contenido\par}%
  \medskip
  \begingroup
    \parindent=0pt\parskip=2pt
    \@starttoc{stc}%
  \endgroup
}

\makeatother
% ===================== fin ToC propio =====================


%%%%%%%%%%%%%%


% --- Datos del tema ---
\title{Teoría de la integración monetaria\\
\large }
\author{Marcelo Ortega}
\date{Marzo 2026}
\newcommand{\TemaCode}{3B16}

\oldtitle{
    B.15. Teoría de la integración monetaria.
}
\changerationale{
    Sin cambios.
}

\begin{document}


\maketitle
\changebox

\newpage

% --- Página de resumen y modificaciones ---
\puntosclave{ %% Escribir puntos clave Apuntes CECO
     
    }
\vspace{1cm}

\modificaciones{
    
    }

\newpage
\MyTOC
\newpage

\mylistofequations
\clearpage

\MyListOfFigures
\clearpage


\pagenumbering{arabic}
\setcounter{page}{1}


\section{Introducción}



\subsection{Contextualización}


\subsubsection{Relevancia}




\subsubsection{Problemática}



\subsection{Estructura de la exposición}
La estructura consiste en dos bloques de teoría económica y uno de evidencia empírica. Los dos primeros están separados, ya que responden a las dos diferentes olas de literatura sobre la materia. Mientras que la primera se enmarca en los años 50 del pasado siglo, y estudia tantocasos de países concretos (por ejemplo, ¿es EEUU una unión monetaria óptima?) como de hipotéticas uniones monetarias internacionales (¿formarían EEUU y Canadá una unión monetaria óptima?), la segunda ola de la literatura está fuertemente ligada a la creación y primeros desarrollos de la Eurozona (es decir, es contemporánea primero al Plan Delors, y después de la crisis del euro). Como resultado de lo anterior, mientras que la primera ola hace una suerte de estudio "costebeneficio" acerca de la deseabilidad de fonnar una unión monetaria o no, la segunda se centra en aspectos concretos que son relevantes para que una unión monetaria funcione correctamente (principalmente, el alineamiento de los ciclos, los mecanismos de integración fiscal en materia de transferencias fiscales, mutualización de riesgos y coordinación de las políticas fiscales, el diseño institucional, la integración financiera o los ajustes estructurales). Entre ambas olas se comprueban además los avances de la literatura macroeconómica en general: desde modelos basados en el lS-LM-BP-OA de autores como Mundell, hasta la microfundamentación actual.

En lo referente al último apartado, se incluyen dos estudios de caso: el de la Unión Económica y Monetaria de la UE y el hipotético caso de una unión monetaria en el Mercosur. Ésta es ocasión de aterrizar los desarrollos teóricos vistos anterionnente. Si bien este terna es indudablemente útil para entender correctamente el 3843, hay que tener cuidado de no entrar a hacer una exposición exhaustiva del diseño institucional de la UEM en este tema.



\clearpage
\section{}
\puntosclave{
    En una unión monetaria se pueden producir choques asimétricos que pueden ser costosos para los países que los sufran al no poder utilizar el tipo de cambio como instrumento de ajuste. 
    
    Si dentro de una unión monetaria existe libertad de movimiento de trabajadores, esto puede ayudar al ajuste en caso de choque asimétrico. Otra posibilidad es la flexibilidad de salarios y precios, pero es una opción más dolorosa. 

    Una unión fiscal que acompañe a la Unión Monetaria puede amortiguar el impacto de los choques asimétricos.
}


\subsection{Costes de una Unión Monetaria}
Los países que se incorporan a una unión monetaria enfrentan costes importantes. Estos costes surgen del hecho de que los países difieren en muchos aspectos y que, cuando un país renuncia a su moneda nacional, también renuncia a un instrumento muy importante de política económica que podría servirle otrorí como instrumento estabilizador. En otras palabras, en una unión monetaria plena el banco central nacional o bien deja de existir o bien no tendrá poder real, el Estado que se una: no podrá modificar el precio de su moneda (mediante devaluaciones y revaluaciones), determinar la cantidad de dinero nacional en circulación, ni modificar el tipo de interés a corto plazo.

Los países pueden utilizar políticas monetarias nacionales, incluidos cambios en el tipo de cambio, para corregir estas diferencias. Encontramos que, en la mayoría de los casos, existen alternativas al uso de la política monetaria nacional como instrumento. Por ejemplo, cuando un país enfrenta una pérdida de competitividad, puede intentar recuperarla mediante reducciones de salarios y precios. Sin embargo, estas alternativas suelen ser más costosas para un miembro de una unión monetaria que para un país independiente que ha mantenido su soberanía monetaria (incluida la capacidad de modificar el tipo de cambio).

Para saber por qué los países pueden encontrar costoso renunciar a sus monedas nacionales y unirse a una unión monetaria, especialmente cuando dicha unión es incompleta, es necesario observar algunos de los trade offs más importantes que enfrenta cualquier proceso de integración monetaria.\footnote{El análisis que sigue en este capítulo se conoce como la Teoría de las Áreas Monetarias Óptimas. Esta teoría, iniciada por Mundell (1961), McKinnon (1963) y Kenen (1969), se ha centrado en el lado de los costes dentro del análisis coste-beneficio de una unión monetaria.
}

\textcolor{red}{\textbf{OJO->} El análisis presentado en este capítulo, basado en la teoría de las áreas monetarias óptimas, ha sido objeto de numerosas críticas. Estas críticas han dado lugar a nuevos e importantes desarrollos. En el Capítulo 2 se abordarán estas críticas.}

\subsubsection{Shocks asimétricos: MUNDELL (1961)}
Considérese el caso de un desplazamiento de la demanda desarrollado por MUNDELL (1961). Supongamos que dos países, a los que llamamos Francia y Alemania, forman una unión monetaria y utilizan una moneda común gestionada por un banco central común, el Banco Central Europeo. Supongamos además que, por alguna razón, los consumidores desplazan sus preferencias desde los productos fabricados en Francia hacia los productos fabricados en Alemania:\footnote{Este es el efecto sustitución de un aumento de precios. En el análisis estándar de demanda agregada, existe también un efecto monetario: cuando el nivel de precios doméstico aumenta, el stock de saldos reales de dinero disminuye, lo que conduce a un aumento del tipo de interés real interno. Esto, a su vez, reduce la demanda agregada (véase De Grauwe, 1983). Aquí ignoramos el efecto monetario y nos centramos en el efecto sustitución.

La curva de oferta expresa la idea de que, cuando el precio de la producción nacional aumenta, las empresas domésticas, en un entorno competitivo, incrementarán su oferta para beneficiarse del mayor precio. Además, cada curva de oferta se traza bajo el supuesto de que el salario nominal y los precios de otros insumos (por ejemplo, energía, insumos importados) permanecen constantes. Cambios en los precios de estos insumos desplazarán estas curvas de oferta.
}

\imagenfit{Fig1.png}{Demanda y Oferta Agregadas en Francia y Alemania: Shocks de demanda}{}

Representando las curvas estándar de demanda y oferta agregadas en una economía abierta, vemos como el desplazamiento de la demanda se representa mediante un movimiento hacia arriba de la curva de demanda en Alemania y un movimiento hacia abajo en Francia.\footnote{Como se analizará más adelante, será importante saber si estos desplazamientos de la demanda son permanentes o temporales. Por el momento, suponemos que estos cambios son permanentes, por ejemplo, debido a un cambio en las preferencias de los consumidores.
} El resultado de estos desplazamientos de la demanda es que la producción disminuye en Francia y aumenta en Alemania, lo que probablemente conducirá a un aumento del desempleo en Francia y a una disminución del desempleo en Alemania. En última instancia, ambos países tendrán un problema de ajuste: (i) Francia se ve afectada por una reducción de la producción y un mayor desempleo; y, (ii) Alemania experimenta un auge, que también genera presiones al alza sobre su nivel de precios. 

¿Existe un mecanismo que conduzca a un restablecimiento automático del equilibrio? No uno, sino dos mecanismos devolverían automáticamente el equilibrio en ambos países:
\begin{la}
    \item \textbf{Flexibilidad salarial}. Si los salarios en Francia y Alemania son flexibles, los trabajadores franceses (que están desempleados) reducirán sus reivindicaciones salariales, mientras que en Alemania el exceso de demanda de trabajo empujará al alza el nivel salarial.\footnote{La reducción del salario en Francia desplazaría la curva de oferta agregada hacia abajo, mientras que el aumento de los salarios en Alemania desplaza la curva de oferta agregada hacia arriba. Estos desplazamientos conducen a un nuevo equilibrio. En Francia, el precio de la producción disminuye, los productos franceses se vuelven más competitivos y se estimula la demanda. Lo contrario ocurre en Alemania.
    } No obstante, deben hacerse dos observaciones sobre este mecanismo de ajuste: (i) en primer lugar, los aumentos salariales y de precios en Alemania hacen que los productos franceses sean más competitivos, lo que conduce a un desplazamiento hacia arriba de la curva de demanda agregada francesa, y viceversa;\footnote{
    Es decir, la disminución de costes y precios en Francia hace que los productos alemanes sean menos competitivos y desplaza la curva de demanda agregada alemana hacia abajo. 
    
    Estos efectos se conocen como de \textbf{segundo orden} sobre la demanda agregada que \textbf{refuerzan el mecanismo de equilibrio}.
    } y, (ii) en segundo lugar, una reducción salarial en Francia implica que los trabajadores franceses verán disminuir su renta real, lo que (probablemente) reducirá la demanda de bienes y servicios;\footnote{Como resultado de este efecto renta de la caída salarial, el efecto positivo sobre la competitividad puede verse compensado por nuevos desplazamientos descendentes de la curva de demanda.
    } y,
    \item \textbf{Movilidad del trabajo}. Un segundo mecanismo que conduce a un nuevo equilibrio es la movilidad del trabajo. Los trabajadores franceses desempleados se trasladan a Alemania, donde existe exceso de demanda de trabajo. Este movimiento de trabajo elimina la necesidad de que los salarios disminuyan en Francia y aumenten en Alemania. Así, el problema de desempleo en Francia desaparece, mientras que las presiones salariales inflacionarias en Alemania se desvanecen.
\end{la}


En cualquier caso, el restablecimiento del equilibrio puede visualizarse de la siguiente forma:

\imagenfit{Fig1_2.png}{Proceso de ajuste automático}{}

Por tanto, el problema de ajuste para Francia y Alemania desaparecerá automáticamente si los salarios son flexibles y/o si la movilidad del trabajo entre ambos países es suficientemente elevada (en principio). 

Sin embargo, si estas condiciones no se cumplen, el problema de ajuste no desaparecerá. Supongamos, por ejemplo, que los salarios en Francia no disminuyen a pesar del  desempleo, y que los trabajadores franceses no se trasladan a Alemania. En ese caso, Francia queda atrapada en la situación de desequilibrio; mientras que en Alemania, el exceso de demanda de trabajo ejerce presión al alza sobre los salarios, produciendo un desplazamiento hacia arriba de la curva de oferta. El ajuste del desequilibrio deberá entonces producirse exclusivamente a través de aumentos de precios en Alemania. Estos aumentos de precios alemanes vuelven a hacer más competitivos los bienes franceses, lo que conduce a un desplazamiento hacia arriba de la curva de demanda agregada en Francia. Así, \textbf{si los salarios no disminuyen en Francia, el ajuste del desequilibrio adoptará la forma de inflación en Alemania}.

¿Qué habría ocurrido si los dos países no hubieran formado una unión monetaria? En ese caso, habrían sido libres de utilizar sus instrumentos nacionales de política monetaria para ajustarse a shocks asimétricos. Existen varias formas en las que los países que mantienen su independencia monetaria pueden utilizar estos instrumentos. Entre los más destacados:
\begin{la}
    \item \textbf{Tipo de Cambio Flexible}. Supongamos primero que Francia y Alemania hubieran elegido un régimen de tipo de cambio flexible. Francia podría haber reducido su tipo de interés, estimulando así la demanda agregada, mientras que Alemania podría haberlo aumentado, reduciendo su demanda agregada. Estas políticas monetarias probablemente habrían conducido a una depreciación del franco francés y a una apreciación del marco alemán, haciendo los productos franceses más baratos en Alemania. Tanto el cambio en los tipos de interés como en el tipo de cambio tenderían a aumentar la demanda agregada en Francia y a reducirla en Alemania;\footnote{En un primer régimen, los países mantienen tipos de cambio flexibles, como hacen Estados Unidos, el Reino Unido y Japón. En ese caso, pueden modificar sus políticas monetarias (mediante cambios en el tipo de interés interno y/o en la oferta monetaria) para alcanzar determinados objetivos. En un segundo régimen, los países fijan su tipo de cambio a otra moneda, por ejemplo Dinamarca al euro o varios países latinoamericanos al dólar. En este caso, pueden devaluar o revaluar sus monedas.
    } y,
    \item \textbf{Tipo de Cambio Rígido}. Si, por el contrario, Francia y Alemania hubieran optado por un tipo de cambio fijo, Francia habría podido devaluar el franco frente al marco, logrando efectos similares sobre la demanda agregada. La devaluación del franco habría incrementado la competitividad de los productos franceses, estimulando la demanda procedente de Alemania.
\end{la}

Los efectos de estas políticas monetarias nacionales pueden visualizarse como:

\imagenfit{Fig1_3.png}{Efectos de una expansió monetaria en Francia y una reducción en Alemania}{}

La política monetaria expansiva en Francia (o, en el segundo régimen, la devaluación del franco francés) desplaza la curva de demanda agregada francesa hacia arriba. En Alemania ocurre lo contrario. La política monetaria restrictiva en Alemania (la apreciación del marco) reduce la demanda agregada, desplazando la curva de demanda hacia la izquierda.

En cualquier caso estos desplazamientos resuelven el problema de desempleo en Francia y Alemania evitaría presiones inflacionarias, que se logra de forma notable utilizando un solo instrumento.

En contraste, cuando Francia forma parte de una unión monetaria con Alemania, renuncia al control de su política monetaria. Si sufre un problema persistente de desempleo, este solo puede desaparecer mediante deflación en Francia. En este sentido, puede decirse que la unión monetaria tiene un coste para Francia cuando enfrenta un shock negativo de demanda. Del mismo modo, Alemania encontrará costoso formar parte de una unión monetaria con Francia, ya que tendrá que aceptar más inflación de la que desearía.

Recapitulando, si los salarios son rígidos y la movilidad laboral es limitada, los países que forman una unión monetaria tendrán mayores dificultades para ajustarse a shocks de demanda asimétricos que los países que mantienen su propia moneda y pueden devaluarla (o revaluarla). Por el contrario, en los países que conservan su moneda, las políticas monetarias nacionales, incluido el tipo de cambio, aportan flexibilidad a un sistema que de otro modo sería rígido. Dicho de otro modo, una unión monetaria entre dos o más países es óptima si se cumple al menos una de las siguientes condiciones: (i) \textbf{existe suficiente flexibilidad salarial}; o, (ii) \textbf{existe suficiente movilidad del trabajo}.

\paragraph{Crítica: }



\subsubsection{Instituciones del mercado laboral: BRUNO y SACHS (1985)}
Otra fuente de asimetría que puede obligar a los países miembros de una Unión Monetaria a enfrentar procesos de ajuste difíciles son las diferencias institucionales en los mercados laborales entre países (p.e. europeos).\footnote{Algunos mercados laborales están dominados por sindicatos altamente centralizados (por ejemplo, Alemania). En otros países, los sindicatos están descentralizados (por ejemplo, el Reino Unido).
} Pero, ¿qué consecuencias pueden tener estas diferencias en el ajuste por shocks? De manera general, existe consenso en que estas diferencias pueden generar costes significativos para una Unión Monetaria, principalmente porque pueden dar lugar a trayectorias divergentes de salarios y precios, incluso cuando los países enfrentan perturbaciones similares. 

Una de las teorías más influyentes al respecto fue desarrollada por BRUNO y SACHS (1985). La idea es sencilla: los shocks (p.e. de oferta), tienen efectos macroeconómicos muy distintos dependiendo del grado de centralización de la negociación salarial. Veamos por qué. Imaginemos que experimentamos un shock negativo de oferta por el encarecimiento del crudo, ¿cómo se articularan los conflictos de interés y las reivindicaciones salariales? 
\begin{la}
    \item \textbf{Mercado laboral corporativista}. Cuando la negociación salarial está centralizada (sistemas corporativistas), es probable que los sindicatos tengan en cuenta el efecto inflacionario de los aumentos salariales. Es decir, saben que aumentos salariales excesivos generarán más inflación, de modo que los salarios reales no aumentarán. Por tanto, es razonable suponer que no tienen incentivos a formular demandas salariales excesivas;
    \item \textbf{Mercado laboral descentralizado}. Por el contrario, con negociación salarial menos centralizada, los sindicatos individuales que negocian aumentos salariales saben que el impacto de estos aumentos sobre el nivel general de precios es pequeño, ya que representan solo una fracción de la fuerza laboral. Surge entonces un problema de free-riding: cada sindicato tiene interés en aumentar los salarios nominales de sus miembros, lo que conduce a un equilibrio no cooperativo con un nivel salarial nominal más alto que en el caso anterior.\footnote{La analogía con los espectadores en un estadio de fútbol es bien conocida. Cuando todos están sentados, cada espectador tiene incentivos a levantarse para ver mejor. El resultado es que todos terminan de pie, sin ver mejor y estando más incómodos. Una vez que están de pie, es igualmente difícil inducirlos a sentarse: quien lo haga primero verá peor si los demás no lo imitan.
    
    Esta historia de cooperación fue ampliada por CALMFORS y DRIFFILL (1988), quienes señalaron que la relación entre centralización de la negociación salarial y resultados no es lineal. En particular, cuando se avanza hacia una mayor descentralización, aparece otra externalidad relevante. En sistemas muy descentralizados (por ejemplo, negociación a nivel de empresa), las demandas salariales afectan directamente a la competitividad de la empresa y, por tanto, a las perspectivas de empleo de los trabajadores. Demandas salariales excesivas provocarán una fuerte reducción del empleo. Por ello, ante un shock de oferta, los sindicatos en sistemas altamente descentralizados pueden mostrar un grado considerable de moderación salarial. Este resultado lleva a la conclusión de que los países con niveles extremos de centralización —muy altos o muy bajos— están mejor preparados para afrontar shocks de oferta (como aumentos en el precio del petróleo) que aquellos con niveles intermedios de centralización. En estos países extremos se observa mayor moderación salarial, lo que se traduce en mejores resultados en términos de inflación y empleo (véase Calmfors y Driffill, 1988; para un análisis más detallado, Hancke, 2014).
    }
\end{la}

Con este sencillo planteamiento, podemos suponer muy posible que, ante un shock de oferta, ocurran ajustes muy distintos en función del tipo de instituciones del mercado laboral. Una negociación corporativista entiende que la pérdida de salario real causada por dicho shock no puede compensarse mediante aumentos del salario nominal; mientras que en una negociación desentralizada, ningún sindicato tiene incentivos a ser el primero en moderar salarios, ya que corre el riesgo de que los demás no lo sigan, reduciendo así el salario real de sus miembros. En definitiva, en sistemas descentralizados es estructuralmente más difícil lograr moderación salarial. 

¿Por qué importan estas diferencias? Un país puede experimentar aumentos de salarios y precios más rápidos que otros, incluso cuando el shock inicial es el mismo. En términos del modelo de dos países utilizado en el análisis anterior, la curva de oferta se desplaza hacia arriba en mayor medida en un país (descentralizado) que en el otro (corporativista):

\textbf{INTRODUCIR GRÁFICO DE AJUSTE}

Por tanto, los países con instituciones del mercado laboral muy diferentes pueden encontrar costoso formar una unión monetaria. Ante cada shock de oferta, los salarios y precios pueden reaccionar de forma distinta, dificultando la corrección de estas divergencias cuando el tipo de cambio está irrevocablemente fijado.

\paragraph{Crítica: }

Las diferencias en el funcionamiento de los mercados laborales entre países están bien documentadas. Sin embargo, estas diferencias se han acumulado a lo largo del tiempo, en parte debido a regímenes de política económica distintos. La cuestión es si la integración monetaria puede modificar sustancialmente el comportamiento de los sindicatos, de modo que las diferencias actuales puedan desaparecer.

Un ejemplo ayuda a ilustrarlo. Representando los mercados laborales de dos países candidatos a una unión monetaria, en base al modelo de McDonald y Solow (1981):\footnote{En el eje vertical se muestra el salario real y en el horizontal el nivel de empleo (N). Las curvas convexas representan las curvas de indiferencia del sindicato. Se supone que existe un único sindicato en cada país, que maximiza su utilidad en función del salario real y el empleo de sus afiliados. La recta decreciente representa la demanda de trabajo agregada. Para el sindicato, esta recta actúa como restricción, y elige el punto que maximiza su utilidad (puntos A y B en la figura).
}

\imagenfit{Fig2_3.png}{El Modelo SOLOW-McDONALD en dos países}{}

Una característica importante del modelo es que la curva de empleo incorpora la reacción de las autoridades.\footnote{Por tanto, si estas otorgan mayor peso al empleo que los sindicatos, reaccionarán ante reducciones del empleo mediante políticas expansivas (monetarias o fiscales, creación de empleo público, etc.). Si los sindicatos anticipan esta reacción, la restricción cambia: la curva de empleo se vuelve más inclinada, ya que aumentos salariales reducen menos el empleo total al ser compensados por la política pública.

Esta curva puede interpretarse como la función de reacción del gobierno. El sindicato actúa como líder tipo STACKELBERG.
} Como podemos ver, la curva del país $B$ es más inclinada reflejando una mayor disposición de las autoridades a acomodar los salarios mediante políticas expansivas. La unión monetaria modifica esta situación, ya que la política monetaria se centraliza, haciendo que ambos países enfrenten una reacción común de la autoridad monetaria, lo que tiende a homogeneizar los resultados salariales y de empleo.\footnote{No obstante, estas diferencias no desaparecen completamente. Los gobiernos nacionales conservan otros instrumentos (por ejemplo, empleo público financiado con deuda). Por tanto, aunque las diferencias se reducen, no se eliminan. Además, el análisis supone sindicatos centralizados, lo cual no es realista en todos los países. No está claro cómo la unión monetaria afecta estas diferencias institucionales.
}

En conclusión, las diferencias institucionales en los mercados laborales persistirán durante un tiempo considerable tras la introducción de una moneda común, pudiendo generar divergencias en salarios y empleo y problemas de ajuste.

\subsubsection{Innovaciones financieras y mercados de crédito}
Otra diferencia estructural que puede tener implicaciones importantes en los rendimientos de una Unión Monetaria, se encuentra en la heterogeneidad de sistemas legales. En el caso europeo se ve claramente como estas diferencias son profundas y, en ocasiones, tienen efectos significativos sobre el funcionamiento de los mercados. Algunos ejemplos claros son:
\begin{lnum}
    \item \textbf{Mercados hipotecarios}. Los mercados hipotecarios operan de forma muy distinta en los países de la UE y, en gran medida, el resultado es que las hipotecas son productos muy diferentes, con distintos niveles de riesgo, según el país.\footnote{La ley protege a los bancos que conceden préstamos hipotecarios en mayor medida en algunos países que en otros. Un ejemplo claro se encuentra en la ratio préstamo-valor de los créditos con garantía hipotecaria. A priori, puede considerarse que el valor del préstamo suele ser inferior al valor de la vivienda (la garantía), por lo que la ratio préstamo-valor es inferior al 100\%. Sin embargo, si se entra en el funcionamiento del mercado en cada país, esta ratio varía considerablemente. Por otro lado, las diferencias legales también influyen en la frecuencia con la que se ajustan los tipos de interés. Así, hay países donde los bancos ofrecen hipotecas con tipos variables, mientras que en otros los tipos son fijos durante toda la vida del préstamo.

    De hecho, diversos estudios empíricos confirman que estas diferencias en la transmisión de los shocks pueden ser sustanciales (véase DORNBUSCH et al., 1998; CECCHETTI, 1999; MACLENNAN et al., 1999; MOJON, 2000; PEERSMAN y SMETS, 2001)
    } ¿Qué consecuencias tiene esto en los procesos de ajuste? Bajo los efectos de un mismo shock, por ejemplo un aumento del tipo de interés por parte del Banco Central Europeo, la transmisión es muy distinta entre Estados miembros de la unión monetaria;
    \item \textbf{Mercados financieros (empresariales)}. Las formas en que las empresas se financian es otra diferencia importante que existe entre los países de la Eurozona.\footnote{En países con una tradición legal anglosajona, las empresas tienden a acudir directamente a los mercados de capitales (mercados de bonos y acciones) para financiar sus proyectos de inversión. Como resultado, estos mercados están bien desarrollados, son sofisticados y altamente líquidos. En países con una tradición legal continental, las empresas obtienen financiación principalmente a través del sistema bancario, por lo que los mercados de capitales están menos desarrollados. 
    } ¿Qué consecuencias tiene esto en los procesos de ajuste? Por ejemplo, ante un aumento del tipo de interés, en países con sistemas financieros de tipo anglosajón donde los consumidores poseen una proporción significativa de bonos y acciones, se producen importantes efectos riqueza negativos ($\uparrow i_t\to \,\downarrow p_A$). En cambio, en países con sistemas financieros de tipo continental, estos efectos son menores y el impacto se transmite principalmente a través del canal del crédito bancario: un aumento suficientemente elevado de los tipos puede llevar a los bancos a racionar el crédito.
\end{lnum}

En conclusión, la forma en que un mismo aumento del tipo de interés afecta al consumo y a la inversión difiere considerablemente entre los países de la UE.

\paragraph{Crítica: }
Los mercados financieros siguen funcionando de manera diferente entre países (p.e. de la UE), lo que implica que los mismos shocks monetarios pueden transmitirse de forma distinta. 

Por ejemplo, antes de la unión monetaria europea, algunos países mantenían baja inflación (como Alemania), mientras que otros experimentaban inflación elevada (como Italia). En entornos de alta inflación, los inversores evitan bonos a largo plazo, ya que su precio es muy sensible a la inflación inesperada. Como resultado, en países con alta inflación predominan los mercados de deuda a corto plazo. Antes de la UEM, gran parte de la deuda italiana era a corto plazo, mientras que en Alemania predominaban los bonos a largo plazo.

Esto generaba asimetrías en la respuesta a cambios en los tipos de interés. Un aumento de tipos afectaba rápidamente al presupuesto italiano (por la corta madurez de la deuda), mientras que en Alemania el impacto era más gradual. Estas diferencias desaparecieron con la unión monetaria, ya que las estructuras de vencimiento convergieron. Sin embargo, esta convergencia se revirtió durante la crisis de deuda soberana, cuando los diferenciales de riesgo aumentaron y los países con mayor riesgo acortaron los vencimientos de su deuda. La lección es que la unión monetaria no elimina automáticamente estas diferencias estructurales.


\subsubsection{Financiación del déficit: Señoreaje}
Lo presentado hasta ahora es la esencia de la teoría tradicional de las Teoría de las Áreas Monetarias Óptimas desarrollada por MUNDELL (1961) y algunos de otros costes (o dificultades) que deben tenerse en cuenta. Sin embargo, esta teoría es incompleta. Pasa por alto otra implicación fundamental de la pérdida de independencia monetaria, probablemente la más importante. 

La entrada en una unión monetaria cambia de manera radical la capacidad de los gobiernos para financiar sus déficits presupuestarios: \textbf{los miembros de una unión monetaria emiten deuda en una moneda sobre la cual no tienen control}.\footnote{Por ejemplo, cuando Francia, Alemania y España entraron en la Eurozona, dejaron de emitir deuda en sus monedas nacionales, sobre las cuales tenían control total. En su lugar, ahora emiten deuda en una moneda que ninguno de estos gobiernos controla. Para cada uno de estos gobiernos, el euro es como una moneda extranjera. Por el contrario, en un país como el Reino Unido, el Gobierno puede garantizar plenamente a los tenedores de deuda pública que recibirán el pago en libras al vencimiento. La razón es que existe un banco central, el Banco de Inglaterra, que estará dispuesto (o será obligado) a proporcionar liquidez (libras) al gobierno si este se enfrenta a una escasez de liquidez que le impida pagar a los bonistas. Ninguno de los gobiernos de los países miembros de una unión monetaria tiene el poder de obligar al banco central común a proporcionar liquidez en momentos de crisis.
} Esto tiene una profunda consecuencia: los gobiernos no pueden ofrecer una garantía absoluta a los tenedores de bonos públicos de que dispondrán del efectivo necesario para pagarles cuando dichos bonos venzan. Este hecho implica que los mercados financieros adquieren el poder de forzar el impago de estos países. ¿Por qué? Veámoslo con un ejemplo.

\paragraph{Gestión de una crisis de confianza: Reino Unido vs. España}
Consideremos primero qué ocurriría si los inversores temieran que el gobierno del Reino Unido pudiera suspender pagos. En ese caso, venderían bonos del gobierno británico, elevando el tipo de interés. Tras vender estos bonos, los inversores tendrían libras, que probablemente intentarían vender en el mercado de divisas. El precio de la libra caería hasta que otros agentes estuvieran dispuestos a comprarla a ese precio más bajo. 

Pero el efecto de este mecanismo es que las libras permanecen dentro del mercado monetario británico para ser invertidas en activos del Reino Unido. Dicho de otro modo, la cantidad de dinero en el Reino Unido permanece inalterada y, probablemente, parte de ese dinero se reinvertirá en deuda pública británica. No obstante, incluso si esto no ocurriera y el gobierno no pudiera refinanciar su deuda a tipos de interés razonables, podría forzar al Banco de Inglaterra a proporcionarle el efectivo necesario para pagar a los bonistas. Así, el gobierno británico tiene garantizada la disponibilidad de liquidez para financiar su deuda, lo que significa que los inversores no pueden desencadenar una crisis de liquidez que obligue al gobierno del Reino Unido a suspender pagos. Es lo que se conoce como \textbf{autoridad de último recurso superior}: el Banco de Inglaterra.

La situación es radicalmente distinta para un miembro de una unión monetaria como España. Supongamos que los inversores temen un impago del gobierno español y, como resultado, venden bonos del gobierno elevando el tipo de interés. Los inversores que ahora poseen euros probablemente decidirán invertirlos en otros activos, por ejemplo en bonos del gobierno alemán. Como resultado, \textbf{los euros abandonan el sistema bancario español y no existe un mercado de divisas ni un tipo de cambio flexible que detenga este proceso}. Por tanto, la cantidad total de liquidez (oferta monetaria) en España se reduce lo que podría, en última instancia, provocar una crisis de liquidez para el Gobierno: no puede obtener financiación para renovar su deuda a tipos de interés razonables. En contraste con la situación anterior, el gobierno español no puede obligar al Banco de España a proporcionarle liquidez, porque este ya no tiene esa autoridad. El banco central común (el BCE en la Eurozona) puede proporcionar liquidez, pero el gobierno español no controla esa institución. Si la crisis de liquidez es suficientemente intensa, puede obligar al gobierno español a suspender pagos, ya que no dispone de efectivo para pagar a los bonistas.

El problema raíz surge en que \textbf{los mercados financieros son conscientes de ello y pondrán a prueba al gobierno español cuando los déficits presupuestarios empeoren}. Así, en una unión monetaria, los mercados financieros adquieren un enorme poder y pueden forzar a cualquier país miembro a una situación crítica. Esto resulta paradójico: una unión monetaria no fortalece la posición de los gobiernos nacionales frente a los mercados financieros; por el contrario, los debilita.\footnote{La situación de España recuerda a la de las economías emergentes que deben endeudarse en moneda extranjera. Estas economías enfrentan el mismo problema: pueden sufrir un \textit{sudden stop}, es decir, una interrupción repentina de los flujos de capital que provoca una crisis de liquidez (véase Calvo, 1988; Eichengreen et al., 2005).
}

El análisis anterior pone de relieve la fragilidad de una unión monetaria. Cuando los inversores desconfían de un gobierno, venden sus bonos, elevando los tipos de interés y desencadenando una crisis de liquidez. Esta puede derivar en un problema de solvencia: el aumento de los tipos incrementa la carga de la deuda, obligando al gobierno a reducir el gasto y aumentar los impuestos. Estas políticas de austeridad forzada tienen costes políticos y pueden llevar al gobierno a dejar de pagar la deuda, produciéndose así una dinámica de profecía autocumplida: cuando los mercados dudan de la capacidad o voluntad de un gobierno para pagar su deuda, venden los bonos, lo que aumenta la probabilidad de impago.

\paragraph{Shocks asimétricos: Amplificación de los ciclos}
Este último problema adquiere una importancia superior en un escenario de ajuste ante un shock asimétrico. ¿Por qué? Pensemos cómo se ve afectado este ajuste cuando se tienen en cuenta las implicaciones presupuestarias. ¿Cuáles son los efectos sobre el presupuesto del Gobierno?
\begin{lnum}
    \item \textbf{Ingresos fiscales}. En primer lugar, la caída del producto interior bruto (PIB) conduce a una disminución de los ingresos fiscales del gobierno. Esta caída es probablemente más que proporcional a la disminución del PIB, porque los impuestos sobre la renta son progresivos;
    \item \textbf{Estabilizadores automáticos}. En segundo lugar, dado que el desempleo aumenta, el gasto público también se incrementa. 
\end{lnum}

Al combinar estos dos efectos: (i) el déficit presupuestario del gobierno aumenta; y, (ii) se produce automáticamente como consecuencia de la caída del PIB. Ahora bien, si la caída de la demanda agregada es lo suficientemente intensa, el aumento automático del déficit presupuestario puede llegar a ser tan grande que los inversores comiencen a dudar de la solvencia del gobierno. Y, ¿qué sucederá si carece de autonomía de ajuste monetario? La pérdida de confianza llevará a los inversores a vender bonos públicos, lo que provocará un aumento del tipo de interés y una crisis de liquidez. 

\imagenfit{Fig1_5.png}{Amplificación de los shocks: asimetría}{}

Como vemos, la curva de demanda agregada se desplazará hacia la izquierda de manera amplificada: (i) con tipos de interés más altos, los residentes franceses reducen su consumo y su inversión ($D_F\to D’_F$); y, (ii) la crisis de deuda añade un efecto amplificado ($D_F'\to D’’_F$). En síntesis, la crisis de deuda amplifica el shock negativo inicial. Y, ¿cuál es el efecto de la crisis de deuda sobre el lugar \textit{seguro} para los inversores (Alemania, p.e.)? Cuando los inversores venden bonos obtienen liquidez (euros) que desean reinvertir. Dado que previamente mantenían bonos públicos, es probable que busquen otros bonos gubernamentales más seguros, lo que hace más probable que compren bonos del gobierno alemán. El efecto de estas compras es que el precio de los bonos alemanes aumenta, lo que reduce su rendimiento (tipo de interés). Esto se traduce en un estimulo adicional para la economía, pues este flujo de liquidez provoca una disminución del tipo de interés en Alemania, lo que a su vez amplifica el incremento de la demanda agregada alemana ($D_G'\to D_G''$).

De este análisis se concluye que la crisis de deuda conduce a una amplificación del shock asimétrico: intensifica los efectos negativos en un país y los efectos positivos en otro(s). En definitiva, estos cambios en los tipos de interés, en lugar de estabilizar el sistema tienden a desestabilizarlo, agravando los problemas de ajuste en ambos países.\footnote{Esto puede resultar sorprendente ya que en una unión monetaria lo tipos de interés deberían ser iguales (o convergerntes). ¿No es característico de una unión monetaria que los tipos de interés sean iguales en todos los países? Esto es cierto para el tipo de interés a corto plazo, bajo el control del banco central común. Sin embargo, los tipos de interés a largo plazo pueden divergir, y son precisamente estos los tipos aplicados a los bonos públicos a largo plazo. Por tanto, lo normal es observar divergencia cuando los inversores perciben distintos niveles de riesgo entre los diferentes emisores soberanos.
}


\subsubsection{Implicaciones de Política Económica: ¿Cuán efectivo es mantener la independencia monetaria?}




\subsection{Beneficios de una Unión Monetaria}

Una moneda común tiene beneficios importantes. En este capítulo se identificaron varios de ellos. Primero, una moneda común en Europa reduce los costes de transacción. Esto produce beneficios no solo directos, sino también indirectos, al estimular la integración económica en Europa.

Segundo, al reducir la incertidumbre de precios, una moneda común mejora la eficiencia asignativa del mecanismo de precios. Esto mejora el bienestar, aunque es difícil cuantificar este efecto.

Tercero, la principal mejora de bienestar derivada de eliminar los tipos de cambio dentro de la unión consiste en que elimina los movimientos extremos de dichos tipos de cambio. Estos cambios extremos en los tipos de cambio, que han ocurrido regularmente dentro de la Unión Europea en el pasado, generan importantes perturbaciones macroeconómicas. Se convierten ellos mismos en fuentes de shocks asimétricos, en lugar de ser instrumentos para ajustarse a shocks asimétricos.

No obstante, también debemos tener presente que la existencia de una moneda común puede crear nuevos riesgos. Estos surgen cuando los países encuentran difícil ajustarse a perturbaciones. Una unión monetaria también crea riesgos específicos para los gobiernos nacionales. Al entrar en una unión monetaria, pierden la capacidad de garantizar que el stock pendiente de bonos públicos será siempre reembolsado. Esto puede perturbar los mercados financieros, creando nuevos riesgos que no existían fuera de la unión monetaria. Como resultado, la disminución del riesgo cambiario no necesariamente reduce el riesgo sistémico en la UEM.

Cuarto, la mayor transparencia de precios proporcionada por el uso de una moneda común probablemente aumentará algo la competencia, beneficiando a los consumidores. Sin embargo, también concluimos que la fuente principal de beneficios no es tanto la mayor transparencia de precios. Es el hecho de que la existencia del euro puede desencadenar integración en otros ámbitos —financieros, institucionales, políticos— lo que conducirá a mayor competencia y ganancias de bienestar.

Quinto, si la nueva moneda común llega a convertirse en una verdadera moneda global, pueden obtenerse beneficios adicionales en forma de ingresos públicos y de expansión de la industria financiera dentro de la unión.




Mientras que los costes de una moneda común tienen mucho que ver con la gestión macroeconómica de la economía, los beneficios se sitúan principalmente en el nivel microeconómico. Cabe esperar que eliminar las monedas nacionales y pasar a una moneda común genere ganancias de eficiencia económica con dos orígenes distintos: (i) uno es la eliminación de los costes de transacción asociados al cambio de monedas nacionales; (ii) el otro es la eliminación del riesgo derivado de los movimientos futuros inciertos de los tipos de cambio. 

\subsubsection{Costes de transacción: Ganancias directas}
Eliminar los costes de cambiar una moneda por otra es, sin duda, la ganancia más visible de una unión monetaria. Todos experimentamos estos costes cada vez que cambiamos moneda y desaparecen cuando los países pasan a una moneda común.

¿Qué magnitud tienen las ganancias derivadas de la eliminación de los costes de transacción? La Comisión Europea ha estimado estas ganancias y obtiene una cifra entre 13.000 y 20.000 millones de euros al año, que representa entre un cuarto y medio punto porcentual del PIB comunitario. Puede parecer poca cosa. Sin embargo, es una ganancia que debe añadirse a las demás ganancias derivadas de una unión monetaria.

Debe señalarse aquí que estas ganancias que recibe el público general tienen una contrapartida en algún lugar. Se encuentran principalmente en el sector bancario. Encuestas en distintos países indican que alrededor del 5 por ciento de los ingresos bancarios son comisiones pagadas a los bancos por cambiar monedas nacionales. Esta fuente de ingresos para los bancos desaparece con una unión monetaria.

Lo anterior no debe dar la impresión de que la ganancia para el público se compensa con la pérdida para los bancos. Los costes de transacción implicados en el cambio de dinero son una pérdida irrecuperable. Son como un impuesto pagado por el consumidor a cambio de nada. Los bancos, sin embargo, tienen un problema de transición: deben buscar otras actividades rentables. Cuando esto se ha hecho, la sociedad ha ganado. Los empleados bancarios, anteriormente dedicados al cambio de dinero, quedan libres para desempeñar otras tareas para la sociedad.

Debe mencionarse aquí un punto importante. Los beneficios de una moneda común surgen no solo por el uso de los mismos billetes, sino también porque el sistema de pagos está integrado. La integración del sistema de pagos en la Eurozona no se produjo inmediatamente al inicio de la Eurozona en 1999. Llevó algún tiempo que el sistema de pagos dentro de la Eurozona estuviera plenamente integrado. Esto se ha logrado hoy.

\begin{specialblock}{Sistema de Pagos -- TARGET2}
    Para que la Eurozona funcione de manera eficiente, el sistema de pagos debe estar integrado, de modo que los pagos transfronterizos puedan gestionarse con la misma fluidez que los pagos dentro de un mismo país. Un elemento clave de este sistema de pagos son las transferencias que los bancos realizan entre sí. Estas transferencias están centralizadas en el banco central. Esta es también la razón por la cual el banco central es denominado el “banco de los bancos”. Antes del inicio de la Eurozona, cada país tenía su propio sistema de pagos centralizado por su banco nacional. Por tanto, era necesario integrar estos sistemas de pagos en la Eurozona. Esta integración se logró mediante el sistema TARGET en 1999. Posteriormente fue mejorado y renombrado como TARGET2. Sus principales características son las siguientes: * Es un sistema en tiempo real. Esto significa que los pagos que los bancos realizan entre sí llegan a su destino “instantáneamente”, es decir, con un retraso de unos pocos segundos o minutos. * Es un sistema de liquidación bruta. Esto implica que el importe bruto de cada pago pasa a través de TARGET2. Esto exige que los bancos que realizan pagos proporcionen colateral para cada transacción. Esto contrasta con un sistema de liquidación neta, utilizado en la mayoría de los sistemas nacionales de pagos. En estos, los bancos pueden acumular posiciones netas deudoras o acreedoras durante el día sin necesidad de aportar colateral, liquidándose dichas posiciones al final del día. La razón por la cual se seleccionó TARGET2 es que elimina el riesgo de que el incumplimiento de un banco genere un efecto dominó sobre otros bancos implicados en la cadena de pagos. Desde la crisis de deuda soberana en 2010, han surgido importantes desequilibrios en el sistema TARGET2. En particular, los países del sur de la Eurozona, que habían acumulado grandes déficits por cuenta corriente, se enfrentaron repentinamente a fuertes salidas de liquidez. Lo contrario ocurrió en los países del norte de la Eurozona, que se convirtieron en receptores de estos flujos de liquidez. Como resultado, los países del sur acumularon importantes pasivos en el sistema TARGET2, compensados por activos en los países del norte, principalmente Alemania. (Véase Sinn y Wollmershäuser, 2012, quienes argumentan que esto genera riesgos significativos para el contribuyente alemán; para críticas, véase De Grauwe y Ji, 2012b, y Whelan, 2012). Entre 2012 y 2014, estos desequilibrios en TARGET2 disminuyeron significativamente debido a flujos de liquidez en sentido inverso —es decir, entradas de liquidez en los países del sur y salidas desde los países del norte—. Sin embargo, desde comienzos de 2015, los desequilibrios en el sistema TARGET2 volvieron a aumentar. Las razones de este fenómeno se analizan en el Capítulo 12.
\end{specialblock}

\subsubsection{Costes de transacción: Ganancias indirectas}
La eliminación de los costes de transacción también tiene una ganancia indirecta: debería conducir a una mayor transparencia de precios; es decir, los consumidores, que ahora pueden ver los precios en la misma unidad monetaria, están en mejores condiciones para comparar precios y buscar mejores ofertas. Esto, genera multitud de efectos indirectos: (i) debería aumentar la competencia; y, en consecuencia, (ii) beneficiar a todos los consumidores, que se enfrentarán a los mismos precios más bajos. 

La pregunta entonces sería, ¿sucede verdaderamente? Aunque existe mucha evidencia de que la discriminación de precios todavía se practica ampliamente en Europa,\footnote{Para una amplia gama de productos de marca en el año 2011, podemos ver como:

\imagenfit{Fig3_1.png}{Diferenciales de precios para bienes de consumo (\%) relativos a la media de la Eurozona}{}

En Finlandia esta cesta de bienes y servicios era un 22 por ciento más cara que en la Eurozona, mientras que en Eslovaquia era un 30 por ciento más barata que la media de la Eurozona. Así, existía un diferencial de precios del 50\% entre el país más barato y el más caro de la Eurozona. Más aún, HASKEL y WOLF (2001) muestran que existen diferenciales de precios similares para exactamente los mismos productos vendidos por la tienda de muebles IKEA en Europa. También los diferenciales de precios en el mercado automovilístico son notorios, y BAYE et al. (2006) han mostrado que los precios listados en Internet para una muestra de artículos no han convergido tras la introducción del euro.
} lo que sorprende es que estos diferenciales de precios entre países suelen ser mucho mayores que dentro de los países. ENGEL y ROGERS (1995) observaron que, para diferenciales de precios de los mismos pares de bienes en distintas ciudades norteamericanas (en Estados Unidos y Canadá), cruzar una frontera equivale a recorrer $2.500$ millas dentro del mismo país. En otras palabras, los diferenciales de precios entre Detroit y Windsor (justo al otro lado de la frontera) son del mismo orden de magnitud que los diferenciales de precios entre Nueva York y Los Ángeles, parece que las fronteras son bastante poderosas a la hora de segmentar mercados e introducir grandes diferenciales de precios. Así, en Europa y Norteamérica, la existencia de fronteras sigue generando fuertes impedimentos al comercio, pese a que los aranceles de importación y otras barreras comerciales explícitas han sido eliminados.

También es útil estudiar la evolución de la convergencia de precios a lo largo del tiempo para ver si aparecen tendencias:

\imagenfit{Fig3_2.png}{Evolución de la dispersión de precios en la Eurozona, medidos por la desviación estándar media de 173 productos idénticos: 1990-2015}{WOLSZCZAK-DERLACZ (2006)}

La observación notable es que la convergencia de precios ocurrió antes de 1999 (principalmente a comienzos de la década de 1990) y, desde el inicio de la Eurozona, la convergencia de precios se ha detenido.\footnote{Este fenómeno también ha sido observado por la Comisión Europea (véase Comisión Europea, 2004). Trabajos empíricos más recientes confirman que esta falta de convergencia de precios en la Eurozona es una característica persistente; es decir, no hay evidencia de que desde 2005 la convergencia de precios haya aumentado (véanse PARSLEY y WEI, 2008; CLEMENTI et al., 2010). Un estudio de los mercados de automóviles en la Eurozona también confirma que el euro ha tenido poco efecto en reducir la considerable dispersión de precios para los mismos automóviles dentro de la Eurozona (Dvir y Strasser, 2013).
} ¿Por qué la introducción del euro parece ser una fuerza débil para generar convergencia de precios?\footnote{Muchos de los productos en la muestra de WOLSZCZAK-DERLACZ son productos de supermercado. En este caso el arbitraje sería claramente prohibitivo en la Eurozona. Pero ¿por qué entonces observamos casi ningún diferencial de precios para productos de supermercado dentro de un mismo país? La razón es que el negocio minorista sigue estando muy segmentado nacionalmente. En la mayoría de los países, unas pocas cadenas de supermercados dominan todo el mercado, llevan a cabo campañas comerciales y publicitarias nacionales y fijan precios para todo el mercado nacional. Otras razones son que la mayoría de estos minoristas siguen siendo compañías muy nacionales y enfrentan diferentes regulaciones, costumbres, lenguas y culturas. 

No obstante, otros productos de la muestra son productos electrónicos (por ejemplo, cámaras, teléfonos móviles, etc.) y, cabe suponer, que los diferenciales de precios de estos productos no se deben tanto a costes de transacción. Sin embargo, este es un sector de productos altamente diferenciados, lo que hace muy difícil para los consumidores realizar comparaciones precisas de precios. Razón por la cual los diferenciales de precios de estos productos siguen siendo elevados incluso dentro de países con la misma moneda. Así, aunque el euro pueda facilitar algo la comparación de precios, es dudoso que contribuya mucho a eliminar los diferenciales observados.
} Existen múltiples razones pero existe consenso en que si el euro contribuye a la convergencia de precios, no será tanto porque facilite las comparaciones directas de precios por parte de los consumidores, sino porque pueda contribuir a una mayor integración económica de otras formas. Por ejemplo, es posible que la introducción del euro estimule la integración financiera, lo que puede poner en marcha una dinámica de integración en otras áreas (p.e. armonización legislativa). Así, la existencia del euro puede convertirse en un importante detonante de una mayor integración en muchos otros ámbitos (político, legislativo, regulatorio). 

\subsubsection{Incertidumbre: Efectos sobre el bienstar}
Otra fuente de posibles ganancias, ligada a la anterios pero ligeramente diferente, tiene que ver con los efectos que la incertidumbre genera sobre los ingresos futuros de las empresas. Se acepta, generalmente, que esto conduce a una pérdida de bienestar en un mundo poblado por individuos aversos al riesgo, ya que estos preferirán un rendimiento futuro más cierto a uno menos cierto, al menos si el valor esperado de esos rendimientos es el mismo. Dicho de otro modo, solo querrán asumir el rendimiento más arriesgado si se les promete que, en promedio, será superior al menos arriesgado. Por tanto, cabe suponer que eliminar el riesgo cambiario reduce una fuente de incertidumbre y debería, por tanto, aumentar el bienestar.

\paragraph{Crítica: DE GRAUWE}
Existe una característica importante de la teoría de la empresa que puede invalidar esa conclusión. Considérese una empresa maximizadora de beneficios que es precio-aceptante en el mercado de producto. Supongamos también que la empresa exporta toda su producción. Representamos su curva de coste marginal y el precio de su producto:

\imagenfit{Fig3_3.png}{Beneficios de la firma bajo certidumbre e incertidumbre de precios}{}

El precio que obtiene la empresa viene dado por el precio en el mercado de exportación multiplicado por el tipo de cambio. Supongamos que existen dos regímenes: (i) en el panel izquierdo, el tipo de cambio es fijo y como consecuencia, el precio obtenido por las empresas es constante y perfectamente predecible (suponiendo que el precio en moneda extranjera es constante); y, (ii) en el panel derecho, el tipo de cambio fluctúa aleatoriamente, y, como consecuencia, genera fluctuaciones aleatorias de precios (supongamos que lo hace simétricamente entre $p_2$ y $p_3$).

En el primer régimen de certidumbre, el beneficio de la empresa en cada período viene dado por el área sombreada menos el área $FGp_1$. En el segundo régimen, incierto, el beneficio fluctuará dependiendo de si prevalece el precio $p_2$ o $p_3$. Entonces, podemos ver que, en promedio, el beneficio será mayor en el régimen incierto: cuando el precio es bajo, el beneficio es menor que en el caso de certeza, cuantía igual al área $p_1BCp_2$; pero cuando el precio es alto, el beneficio es mayor que en el caso de certeza, por una cuantía igual al área $p_3EBp_1$. Y, como $p_3EBp_1>p_1BCp_2$, esto implica que en promedio sea mayor.\footnote{Como puede parecer contraintuitivo, podemos interpretarlo del siguiente modo. Cuando el precio es alto, la empresa aumenta la producción para aprovechar el mayor ingreso por unidad producida. Así, obtiene un mayor beneficio por cada unidad de producción que habría producido de todos modos y, además, expande su producción. Este último efecto se mide mediante el triángulo sombreado superior. Cuando el precio es bajo, sin embargo, la empresa hará lo contrario: reducirá la producción. Al hacerlo, limita la reducción de su beneficio total, efecto que está representado por el triángulo sombreado inferior.
}

Si se quieren realizar comparaciones de bienestar entre un régimen de certeza de precios y uno de incertidumbre de precios, el efecto positivo de la incertidumbre sobre el beneficio medio debe compararse con la mayor incertidumbre sobre esos beneficios. El mayor beneficio medio aumenta la utilidad de la empresa, mientras que la mayor incertidumbre sobre esos beneficios reduce la utilidad de la empresa aversa al riesgo. Por tanto, no está claro si el bienestar disminuye cuando aumenta la incertidumbre cambiaria o, inversamente, si podemos afirmar con gran confianza que el bienestar de las empresas aumenta cuando se eliminan las monedas nacionales y se introduce una moneda común.\footnote{Otra forma de formular este análisis es reconocer que las variaciones del tipo de cambio crean oportunidades de obtener beneficios. Cuando el tipo de cambio se vuelve más variable, aumenta la probabilidad de obtener grandes beneficios. En cierto sentido, exportar puede verse como una opción. Cuando el tipo de cambio se vuelve muy favorable, la empresa ejerce su opción de exportar. Con un tipo de cambio desfavorable, la empresa no ejerce esta opción. Es bien sabido por la teoría de opciones que el valor de la opción aumenta cuando aumenta la variabilidad del activo subyacente. Así, la empresa que tiene la opción de exportar mejora su situación cuando el tipo de cambio se vuelve más variable. Pueden añadirse muchas complicaciones a esta teoría. Puede introducirse el supuesto de competencia imperfecta, o el supuesto de que las empresas enfrentan costes de ajuste cuando varían la producción. Sin embargo, aunque estas complicaciones proporcionan importantes ideas sobre el efecto de la incertidumbre de precios sobre los beneficios medios, en estos modelos suele ocurrir que la incertidumbre de precios puede aumentar los beneficios medios de la empresa.
}

\subsubsection{Incertidumbre: efectos sobre el crecimiento}
Otro aspecto importante relativo a la incertiodumbre es que la eliminación del riesgo cambiario conducirá a un aumento del crecimiento económico. Sus implicaciones pueden observarse utilizando el modelo neoclásico de crecimiento y su extensión a situaciones de economías dinámicas de escala.\footnote{Este análisis tuvo un papel destacado en el informe de la Comisión Europea One Market, One Money (1990), que a su vez estuvo muy influido por BALDWIN (1989). Este análisis fue muy influyente para vender la idea de una unión monetaria en Europa como herramienta para impulsar el crecimiento económico.

El eje horizontal muestra el stock de capital por trabajador, y el eje vertical la producción por trabajador. La línea $f(K)$ es la función de producción, que tiene la forma cóncava habitual, lo que implica productividades marginales decrecientes. 
} 

\imagenfit{Fig3_5.png}{El modelo neoclásico de crecimiento económico}{}

Como sabemos, el equilibrio se obtiene donde la productividad marginal del capital es igual al tipo de interés que los consumidores utilizan para descontar el consumo futuro, representado mediante el punto $A$, donde la línea $rr$ (cuya pendiente es igual a la tasa de descuento) es tangente a la función de producción $f(k)$. 

Supóngase que la eliminación del riesgo cambiario reduce el riesgo sistémico. Esto tendrá el efecto de reducir el tipo de interés real: en un entornos menos arriesgado, los inversores requerirán una menor prima de riesgo para realizar la misma inversión. Además, cuando los agentes descuenten el futuro, estarán dispuestos a utilizar una tasa de descuento menor.

\imagenfit{Fig3_6.png}{Los efectos de menor riesgo en el modelo neoclásico de crecimiento}{}

La reducción de la tasa de descuento ajustada por riesgo hace que la línea $rr$ sea más plana, girándola hasta $r’r’$. Como resultado, el equilibrio se desplaza de $A$ a $B$. Habrá acumulación de capital y un aumento de la tasa de crecimiento mientras la economía se mueve de $A$ a $B$. En el nuevo equilibrio, la producción por trabajador y el stock de capital a disposición de cada trabajador habrán aumentado. Así, en este modelo neoclásico de crecimiento, la reducción del tipo de interés debida a la unión monetaria aumenta temporalmente la tasa de crecimiento de la producción. En el nuevo equilibrio, habrá aumentado el nivel de producción por trabajador.\footnote{Obsérvese, sin embargo, que la tasa de crecimiento de la producción vuelve entonces a su nivel inicial, determinado por la tasa exógena de cambio tecnológico y la tasa de crecimiento de la población. Además, la productividad del capital ha disminuido.
} 

\paragraph{Crítica: Eviencia Empírica}
¿Cuánto se ha materializado esta promesa? Comparando las tendencias del PIB real durante $2000-18$ de la zona euro con la tendencia del PIB real en Estados Unidos y en los países de la UE que no son miembros de la zona euro (denominados UE-10):

\imagenfit{Fig3_8.png}{}{}

Es llamativo que el PIB haya aumentado a una tasa significativamente menor en la zona euro que en Estados Unidos y en la UE-10. Puede haber muchas razones por las que la tasa de crecimiento fue tan baja en la zona euro, pero si una unión monetaria fuera una máquina impulsora del crecimiento, probablemente no habríamos observado esto. La impresión que obtenemos es que existe muy poca evidencia de que el euro haya impulsado el crecimiento como prometía el análisis teórico anterior. ¿Por qué parece que la teoría que predecía efectos impulsores del crecimiento de la unión monetaria no se ha visto confirmada? La razón principal probablemente es que la reducción de la incertidumbre cambiaria dentro de la unión no parece haber conducido a una disminución significativa del tipo de interés real. Hay muy poca evidencia de que el tipo de interés real en la Eurozona en su conjunto haya disminuido.\footnote{Obsérvese que desde la crisis financiera iniciada en 2007 los tipos de interés reales han mostrado una tendencia a disminuir. Este fenómeno, sin embargo, ha ocurrido en la mayoría de los países industrializados, independientemente de si pertenecen o no a la Eurozona (véanse SUMMERS, 2014, y TEULINGS y BALDWIN, 2014).

Solo en los países en proceso de convergencia, como Irlanda, España, Portugal y Grecia, el tipo de interés real disminuyó significativamente. Es en estos países (con la excepción de Portugal) donde observamos una aceleración del crecimiento económico antes de la crisis financiera de 2007-08, como predecía la teoría. Gran parte de este crecimiento, sin embargo, resultó ser temporal. Además, el crecimiento significativamente menor observado en la Eurozona en comparación con Estados Unidos y la UE-10 desde 2008 está muy probablemente relacionado con la aplicación de intensos programas de austeridad en un gran número de países de la Eurozona.
} El débil vínculo entre incertidumbre cambiaria, tipo de interés real y crecimiento también puede deberse al hecho de que la reducción de la incertidumbre cambiaria no reduce necesariamente el riesgo sistémico. Una menor incertidumbre cambiaria puede verse compensada por una mayor incertidumbre en otros ámbitos, por ejemplo, incertidumbre sobre la producción y el empleo, e incertidumbre sobre la sostenibilidad de la deuda pública. Como resultado, las empresas que operan en una zona monetaria mayor pueden no operar, en promedio, en un entorno menos arriesgado. Existe toda una literatura teórica, iniciada por William Poole (1970), que ha analizado este problema. Presentamos los principales resultados en el Recuadro 3.2.

\subsubsection{Unidad monetaria: Una moneda internacional}
Cuando los países forman una unión monetaria, es probable que la nueva moneda resultante tenga más peso en las relaciones monetarias internacionales que la suma de las monedas individuales antes de la unión. Como resultado, es probable que la nueva moneda común encuentre un uso creciente fuera de la unión, lo que crea beneficios adicionales de la unión monetaria. 

¿Por qué? Las ventajas de tener una moneda que se utiliza como unidad de cuenta y medio de intercambio en el resto del mundo son significativas. Distinguimos tres fuentes de beneficios:
\begin{lnum}
    \item \textbf{Efectos sobre el Balance del Banco Central}. Primero, cuando una moneda se usa internacionalmente, el emisor de esa moneda obtiene ingresos adicionales. Por ejemplo, en 1999 más de la mitad de los dólares emitidos por la Reserva Federal se utilizaban fuera de Estados Unidos. Esta situación tiene el efecto de más que duplicar el tamaño del balance de la Reserva Federal respecto a una situación en la que el dólar se utilizara solo domésticamente. De ello se sigue que los beneficios potenciales de la Reserva Federal también se más que duplican. Dado que estos beneficios van al gobierno de Estados Unidos, los ciudadanos estadounidenses disfrutan de los beneficios del uso mundial del dólar en forma de menores impuestos necesarios para financiar un determinado nivel de gasto público. Si el euro se convierte en una moneda mundial como el dólar, los ciudadanos de la Eurozona disfrutarán de beneficios similares. De hecho, existe cierta evidencia de que el euro se utiliza en Europa Central y Oriental, creando ingresos para los ciudadanos de la Eurozona. No obstante, no debe exagerarse la magnitud de estos beneficios. Los beneficios totales de la Reserva Federal de Estados Unidos ascienden a menos del 0,5 por ciento del PIB estadounidense. Por tanto, los ingresos adicionales derivados de tener una moneda internacional siguen siendo relativamente pequeños.;
    \item \textbf{Moneda de reserva internacional}. Una segunda fuente de beneficios se relaciona con el hecho de que una moneda internacional también es mantenida como reserva internacional por bancos centrales extranjeros. Típicamente, estas reservas no se mantienen en forma de efectivo sino como títulos del Tesoro. Así, el Banco Central de China mantiene más de un billón de dólares en forma de títulos del Tesoro estadounidense. Otros bancos centrales asiáticos también mantienen muchos cientos de miles de millones de dólares en títulos del Tesoro estadounidense. Estas tenencias extranjeras han sido una fuente importante de financiación fácil para los déficits presupuestarios de Estados Unidos durante la última década. La peculiaridad de esta financiación es que los tenedores extranjeros soportan el riesgo cambiario. Es en este sentido que el presidente francés De Gaulle habló del privilegio exorbitante de Estados Unidos (véase EICHENGREEN, 2012);
    \item \textbf{Efectos sobre el Mercado financiero doméstico}. Un tercer beneficio, probablemente el mayor, pero más difícil de cuantificar: cuando una moneda se convierte en internacional, esto impulsará la actividad de los mercados financieros domésticos. Los residentes extranjeros querrán invertir en activos y emitir deuda en esa moneda. Como resultado, los bancos domésticos atraerán negocio, y también lo harán los mercados de bonos y acciones. Esto, a su vez, crea conocimiento especializado y empleo. Así, si el euro se convierte en una moneda internacional como el dólar, es probable que cree nuevas oportunidades para las instituciones financieras de la Eurozona. Aquí también es necesaria una advertencia. Algunos países, como el Reino Unido, han sido capaces de atraer actividades financieras del resto del mundo sin el respaldo de una moneda local que sea verdaderamente internacional. La City de Londres es hoy un importante centro de finanzas internacionales, pese a que la libra esterlina ya no desempeña un papel importante en el mundo. Así, tener una moneda internacional no es una condición necesaria para generar servicios financieros por los que el resto del mundo esté dispuesto a pagar. Tampoco es, por lo demás, una condición suficiente.
\end{lnum}

En este sentido, el euro se utiliza ahora como moneda de reserva por bancos centrales extranjeros. Mostramos la evolución del uso del euro como moneda de reserva y la comparamos con el uso del dólar en la Figura 3.11. Observamos que la participación del euro como moneda de reserva ha fluctuado entre el 20 y el 25 por ciento. Esto contrasta con el uso del dólar como moneda de reserva, que ha fluctuado entre el 60 y el 70 por ciento.

\imagenfit{Fig3_11.png}{Participación en reservas extranjeras de algunas monedad internacionales}{}




\subsubsection{Unidad Monetaria: Apertura de los países}
Antes se analizaron los costes de la unión monetaria y la apertura del país, también podemos derivar una relación entre los beneficios de una unión monetaria y la apertura de un país. Es probable que las ganancias de bienestar de una unión monetaria identificadas en este capítulo aumenten con el grado de apertura de una economía. Por ejemplo, la eliminación de los costes de transacción tendrá más peso en países donde las empresas y los consumidores compran y venden una gran proporción de bienes y servicios en países extranjeros. De manera similar, los consumidores y las empresas de estos países están más expuestos a errores de decisión porque se enfrentan a grandes mercados extranjeros con monedas distintas. Eliminar estos riesgos generará una mayor ganancia de bienestar (per cápita) en economías pequeñas y abiertas que en países grandes y relativamente cerrados.

\imagenfit{Fig3_12.png}{}{}

Si mostramos en el eje horizontal mostramos la apertura del país respecto de sus socios potenciales en la unión monetaria, medida por la participación de su comercio bilateral en el PIB del país considerado, y en el eje vertical representamos los beneficios, como porcentaje del PIB, podemos ver como, a medida que aumenta la apertura hacia los demás socios de la unión, aumentan las ganancias de una unión monetaria por unidad de producción.
















































\clearpage
\section{Evidencia empírica: Evaluaciones de impacto}
\puntosclave{
    En la actualidad existen diversas uniones monetarias en la economía mundial que se han materializado. Se recomienda una lectura ponnenorizada del Anexo 2.
    
    La unión monetaria europea ha aumentado los flujos comerciales de los países miembros.
    
    Al mismo tiempo la Unión Económica y Monetaria (Eurozona) aumenta la probabilidad de sincronización del ciclo económico, en lo que se refiere a la transmisión de shocks asimétricos. 
    
    En Mercosur, se demuestra que una cesta compuesta por tipos de cambio de diferentes monedas de países miembros tiene un comportamiento similar a una moneda agregada de carácter estable, constituyendo bases para una futura integración monetaria. 
    
    Igualmente, para el caso de Mercosur surgen opciones adicionales de integración monetaria, tales como implementar una caja de convertibilidad dual, ya que concentra sus flujos comerciales con las áreas monetarias del euro y el dólar.
}

Desde una vertiente académica, el impacto de la unión monetaria en la UE sobre el comercio de los Estados miembros habría sido positivo, a partir de la utilización de modelos de gravedad de comercio bilateral. El estudio que se puede considerar como seminal (Rose y Engel, 2002), estima que pertenecer a una unión monetaria aumenta los flujos de comercio de un país 6,5 veces más en comparación con aquellos países que no pertenecen a la unión monetaria. Esta evidencia hasido corroborada posterionnente por diversos autores (Baldwin y Taglioni, 2007; Olivero y Yotov, 20 12; Gunnella et al., 202 1, entre otros).\footnote{El modelo de gravedad de comercio parte de la Segunda Ley de Newton, donde la fuerza de atracción entre dos planetas es directamente proporcional a la masa de los mismos e inversamente proporcional a la distancia que separa ambos planetas. Tinbergen ( 1 962) extrapola este paradigma a la economía. considerando que el comercio bilateral entre dos países mantiene una relación positiva con el tamaño de los países y negativa con la distancia que los separa. Así, se espera que los países concentren sus relaciones comerciales con aquellos países más ricos en términos de renta y más cercanos geográficamente.
} El contexto teórico asume que una unión monetaria reduce los costes de transacción, al compartir los países una moneda común (Alesina et al., 2002), lo que aumenta los flujos comerciales (Larch et al., 20 1 9). Resulta convenient recordar que los costes de transacción son uno de los componentes más importantes de los costesde comercio (Anderson y Van Wincoop, 2004), con lo que una reducción sustancial de los mismos se espera que se traduzca en una importante mejora de los fhtjos comerciales. Sin embargo, al hablar de la unión monetaria europea resulta necesario matizar esta afirmación previa y rebajar el optimismo por diferentes motivos. En primer lugar, debido a la crítica que suscita la utilización de modelos de gravedad de comercio para estimar los impactos de decisiones de política económica. Por ejemplo, de cara a evaluar el impacto de la salida del Reino Unido de la Unión Europea, el Brexit, sobre los flujos comerciales comunitarios, surge el estudio de Gudgin et al. (20 1 7) realizan una revisión bibliográfica sobre las estimaciones realizadas a partir de modelos de gravedad y concluyen que en todos los estudios se reportó un impacto menor del que 





\subsubsection{Unión monetaria y comercio: la evidencia empírica}

Concentrémonos ahora en los efectos comerciales de una unión monetaria. Identificamos dos mecanismos mediante los cuales una unión monetaria podría aumentar el comercio entre sus miembros. El primero tiene que ver con los costes de transacción: la unión monetaria reduce los costes de transacción y, por tanto, estimula el comercio internacional. El segundo mecanismo se relaciona con la incertidumbre cambiaria. Esta también podría ser un factor que estimule el comercio entre los miembros de una unión monetaria.

¿Cuál es la evidencia empírica sobre los efectos de una unión monetaria en el comercio? La primera generación de estudios econométricos generalmente encontró muy poco. Estos estudios se basaban típicamente en análisis de series temporales, en los que los flujos comerciales bilaterales se relacionaban con medidas de variabilidad del tipo de cambio y otras variables de control. La mayoría de las veces esta relación resultaba débil e insignificante. Así, se concluyó que eliminar la variabilidad cambiaria en el marco de una unión monetaria tendría poco efecto sobre los flujos comerciales —para una revisión, véase Fondo Monetario Internacional, 1984—.

Existe una segunda generación de estudios econométricos, iniciada por Andy Rose, que llegó a conclusiones muy distintas. Utilizando datos de corte transversal y controlando por una multitud de otras variables que afectan los flujos comerciales —por ejemplo, renta, distancia, restricciones comerciales, idioma y muchas más—, Rose (2000) encontró que los pares de países que forman parte de una unión monetaria tienen entre sí flujos comerciales que, en promedio, son un 200 por ciento mayores que los de pares de países que no forman parte de una unión monetaria. Analizamos estos estudios en el Capítulo 2 y llegamos a la conclusión de que estos efectos comerciales de una unión monetaria están sobreestimados. Además, en ausencia de una buena teoría sobre cómo una unión monetaria impulsa el comercio, las estimaciones agregadas de la correlación entre unión monetaria y comercio son poco fiables —véase Baldwin, 2006—.

Varios estudios han intentado superar esta crítica examinando evidencia sectorial y microeconómica (Flam y Nordström, 2006; Baldwin et al., 2008; Berger y Nitsch, 2008; Nitsch y Pisu, 2008). Las estimaciones del efecto del euro sobre el comercio dentro de la Eurozona encontradas en estos estudios varían entre el 5 y el 20 por ciento. El mecanismo mediante el cual el euro ha impulsado el comercio se origina en el hecho de que la existencia del euro ha reducido los costes fijos y variables de las empresas exportadoras. Esto ha permitido que empresas que antes solo atendían mercados domésticos comenzaran a exportar a otros países de la Eurozona. Las pequeñas empresas en particular parecen haberse beneficiado de este efecto —véase Nitsch y Pisu, 2008—. Además, las empresas que ya exportaban han ampliado la gama de productos que venden en el extranjero.

Finalmente, en 2015 Glick y Rose publicaron una sorprendente mea culpa en la que admitieron que, utilizando datos posteriores a la Unión Monetaria Europea —UEM— y diversos métodos econométricos, “no encontramos ningún efecto sustantivo, fiable y robusto de la unión monetaria sobre el comercio” —véase Glick y Rose, 2015—. Sorprendentemente, los mismos autores publicaron un estudio en 2016 que parecía retractarse de esta mea culpa (Glick y Rose, 2016).

De esta breve revisión concluimos que la evidencia sobre un efecto positivo de las uniones monetarias en el comercio es débil.



\subsection{Costes y beneficios: una visión de conjunto}

En los capítulos anteriores, se identificaron los costes y beneficios de una unión monetaria. En este capítulo, concluimos esta discusión comparando los beneficios con los costes de una manera sintética. Esto nos permitirá evaluar la sensatez de los países de la Unión Europea (UE) cuando decidieron lanzar la Unión Económica y Monetaria (UEM), y los riesgos que asumieron. Además, hará posible extraer algunas conclusiones sobre la conveniencia económica de unirse a la UEM para los nuevos Estados miembros de la UE que están esperando entrar y para aquellos que hoy todavía dudan en hacerlo. También aplicaremos este análisis de costes y beneficios a otras partes del mundo: América Latina, Asia Oriental y África Occidental.

4.1 Comparación entre costes y beneficios

Es útil combinar las figuras (derivadas en los capítulos anteriores) que relacionan los beneficios y los costes con la apertura de un país. Esto se hace en la Figura 4.1. El punto de intersección de las líneas de beneficios y de costes determina el nivel crítico de apertura que hace que valga la pena para un país unirse a una unión monetaria con sus socios comerciales. A la izquierda de ese punto, el país está mejor manteniendo su moneda nacional. A la derecha está mejor cuando renuncia a su moneda nacional y la sustituye por la de sus socios comerciales.

La Figura 4.1 nos permite extraer algunas conclusiones cualitativas sobre la importancia de los costes y beneficios. La forma y la posición de la curva de costes dependen en gran medida de la visión que se tenga sobre la eficacia de las políticas monetarias nacionales, incluidas las políticas de tipo de cambio, para corregir los efectos de diferentes desarrollos en la demanda y en los costes entre los países implicados.

En un extremo, existe una visión, que será llamada “monetarista”, que sostiene que las políticas monetarias nacionales son ineficaces como instrumentos para corregir shocks asimétricos, sean permanentes o temporales. E incluso si son eficaces, según esta visión, el uso de estas políticas típicamente empeora la situación de los países. En esta visión “monetarista”, la curva de costes está muy cerca del origen. Representamos este caso en la Figura 4.2(a). El punto crítico que hace que valga la pena formar una unión está cerca del origen. Así, en esta visión, muchos países en el mundo obtendrían beneficios al renunciar a sus monedas nacionales y unirse a una unión monetaria.

En el otro extremo, está la visión “keynesiana” de que el mundo está lleno de rigideces (los salarios y los precios son rígidos; el trabajo es inmóvil), de modo que las políticas monetarias nacionales y el tipo de cambio son instrumentos poderosos para absorber shocks asimétricos. Esta visión está bien representada por el modelo original de Mundell discutido en el Capítulo 1. En esta visión, la curva de costes está lejos del origen, como se muestra en la Figura 4.2(b), de modo que relativamente pocos países deberían encontrar en su interés unirse a una unión monetaria. De esta visión también se sigue que muchos países grandes que actualmente tienen una sola moneda estarían mejor (económicamente) dividiendo el país en diferentes zonas monetarias.

Desde principios de la década de 1980, la visión “monetarista” ha ganado adeptos y ha cambiado la visión que muchos economistas solían tener sobre la conveniencia de una unión monetaria. La popularidad del monetarismo ayuda a explicar por qué la UEM se convirtió en una realidad en la década de 1990. Sin embargo, la crisis de deuda soberana que estalló en 2010 y las dificultades experimentadas en muchos países de la Eurozona para ajustarse a esta crisis parecen haber reducido esta popularidad en varios países. No se puede excluir que esto pueda llevar a una reevaluación de la conveniencia de una unión monetaria por parte de algunos de los miembros actuales de la Eurozona.

¿Qué nos enseña este análisis sobre la cuestión de si la UEM es un área monetaria óptima? Para responder a esta pregunta, primero presentamos algunos datos sobre la importancia del comercio intra-UE para cada país de la UE. Los datos se presentan en la Tabla 4.1, cuya característica más llamativa son las grandes diferencias en apertura entre los países de la UE. Esto conduce inmediatamente a la conclusión de que el cálculo de costes y beneficios es probable que produzca resultados muy diferentes para los distintos países de la UE. Para algunos países con un alto grado de apertura relativo a los otros socios de la UE, el cálculo de costes y beneficios es probable que muestre beneficios netos de estar en la UEM. Este es muy probablemente el caso en los países del Benelux, Austria, Irlanda y los nuevos Estados miembros de la UE, las Repúblicas Checa y Eslovaca, Estonia, Hungría y Eslovenia. Es muy llamativo encontrar que los países de Europa Central que se unieron a la UE en 2004 están al menos tan integrados con el resto de la UE como los países miembros más antiguos. Para los países en el otro extremo de la clasificación, el Reino Unido, Chipre y Grecia, está menos claro que pertenezcan a un área monetaria óptima con el resto de la UE.

Seguirá siendo difícil, sin embargo, trazar una línea precisa y concluir que aquellos por encima de la línea forman parte de un área monetaria óptima y aquellos por debajo no la forman con los otros Estados miembros de la UE. La razón es doble. En primer lugar, hay otros parámetros importantes que deben incorporarse al análisis, por ejemplo, el grado de flexibilidad y el grado de asimetría de los shocks. Volveremos sobre esta cuestión en las siguientes secciones. En segundo lugar, el grado de “completitud” de una unión monetaria importa para sus costes y beneficios. Retomamos esta cuestión en la sección 4.4.

Además, algunos países con una baja proporción de comercio pueden, sin embargo, encontrar ventajoso estar en una unión monetaria. Nuestro análisis de las cuestiones de credibilidad deja claro que países tradicionalmente de alta inflación, como Grecia o Italia, podrían haber decidido que era de su interés estar en la UEM a pesar del hecho de que su proporción de comercio con los miembros de la unión es relativamente baja. En términos del análisis de la Figura 4.2, esto implica que las autoridades griegas e italianas no consideraron costosa la pérdida de sus instrumentos de política monetaria nacional, de modo que la proporción mínima de comercio que hace ventajosa la unión es muy baja. Dicho de otro modo, si uno es suficientemente “monetarista”, podría argumentar que para países con bajos grados de apertura, los beneficios podrían aún superar los costes, y que estar en una unión monetaria también podría tener sentido para ellos desde un punto de vista económico.

4.2 Unión monetaria, rigideces de precios y salarios, y movilidad laboral

El cálculo de costes y beneficios de una unión monetaria también está muy influido por el grado de rigidez de salarios y precios. Como se recordará de nuestra discusión en los Capítulos 1 y 2, cuando los países se enfrentan a shocks asimétricos permanentes, que requieren cambios en los precios relativos, la pérdida del tipo de cambio puede ser un obstáculo en la medida en que hace que el ajuste a estos shocks sea más difícil. Como resultado, los países en los que el grado de rigidez de salarios y precios es bajo experimentan menores costes cuando se mueven hacia una unión monetaria. Mostramos esto en la Figura 4.3.

Una disminución en las rigideces de salarios y precios tiene el efecto de desplazar la curva de costes en la Figura 4.3 hacia abajo. Como resultado, el punto crítico en el que se vuelve ventajoso para un país renunciar a su moneda nacional se reduce. Más países se convierten en candidatos para una unión monetaria. De manera similar, un aumento en el grado de movilidad del trabajo desplaza la curva de costes hacia la izquierda y hace que una unión monetaria sea más atractiva. Sin embargo, debe señalarse que no todas las formas de integración tienen estos efectos. Como se destacó en el Capítulo 2, la integración económica también puede conducir a una mayor concentración regional de las actividades industriales. Esta característica del proceso de integración cambia el cálculo de costes y beneficios, en el sentido de que desplaza la curva de costes hacia la derecha y hace que una unión monetaria sea menos atractiva.

4.3 Shocks asimétricos y flexibilidad del mercado laboral

No es solo el grado de flexibilidad del mercado laboral (flexibilidad salarial y movilidad del trabajo) lo que importa para determinar si una unión monetaria será atractiva para los países. También son importantes el tamaño y la frecuencia de los shocks asimétricos a los que están sometidos. Esto significa que los países que experimentan shocks de demanda y oferta muy diferentes (porque sus estructuras industriales difieren mucho) encontrarán más costoso formar una unión monetaria. En el marco de la Figura 4.3 esto significa que la curva de costes se desplaza hacia la derecha.

Ahora estamos en condiciones de analizar la relación entre la flexibilidad del mercado laboral y los shocks asimétricos en una unión monetaria. Esto se hace gráficamente de la siguiente manera (véase la Figura 4.4). En el eje vertical representamos el grado de simetría entre las regiones (países) que son candidatos a formar una unión monetaria. Por simetría se entiende aquí el grado en que las tasas de crecimiento de la producción y del empleo están correlacionadas. Así, a medida que nos movemos hacia arriba verticalmente, la simetría aumenta, o dicho de otro modo, disminuye la medida en que ocurren shocks asimétricos. En el eje horizontal tenemos el grado de flexibilidad de los mercados laborales en estas regiones (países). La flexibilidad aquí se refiere a la flexibilidad salarial y a la movilidad interregional (internacional) del trabajo.

La idea central de la teoría de las áreas monetarias óptimas es que los países o regiones que experimentan mucha asimetría en el crecimiento de la producción y el empleo necesitan mucha flexibilidad en sus mercados laborales si quieren beneficiarse de la unión monetaria y si desean evitar grandes problemas de ajuste. Dicho de otro modo, cuanto menor sea el grado de simetría, mayor es la necesidad de flexibilidad en los mercados laborales para hacer posible una unión monetaria que funcione sin problemas. Esta relación entre simetría y flexibilidad está representada por la línea decreciente OCA. Muestra las combinaciones mínimas de simetría y flexibilidad que los países deben tener para que una unión monetaria proporcione más beneficios que costes. Los países o regiones situados por debajo de la línea OCA no tienen suficiente flexibilidad dado el nivel de simetría al que se enfrentan. Es probable que experimenten grandes costes de ajuste como resultado de shocks asimétricos. No forman un área monetaria óptima. Por lo tanto, es aconsejable que mantengan cierto grado de flexibilidad del tipo de cambio. Por supuesto, estos países siguen siendo libres de formar una unión monetaria. La teoría, sin embargo, predice que sufrirán económicamente por esta decisión. Por el contrario, los países a la derecha de la zona OCA tienen mucha flexibilidad dado el nivel de simetría al que se enfrentan. En otras palabras, podrán ajustarse a shocks asimétricos sin incurrir en grandes costes de ajuste. Los beneficios de una unión monetaria superan los costes para estos países. Forman un área monetaria óptima. Llamamos a la zona a la derecha de la línea OCA la zona OCA.

Continuación de la traducción estrictamente literal:

⸻

¿Dónde debería situarse la UE en la Figura 4.4? Desde 2013, la UE ha consistido en veintiocho países miembros. Diez nuevos países miembros se añadieron a la unión en mayo de 2004, dos más en 2007 y uno en 2013. Cuando el Reino Unido abandone la UE en 2020 habrá veintisiete miembros. La mayoría de los estudios empíricos han examinado esta cuestión desde el punto de vista de la UE compuesta por quince países miembros (UE-15). Existe un amplio consenso entre los economistas que han intentado aplicar empíricamente la teoría en que la UE-15 no es un área monetaria óptima (véase Eichengreen 1990; Neumann y von Hagen 1991; Bayoumi y Eichengreen 1993, 1997; De Grauwe y Heens 1993; De Grauwe y Vanhaverbeke 1990; Beine et al. 2003).

Algunos estudios se llevaron a cabo analizando la optimalidad de una unión monetaria que involucrara a la UE de veinticinco países miembros (UE-25) que existía antes de la entrada de Bulgaria y Rumanía en 2007 y Croacia en 2013 (véase Korhonen y Fidrmuc 2001). Estos estudios llegaron a la misma conclusión de que la UE-25 probablemente no era un área monetaria óptima. Así, según estos estudios empíricos, la UE-25 se situaba por debajo de la línea OCA. Como resultado, desde un punto de vista económico, una unión monetaria que involucre a todos los países miembros de la UE es una mala idea. Es probable que los costes económicos de una unión monetaria sean mayores que los beneficios para un número significativo de países. (En el Recuadro 4.1 presentamos un estudio de caso que ilustra la metodología utilizada en muchos estudios empíricos.)

Mientras que existe un fuerte consenso entre los economistas en que la UE-27 no debería formar una unión monetaria, existe una convicción igualmente fuerte de que hay un subconjunto de países de la UE que forman un área monetaria óptima. El conjunto mínimo de países para los cuales una unión monetaria es óptima se cree generalmente que incluye Alemania, los países del Benelux, Austria y Francia. Esta conclusión está respaldada por los mismos estudios empíricos citados anteriormente.

Algunos investigadores, sin embargo, han tendido a ampliar el grupo de países de la UE que se beneficiarían de la unión monetaria; véanse los estudios de Artis y Zhang (1997). Fidrmuc (2004), Melitz (2004) y Boeri y Garibaldi (2006) llegan a conclusiones similares.

Otros estudios empíricos han puesto en duda la visión núcleo-periferia de la integración monetaria en la UE. Erkel-Rousse y Melitz (1995) y Canzoneri et al. (1996a) encuentran que en la mayoría de los países de la UE las políticas monetarias son incapaces de afectar variables reales como la producción y el empleo. Así, incluso si los países de la UE se enfrentan a shocks asimétricos, sus instrumentos de política monetaria nacional no pueden utilizarse para hacerles frente de manera eficaz. Como resultado, la pérdida de estos instrumentos para la mayoría de los países de la UE no es muy costosa.

Más recientemente, Campos y Macchiarelli (2016), actualizando el análisis de Bayoumi y Eichengreen (1993), llegan a la conclusión de que la división núcleo-periferia en la Eurozona se ha debilitado desde la década de 1990, cuando Bayoumi y Eichengreen realizaron su análisis.

Finalmente, otra serie de estudios empíricos ha encontrado que una gran parte de los shocks asimétricos en los países de la UE ocurre a nivel sectorial y no tanto a nivel nacional. Dicho de otro modo, muchos de los cambios en la producción y el empleo en un país son el resultado de desarrollos diferentes entre sectores (por ejemplo, debido a desplazamientos de la demanda o a cambios tecnológicos diferenciales). Estos shocks no pueden ser afrontados mediante cambios en el tipo de cambio (véase Bini-Smaghi y Vori 1993; Bayoumi et al. 1995; y Gros 1996).

De este breve panorama de estudios empíricos, quedará claro que sigue siendo muy incierto cuán grande es el área monetaria óptima en Europa. Sin embargo, estos estudios empíricos no parecen socavar nuestra conclusión de que la UE-27 en su conjunto muy probablemente no constituye un área monetaria óptima. No existe consenso, sin embargo, sobre el tamaño del subconjunto de países que se beneficiarán de la unión monetaria. ¿Dónde deberíamos situar a la actual Eurozona, compuesta por diecinueve países miembros de la UE, en la Figura 4.4? Hemos situado a este grupo de países (Eurozona) por encima de la línea OCA, pero muchos economistas pueden objetar esto y situarían a la Eurozona por debajo de la línea OCA. La reciente crisis de la Eurozona da cierto apoyo a la opinión de que la Eurozona actual no es un área monetaria óptima y que, por lo tanto, debería situarse a la izquierda de la línea OCA. Volveremos sobre esta cuestión.

Debido a las muchas dificultades para cuantificar los costes de una unión monetaria, seguirá siendo difícil, sin embargo, obtener resultados claros en este ámbito. Esta situación genera un amplio margen para el juicio subjetivo.

En la Figura 4.4 también hemos situado a Estados Unidos por encima de la línea OCA. No estamos, por supuesto, realmente seguros de que Estados Unidos forme un área monetaria óptima. Sin embargo, estamos mucho menos inseguros sobre las posiciones relativas de la UE y Estados Unidos. Nótese que hemos situado a Estados Unidos aproximadamente al mismo nivel vertical que la UE-28, expresando el hecho de que el grado de simetría entre regiones en Estados Unidos no es muy diferente de la simetría observada entre países en la UE (véase Krugman 1993). La principal diferencia entre Estados Unidos y la UE parece ser el grado de flexibilidad de los mercados laborales. Muchos estudios empíricos han documentado esta diferencia en el grado de flexibilidad de los mercados laborales en Estados Unidos y en Europa. Por ejemplo, existe abundante evidencia de que los salarios reales en Europa responden menos al desempleo que en Estados Unidos. De manera similar, existe abundante evidencia de que la movilidad laboral es mucho mayor dentro de Estados Unidos que entre los países miembros de la UE (véase Beyer y Smets 2015 y Arpaia et al. 2016).

Debe subrayarse aquí que el análisis subyacente a la Figura 4.4 se basa en la teoría tradicional de las áreas monetarias óptimas. No aborda algunos de los problemas que discutimos en el Capítulo 2. Por ejemplo, los países pueden encontrar difícil, por razones de credibilidad, seguir políticas de baja inflación. La formación de una unión monetaria puede reducir estos problemas (véase Chari et al. 2015). Además, algunos de los shocks asimétricos que se observan hoy en Europa pueden ser el resultado de la ausencia de una unión monetaria. Por ejemplo, la naturaleza no sincronizada del ciclo económico entre la Eurozona y el Reino Unido puede deberse al hecho de que el Reino Unido sigue una política monetaria independiente de las políticas monetarias en el continente. Como resultado, no podemos estar realmente seguros de que la UE-27 no se beneficiaría de una unión monetaria. No obstante, con el estado actual de nuestro conocimiento, no es irrazonable mantener nuestra conclusión de que la UE-28 puede no ser un área monetaria óptima.

El desafío para la UE-27 es moverse hacia la zona OCA, es decir, hacer que una unión monetaria sea menos costosa. ¿Cómo puede lograrse esto? Existen esencialmente dos estrategias. Una consiste en reducir el grado de asimetría de los shocks (aumentar la simetría), y la otra consiste en aumentar el grado de flexibilidad.

La dificultad con la primera estrategia es que el grado de asimetría de los shocks depende en gran medida de factores sobre los cuales los responsables de la política tienen poca influencia; por ejemplo, el grado de especialización industrial importa para determinar cuán importantes son los shocks asimétricos. Sin embargo, los responsables de la política pueden hacer muy poco para cambiar los patrones de especialización regional. Existe, no obstante, un ámbito en el que pueden hacer algo para reducir el grado de asimetría. Este es el campo de la unificación política. Hemos argumentado anteriormente que una fuente importante de shocks asimétricos es la persistencia de los Estados-nación con sus políticas independientes de gasto e impuestos, y con sus propias particularidades en lo que respecta a las instituciones económicas. Para reducir los shocks asimétricos, será necesaria una mayor coordinación de las políticas económicas y una simplificación institucional. (Aquí surge un problema especial: ¿cómo deberían organizarse los sindicatos en una unión monetaria? Abordamos esta cuestión en el Recuadro 4.2.)



La otra estrategia para trasladar a toda la UE a la zona OCA consiste en aumentar el grado de flexibilidad de los mercados laborales (salarios reales y/o movilidad laboral). Esta estrategia implica una reforma de las instituciones del mercado de trabajo. Aunque tales reformas son difíciles de implementar, son necesarias si se quiere tener una unión monetaria que abarque a toda la UE (véase Campos et al. (2018)).

Volvemos a la cuestión de cómo moverse hacia el área OCA en la sección 4.5, donde introducimos otra dimensión de la unión política, es decir, la unión presupuestaria.

4.4 El grado de completitud de una unión monetaria

Argumentamos en el Capítulo 1 que el grado de completitud de las uniones monetarias importa para el coste de estas uniones. En una unión monetaria incompleta, es decir, una en la que no existe unión presupuestaria, habrá una gran fragilidad de los mercados de bonos gubernamentales. Además, en tales uniones es probable que los shocks asimétricos se intensifiquen por perturbaciones en los mercados de bonos gubernamentales que imponen altos tipos de interés en los países que experimentan shocks negativos, obligándolos a aplicar medidas de austeridad en medio de recesiones. En uniones monetarias que están integradas en una unión presupuestaria, estos problemas desaparecen y es probable que el coste de una unión monetaria sea menor. Presentamos estos dos tipos de unión monetaria en la Fig. 4.6.

El gráfico de la izquierda en la Fig. 4.6 presenta el caso de una unión monetaria incompleta. En esta unión, los costes de ajuste tras shocks asimétricos son más altos. Así, hemos dibujado la curva de costes más alejada del origen. El gráfico de la derecha presenta el caso de una unión monetaria completa donde, debido a la integración presupuestaria, la curva de costes está más cerca del origen.

El análisis anterior conduce a la predicción de que es más probable observar uniones monetarias que sean al mismo tiempo uniones presupuestarias. Una unión presupuestaria es un componente esencial de una unión política. Así, una combinación de unión monetaria y política tiene más probabilidades de tener menores costes y, por lo tanto, de funcionar mejor que las uniones monetarias que no están integradas en una unión política. De hecho, la evidencia abrumadora es que las uniones monetarias casi siempre están integradas en una unión política. La Eurozona es una gran excepción a esta regla. Volveremos a los problemas que genera una unión tan incompleta.

4.5 El trade-off entre unión presupuestaria y flexibilidad

El análisis anterior nos permite derivar una nueva línea OCA que muestra el trade-off entre flexibilidad y unión presupuestaria (Fig. 4.7). La idea de un trade-off entre flexibilidad y unión presupuestaria proviene de Andre Sapir de la Universidad de Bruselas (véase Sapir 2015). Por eso la llamamos línea OCAs. Nos permite obtener ideas interesantes sobre cómo los países pueden acercarse a una disposición monetaria óptima.

La forma en que se construye este trade-off es la siguiente. En el eje vertical representamos el grado de unión presupuestaria. Cuanto mayor es el grado de unión presupuestaria, más nos movemos hacia arriba a lo largo de la línea vertical. En el eje horizontal representamos la misma medida de flexibilidad que utilizamos en figuras anteriores. La línea OCAs mide ahora las combinaciones mínimas de unión presupuestaria y flexibilidad necesarias para que una unión monetaria sea económicamente atractiva (beneficios mayores que costes). Tiene pendiente negativa por la siguiente razón. Cuando la unión presupuestaria aumenta, el seguro frente a shocks asimétricos aumenta, haciendo que la unión monetaria sea menos costosa. Como resultado, hay menos necesidad de flexibilidad. Nos movemos hacia arriba a lo largo de la línea OCAs de pendiente negativa.

Esta es una idea importante. La flexibilidad puede sonar bien para muchos economistas y banqueros centrales. Sin embargo, es costosa para la mayoría de las personas que se ven obligadas a ser flexibles. La flexibilidad significa que estas personas pueden tener que aceptar una reducción salarial o verse obligadas a emigrar. Aprendemos de la Fig. 4.7 que un movimiento hacia la unión presupuestaria hace que la unión monetaria sea más aceptable para amplios segmentos de la población.

Podemos utilizar los conocimientos de la Fig. 4.7 para analizar la importancia de la naturaleza de los shocks asimétricos. Hemos hecho la distinción entre shocks asimétricos exógenos y endógenos. Los primeros son shocks permanentes como un aumento de los precios del petróleo; los segundos son el resultado de movimientos no sincronizados del ciclo económico y son temporales. Hemos argumentado que cuando ocurre un shock permanente (exógeno) la flexibilidad es la única opción para ajustarse a este shock. Por el contrario, cuando los movimientos del ciclo económico están desincronizados no es óptimo utilizar flexibilidad. En ese caso las transferencias fiscales son la respuesta adecuada. Esto es lo que proporciona una unión presupuestaria.

Ahora puede mostrarse que la naturaleza de los shocks influye en la pendiente del trade-off. Cuando los shocks son principalmente de tipo permanente, obtenemos un trade-off pronunciado. Mostramos esto en la Fig. 4.8. También hemos situado a la Eurozona de diecinueve miembros por debajo de la línea OCAs, sugiriendo que la Eurozona actual no es un área monetaria óptima. El trade-off pronunciado implica que un pequeño aumento en la flexibilidad nos lleva más rápidamente a la zona OCA que una unión presupuestaria. En el caso más extremo, es decir, cuando todos los shocks son de naturaleza permanente, el trade-off se vuelve vertical. En ese caso ninguna cantidad de unión presupuestaria nos llevará a la zona OCA. No hay otra manera que aumentar la flexibilidad.

Las cosas son muy diferentes cuando los shocks son temporales, impulsados por movimientos del ciclo económico. En ese caso el trade-off es plano (Fig. 4.6). Como resultado, se necesita mucha flexibilidad para trasladar la Eurozona al área OCA en comparación con la unión presupuestaria. Un aumento relativamente pequeño en la unión presupuestaria nos llevará a la zona OCA. En el caso más extremo, es decir, cuando todos los shocks son de naturaleza temporal, el trade-off es horizontal. En ese caso ninguna cantidad de flexibilidad logrará llevar la Eurozona a la zona OCA. La única manera de lograr la optimalidad será mediante una unión presupuestaria.

Una complicación que surge aquí tiene que ver con la histéresis. A veces los shocks temporales pueden dar lugar a efectos de histéresis. Por ejemplo, una recesión típicamente conduce al cierre de plantas y al despido de trabajadores. En la medida en que estos trabajadores han desarrollado habilidades específicas de la empresa que se pierden cuando la empresa desaparece, los trabajadores pierden parte de su capital humano, lo que dificulta encontrar otro empleo (comparable). El desempleo puede entonces prolongarse. Otro ejemplo se refiere a la naturaleza del auge. Si, como fue el caso en Irlanda y España, el auge se concentra en el mercado de la vivienda, muchos trabajadores son atraídos a este sector durante el auge. Después del colapso son despedidos. Pueden tener dificultades para utilizar las habilidades adquiridas en el mercado de la vivienda en otros sectores de la economía. Existe una amplia literatura sobre las fuentes de la histéresis (véase Blanchard et al. 1986; Ball 2009; Fatas y Summers 2016).

La existencia de histéresis tiene implicaciones para nuestra discusión. Implica que si ocurre un shock del ciclo económico es muy importante intentar utilizar políticas de estabilización para evitar efectos de histéresis. Si los shocks temporales del ciclo económico tienen efectos permanentes, la necesidad de establecer mecanismos que mitiguen el impacto de estos shocks se vuelve aún más importante.

Las Figuras 4.8 y 4.9 conducen a otra idea interesante. La flexibilidad en los mercados laborales es algo que los gobiernos nacionales pueden hacer. No hay necesidad de mayor integración para aumentar la flexibilidad. La unión presupuestaria, sin embargo, es de naturaleza diferente. Requiere integración política. En otras palabras, mientras que la flexibilidad está en el ámbito de los gobiernos nacionales, la unión presupuestaria es un asunto europeo (Sapir 2015). Así, cuando los shocks son permanentes, deben ser abordados a nivel nacional, mientras que cuando los shocks son temporales la respuesta debería estar al nivel de la Eurozona.

La cuestión crucial ahora es empírica. ¿Cuál ha sido la naturaleza de los shocks en la Eurozona desde su inicio? La respuesta a esta pregunta es importante ya que influye en nuestra visión sobre cuánto debemos confiar en la flexibilidad y en la unión presupuestaria (seguro) para hacer frente a los shocks en la unión monetaria. En el Capítulo 2 mostramos que los movimientos cíclicos temporales han sido muy importantes en la Eurozona. Estos movimientos cíclicos tendieron a estar fuertemente sincronizados. La asimetría se encontraba en las muy diferentes amplitudes de estos movimientos del ciclo económico entre países (véase Fig. 2.4). Algunos países como Grecia, Irlanda y España experimentaron auges y caídas muy fuertes, mientras que otros países experimentaron movimientos cíclicos relativamente suaves aproximadamente al mismo tiempo. En el Recuadro 4.3 profundizamos en este análisis empírico. En el Capítulo 7 analizamos las implicaciones de estos resultados para la futura gobernanza de la Eurozona en general y para la necesidad de una unión presupuestaria en particular.

4.6 Costes y beneficios a largo plazo

En secciones anteriores, discutimos los costes y beneficios de una unión monetaria. Nuestro análisis fue principalmente estático. Será útil añadir algo de dinámica a este análisis para obtener una mejor comprensión de la cuestión de cómo estos costes y beneficios de la unión monetaria pueden evolucionar a lo largo del tiempo.

Para analizar esta cuestión, utilizamos algunas de las herramientas introducidas en el Capítulo 2, donde discutimos la relación entre el grado de integración económica y la aparición de shocks asimétricos. Esta relación predice si el progreso hacia la integración económica conduce a la convergencia económica. Ahora añadiremos el análisis coste-beneficio.

En la Fig. 4.10 representamos el modelo. En el eje vertical representamos, como antes, el grado de simetría entre grupos de países. En el eje horizontal tenemos el grado de integración comercial entre los mismos grupos de países. La línea con pendiente positiva (TT) dice que, a medida que aumenta la integración comercial, aumenta el grado de simetría entre los países implicados, es decir, los países se vuelven más parecidos y se enfrentan a menos shocks asimétricos. (Llamamos a esto la visión de la Comisión Europea en el Capítulo 2. Volveremos a la visión de Krugman más adelante). La línea con pendiente negativa (llamada OCA) representa las combinaciones mínimas de simetría e integración comercial que hacen que la unión monetaria sea una operación de equilibrio (costes = beneficios).

Se deriva de una manera similar a la otra línea OCA que hemos utilizado en este capítulo. Tanto la integración como la simetría son cualidades de los países que aumentan los beneficios (reducen los costes) de una unión monetaria. Así, si, por ejemplo, aumenta la simetría (reduciendo los costes de una unión), estos países pueden permitirse tener menos integración (que aumenta los beneficios). Todos los puntos sobre la línea OCA son entonces combinaciones de simetría e integración para las cuales la unión monetaria tiene una ganancia neta cero. De ello se sigue que todos los puntos a la derecha de la línea OCA son puntos para los cuales los beneficios de la unión monetaria superan los costes. La llamamos la zona OCA.

En la Fig. 4.10 hemos situado a la UE-27 sobre la línea TT con pendiente positiva, a la izquierda de la línea OCA, reflejando la conclusión a la que llegamos en la sección 4.3, es decir, que hoy (en 2020) la UE-27 probablemente no es un área monetaria óptima. Sin embargo, a medida que avanza la integración comercial dentro de la UE, este punto se moverá hacia arriba a lo largo de la línea TT. Esto inevitablemente nos llevará a la zona OCA, al menos si podemos suponer que la dinámica de integración continuará funcionando dentro de la UE. Así, en esta visión, la unificación monetaria será percibida con el tiempo como beneficiosa para todos los países de la UE. En este sentido, la unión monetaria entre todos los países de la UE es inevitable.

Esta es la visión optimista sobre las perspectivas a largo plazo de la integración monetaria en Europa. Sin embargo, existe también una visión pesimista, que puede derivarse del análisis de Krugman. Se recordará que en el análisis de Krugman la integración económica conduce a menos simetría entre países. Esto se representa en la Fig. 4.11 mediante las líneas TT y TT con pendiente negativa.

Ahora tenemos que considerar dos posibilidades para las perspectivas a largo plazo de la unión monetaria. Una está representada por la línea TT, cuya pendiente es más plana que la pendiente de la línea OCA. En este caso, aunque hoy la UE-27 puede no ser un área monetaria óptima, se moverá hacia la zona OCA con el tiempo. En este caso, más integración conduce a más especialización y, por tanto, a más shocks asimétricos (menos simetría). Sin embargo, los beneficios de una unión monetaria también aumentan fuertemente con el grado de integración. Como resultado, pese a la disminución de la simetría, más integración nos llevará a la zona OCA.

El segundo caso está representado por la línea TT pronunciada. Aquí la integración nos lleva cada vez más lejos de la zona OCA. Esto es así porque las ganancias netas de una unión monetaria no aumentan lo suficientemente fuerte con el grado de integración. Como resultado, los costes de la disminución de la simetría superan todos los demás beneficios que una unión monetaria pueda tener. A largo plazo, las perspectivas para una unión monetaria de la UE-27 son pobres. Debe señalarse que este caso conduce a una anomalía. Implica que una reducción de la integración comercial puede llevarnos a la zona OCA. Así, si los países de la UE retrocedieran y se desintegraran, la unión monetaria se volvería atractiva. Un resultado extraño.

De la discusión del modelo de Krugman, concluimos que incluso si la integración conduce a más shocks asimétricos, esto todavía puede conducir a ganancias netas crecientes de una unión monetaria para la UE-27.

Un punto final pero importante sobre la dinámica a largo plazo de la integración monetaria es el siguiente. La decisión de un país individual de unirse a la unión monetaria puede acelerar el proceso de integración. En el Capítulo 2 documentamos esto refiriéndonos a estudios empíricos que indican cómo la existencia de diferentes monedas nacionales ayuda a segmentar los mercados nacionales. Una decisión de un país individual de unirse a la UEM, incluso si no satisface los criterios OCA, tendría un carácter autocumplido, al menos si la Fig. 4.10 es la visión correcta del mundo. En este caso, el proceso de integración se aceleraría por la propia decisión de unirse a la unión monetaria, de modo que esta nueva agrupación de países se mueve más rápido hacia la zona OCA. En este sentido, puede decirse que existe un componente endógeno en los criterios de un área monetaria óptima: es decir, cuando los países deciden formar una unión monetaria, los criterios se vuelven más favorables. Dicho de otro modo, la propia decisión de formar una unión monetaria cambia el resultado coste-beneficio y reduce los costes en relación con los beneficios. A la inversa, al no formar una unión monetaria, un balance coste-beneficio desfavorable sigue siendo desfavorable. La naturaleza endógena de los criterios de un área monetaria óptima ha sido analizada por Frankel y Rose (1998).

¿Existe alguna evidencia de que la dinámica endógena haya estado funcionando dentro de la Eurozona? Dado que la Eurozona existe desde 1999, uno puede ver ya esta dinámica en funcionamiento. Consideramos aquí los tres criterios OCA: integración económica, simetría y flexibilidad. Revisamos la literatura sobre los efectos del euro en la integración comercial en el Capítulo 3, y llegamos a la conclusión de que la existencia del euro puede haber añadido entre un 5 y un 20 por ciento al comercio intra-Eurozona, aunque investigaciones más recientes ponen en duda esta conclusión (Glick y Rose 2015). Esto sugiere que el criterio de integración no se ha estado moviendo mucho para acercar a los países miembros a la zona OCA (o para introducirlos más profundamente en ella).

¿Qué ocurre con la asimetría de los shocks? Respondemos a esta pregunta en la Fig. 4.12. La figura presenta los coeficientes medios de correlación bilateral de los componentes del ciclo económico de la producción industrial dentro de la Eurozona durante diferentes períodos. [Nota al pie: El procedimiento consiste primero en extraer el componente cíclico de la serie de producción industrial utilizando un filtro HP. Estos componentes cíclicos se correlacionan luego entre pares de países. A continuación se calcula el coeficiente medio de correlación. Debe señalarse que este procedimiento es sensible al parámetro de suavizado utilizado en el filtro HP. Como resultado, los niveles de los coeficientes de correlación obtenidos pueden variar significativamente. Esto explica por qué los coeficientes de correlación en la Figura 4.12 son inferiores a los obtenidos en la Tabla 4.3.]

Observamos lo siguiente. Durante el período previo a la crisis (1999-2007), el coeficiente medio de correlación es ligeramente menor que en el período previo a la Eurozona (1987-98). Durante el período posterior a la crisis (2008-14) la correlación aumenta. Esto puede sugerir que en el período posterior a la crisis la convergencia del producto ha aumentado. Sin embargo, es probable que esto se deba principalmente a la fuerte caída de la producción que ocurrió en 2009 y que afectó a todos los países simultáneamente como resultado de la “Gran Recesión”. En un trabajo reciente, no obstante, Campos et al. (2017) realizaron un metaanálisis de sesenta y tres artículos que estudian la correlación de los ciclos económicos en la Eurozona. Sus principales resultados fueron que (1) los coeficientes de correlación de los ciclos económicos aumentaron de aproximadamente 0,42 antes del euro a aproximadamente 0,6 después; (2) este aumento ocurrió tanto en países del euro (núcleo y periferia) como en países no pertenecientes al euro. Así, parece que el grado de correlación de los movimientos del ciclo económico puede haber aumentado desde el inicio de la Eurozona.

El tercer criterio OCA que se destacó en nuestro análisis anterior fue la flexibilidad. Aquí la evidencia apunta a un aumento significativo de la flexibilidad del mercado laboral (véase Fig. 4.13). Esta muestra la evolución de un indicador de rigideces del mercado laboral, es decir, el índice de legislación de protección del empleo elaborado por la OCDE. Este mide el grado en que la legislación dificulta a las empresas despedir trabajadores. Observamos en la Fig. 4.13 que la Eurozona ha progresado en la flexibilización de su legislación de protección del empleo, especialmente desde el inicio de la crisis de deuda soberana de 2010. La Eurozona sigue situándose por encima de la media de los países de la OCDE en lo que respecta a la protección del empleo, pero la brecha se ha reducido considerablemente. Los países más afectados por la crisis de deuda soberana (Irlanda, Grecia, Portugal y España) fueron los que más avanzaron en la reducción de la protección del empleo. Estos países también han impuesto condiciones más estrictas para recibir prestaciones por desempleo y se ha introducido mayor flexibilidad en los horarios de trabajo (véase Comisión Europea 2008, pp. 77-80).

En conjunto, la evidencia de que existe un componente endógeno en los criterios OCA recibe solo un apoyo mixto. El criterio de integración no muestra mucha evidencia de moverse en la dirección correcta, mientras que el criterio de flexibilidad parece estar avanzando en la dirección OCA. En cuanto al criterio de simetría, parece haber evidencia de que la simetría puede haber aumentado.

















\clearpage
\section{Conclusión} 


\subsection{Recapitulación}


\subsection{Extensiones y relación con otras partes del Temario}



\subsection{Opinión}

\subsubsection{Idea final (Salida o cierra)}

\clearpage
\pagenumbering{gobble}
\MyBibliography



\MyBibItem{DAWID y GATTI}{Chapter 2 - Agent-Based Macroeconomics}{2018}{https://doi.org/10.1016/bs.hescom.2018.02.006}


% ANEXOS
\clearpage
\anexo{\textbf{Preguntas Test}}
%\input{\TemaCode/questions.tex} 


\clearpage
\anexo{Ampliación: Crítica a la Teoría de las Areas Monetarias Óptimas}\label{anx:critica_tamo}
En el presente anexo se amplían las críticas (o contracríticas) que se han expuesto sobre la Teoría de las Áreas Monetarias Óptimas.

$$\text{\textbf{1. Presencia de shocks asimétricos}}$$
En el apartado correspondiente vimos como MUNDELL enfatizaba la importancia de contar con mecanismos autónomos de ajuste a los shocks por los efectos perversos que se podrían derivar de la presencia de shocks asimétricos. También vimos qué puede ser fuente de estos shocks y cómo pueden agravarse atendiendo a las dinámicas de la deuda. 

Como crítica, observamos el argumento de la Comisión Europea que enfatizaba como la integración (económica y monetaria) hacía que estos shocks asimétricos fueron cada vez menos frecuentes y las razones que argumentaban al respecto (p.e. la importancia del comercio intraindustrial). 

No obstante, existe otro argumento (contracítica) propuesto por KRUGMAN (1991). Según este autor, el análisis de MUNDELL no puede descartarse ya que existe una característica implícita en la dinámica del comercio con economías de escala que puede hacer dicho análisis muy relevante: \textbf{la integración comercial, que surge como resultado de las economías de escala, también conduce a concentraciones regionales de la actividad industrial}.\footnote{Esta es una idea antigua desarrollada por MYRDAL (1957) y KALDOR (1966). Para una revisión, véase BALASSA (1961). KRUGMAN (1991) proporciona una fundamentación más rigurosa.
}

El argumento básico es que cuando disminuyen los obstáculos al comercio, se generan dos efectos opuestos sobre la localización de la producción. Por un lado, facilita producir cerca de los mercados finales; por otro, permite concentrar la producción para aprovechar economías de escala. Esto explica por qué la integración comercial puede conducir a una mayor concentración regional de la actividad económica.

El hecho de que el comercio pueda generar concentración regional se ilustra comparando la distribución geográfica de la producción de automóviles en Estados Unidos y en Europa. La característica más destacada es que la producción en Estados Unidos está mucho más concentrada regionalmente que en la Unión Europea. Este patrón se observa en muchos otros sectores industriales (véase KRUGMAN, 1991). Aunque el sector del automóvil ha cambiado desde entonces, esta diferencia en la concentración regional se ha mantenido (KLIER y RUBENSTEIN, 2013). La mayoría de las plantas de ensamblaje en EE. UU. siguen localizadas en el Medio Oeste, con cierta expansión hacia el sur. En Europa, por el contrario, la expansión hacia el este ha mantenido (e incluso aumentado) la dispersión geográfica.

\imagenfit{Table2_1.png}{Distribución regional de la producción de automóviles}{KRUGMAN (1991)}

No hay duda de que el mercado estadounidense está más integrado que el europeo, es decir, presenta menos barreras al comercio interno. La evidencia sugiere que, a medida que la UE avance hacia un mercado más integrado, podría experimentar patrones similares de concentración regional. Esto implica que shocks sectoriales podrían convertirse en shocks específicos de país. En tales casos, los países podrían preferir utilizar el tipo de cambio como instrumento de ajuste.

Estas dos visiones pueden representarse gráficamente. En el eje vertical se mide el grado de simetría de los movimientos de producción y empleo entre países (por ejemplo, mediante la correlación de sus tasas de crecimiento). En el eje horizontal se mide el grado de integración comercial (por ejemplo, el comercio bilateral como proporción del PIB).

Según la Comisión Europea, la relación es positiva: a mayor integración, mayor simetría. Según la visión de KRUGMAN, la relación es negativa: a mayor integración, mayor especialización y, por tanto, más shocks asimétricos.\footnote{Esta visión se relaciona con KENEN (1969), quien argumenta que economías menos diversificadas son más vulnerables a shocks asimétricos. Esto genera la paradoja de que países pequeños y abiertos deberían mantener su moneda (FRANKEL y ROSE, 1998).
}

¿Cuál de estas visiones es correcta? No existe una respuesta concluyente. Sin embargo, puede argumentarse que existe una presunción a favor de la visión de la Comisión Europea. Aunque la integración puede generar concentración, también reduce la importancia de las fronteras nacionales en la localización de la actividad económica. Esto aumenta la probabilidad de que los clusters productivos atraviesen fronteras. Por ejemplo, si la producción de automóviles se concentra en una región que abarca el sur de Alemania y el norte de Italia, un shock afectará a ambos países simultáneamente, haciendo inútil el tipo de cambio bilateral como mecanismo de ajuste.

El argumento no es que la integración no genere concentración, sino que las fronteras nacionales pierden relevancia en la configuración de esa concentración. Por tanto, aunque pueden existir shocks asimétricos a nivel regional, es más probable que estos no coincidan con fronteras nacionales. En consecuencia, la integración económica reduce la eficacia del tipo de cambio como instrumento de ajuste.

No obstante, esta cuestión sigue siendo esencialmente empírica. Una primera dimensión empírica es si las uniones monetarias aumentan la integración económica. Estudios de ROSE (2000), ROSE y VAN WINCOOP (2001) y GLICK y ROSE (2002), utilizando datos de panel para numerosos países, encuentran que pertenecer a una unión monetaria puede duplicar el volumen de comercio bilateral, incluso tras controlar por otros factores.

Estos resultados han generado una amplia literatura, parte de la cual ha criticado los hallazgos de Rose por razones econométricas (por ejemplo, Nitsch, 2001; Persson, 2001; Baldwin).

Otros estudios se han centrado en los efectos comerciales de la integración monetaria en Europa (véase Bun y Klaasen, 2002; Micco et al., 2003; De Nardis y Vicarelli, 2003; Flam y Nordström, 2003; Berger y Nitsch, 2006; Glick y Rose, 2016). La conclusión de estos trabajos es que el inicialmente fuerte “efecto Rose” probablemente está sesgado al alza y que el efecto del comercio derivado de la unión monetaria en Europa es mucho menor, situándose en torno al 5-20 por ciento. Dicho de otro modo, la unión monetaria en Europa podría generar una expansión del comercio de entre el 5\% y el 20\%. No obstante, incluso esta estimación es incierta, como reconocen Glick y Rose (2015).

El segundo aspecto empírico se refiere a cómo esta mayor integración afecta a la asimetría de los shocks. FRANKEL y ROSE (1998) llevaron a cabo un análisis empírico relevante sobre esta cuestión. Estudiaron el grado en que la actividad económica entre pares de países está correlacionada en función de la intensidad de sus vínculos comerciales. Concluyen que una mayor integración comercial entre dos países está fuertemente asociada con una mayor correlación de su actividad económica. Este resultado es confirmado también por ROSE (2000), ROSE y ENGEL (2001), BAXTER y KOUPARITSAS (2005), y ROSE y YI (2006). De manera similar, ARTIS y ZHANG (1997) muestran que, a medida que los países europeos se integraron durante los años ochenta y noventa, sus ciclos económicos se volvieron más correlacionados. 

Existe otra evidencia empírica que refuerza la idea de que la integración económica no necesariamente genera más shocks asimétricos. Esta se relaciona con la creciente importancia del sector servicios. Las economías de escala parecen menos relevantes en los servicios que en la industria. Como resultado, la integración no conduce a una concentración regional de los servicios en la misma medida que en la producción industrial. Dado que los servicios representan hoy más del 70\% del PIB en muchos países de la UE, la tendencia hacia la concentración regional podría detenerse incluso si la integración continúa avanzando. Existe evidencia de que esto ya ocurre en Estados Unidos. La OCDE (1999a) concluyó que la concentración regional en EE. UU. comenzó a disminuir tras décadas de aumento. Este resultado también es confirmado por Aiginger y Leitner (2002), quienes concluyen que la concentración regional está disminuyendo tanto en EE. UU. como en la UE.

En resumen, la evidencia sugiere, primero, que la unión monetaria intensifica la integración comercial (aunque la magnitud del efecto es incierta) y, segundo, que dicha integración reduce la frecuencia de shocks asimétricos. Esto constituye una conclusión favorable para los países que deciden formar una unión monetaria: el propio proceso puede generar condiciones que favorezcan su funcionamiento. Volveremos sobre esta dinámica en el Capítulo 4.

$$\text{\textbf{ } Teoría de las Areas Monetarias Óptimas: Crítica}$$
Hemos analizado las razones por las cuales los países pueden encontrar costoso incorporarse a una unión monetaria. Este análisis, conocido como la teoría de las áreas monetarias óptimas (OCA), ha sido objeto de críticas. Estas críticas se han formulado en distintos niveles. En primer lugar, puede cuestionarse la idea de que las diferencias entre países sean lo suficientemente importantes como para justificar preocupación. En segundo lugar, el uso de políticas monetarias nacionales, incluido el tipo de cambio, puede no ser muy eficaz para corregir dichas diferencias. En tercer lugar, no solo es posible que las políticas monetarias y cambiarias sean ineficaces, sino que también pueden generar más perjuicios que beneficios en manos de los responsables políticos.


$$\text{\textbf{} ¿Qué tan relevantes son las diferencias entre países?}$$
No hay duda de que los países son diferentes. La cuestión es si estas diferencias son lo suficientemente importantes como para constituir un obstáculo para la unificación monetaria. ¿Es probable que se produzcan shocks de demanda concentrados en un solo país?

El análisis clásico de Mundell parte del escenario en el que se produce un desplazamiento de la demanda desde los productos de un país hacia los de otro. ¿Es probable que este tipo de shock ocurra con frecuencia entre los países europeos que forman una unión monetaria? Existen dos enfoques para responder a esta cuestión. Denominaremos al primero la visión de la Comisión Europea, defendida en el informe de 1990 One Market, One Money. El segundo enfoque está asociado a Paul Krugman y puede consultarse en [\authorfont{Ver Anexo \ref{anx:critica_tamo}}]

Según la Comisión Europea, los shocks diferenciales de demanda serán menos frecuentes en una unión monetaria. La razón es que el comercio entre los países industriales europeos es en gran medida intraindustrial. Este comercio se basa en la existencia de economías de escala y competencia imperfecta (diferenciación de productos). Da lugar a una estructura en la que los países compran y venden entre sí los mismos tipos de bienes. Así, Francia vende automóviles a Alemania y compra automóviles de Alemania, y viceversa. Esta estructura implica que la mayoría de los shocks de demanda afectarán a estos países de manera similar. Por ejemplo, si los consumidores reducen su demanda de automóviles, comprarán menos automóviles tanto franceses como alemanes. Por tanto, la demanda agregada de ambos países se verá afectada de manera similar. La eliminación de barreras con la culminación del mercado único reforzará estas tendencias. Como resultado, los shocks tenderán a ser más simétricos que asimétricos.

No obstante, cabe destacar que los resultados aún no son concluyentes. La evidencia sugiere que la unión monetaria intensifica la integración comercial (aunque la magnitud del efecto es incierta) y, segundo, que dicha integración reduce la frecuencia de shocks asimétricos (pero no elimina con certeza). Esto constituye una conclusión favorable para los países que deciden formar una unión monetaria: el propio proceso puede generar condiciones que favorezcan su funcionamiento. Volveremos sobre esta dinámica en el Capítulo 4.










$$\text{\textbf{} Auge y caída y el Estado-nación}$$
Anteriormente se argumentó que la integración económica reduce la probabilidad de shocks asimétricos entre países. Sin embargo, existen shocks endógenos derivados de la dinámica del capitalismo, caracterizada por ciclos de optimismo y pesimismo (animal spirits) que generan auges y caídas.

Cuando estos ciclos están desincronizados, pueden surgir problemas importantes en una unión monetaria. En la Eurozona, los ciclos han estado relativamente correlacionados, pero con amplitudes muy diferentes. Países como Irlanda, Grecia y España experimentaron ciclos mucho más intensos que otros como Alemania o Austria.

\imagenfit{Fig2_4.png}{Ciclo de negocios como componente del PIB: Eurozona}{}

Como vemos, la asimetría no se encuentra tanto en la dirección de los ciclos como en su intensidad: estas diferencias generaron efectos divergentes sobre las finanzas públicas y provocaron fuertes salidas de capital en los países más afectados durante las crisis. Entonces, ¿por qué estas divergencias ocurrieron y por qué la existencia de una moneda común no disciplinó estos animal spirits? La respuesta se encuentra en el carácter incompleto de la UM. 

En la UEM, las políticas monetarias están centralizadas y, por tanto, dejan de ser una fuente de shocks asimétricos. Sin embargo, los países miembros de la UM siguen ejerciendo una soberanía considerable en varias áreas económicas, la más importante siendo el ámbito presupuestario. En la UM, la mayor parte de las competencias de gasto y tributación sigue estando en manos de las autoridades nacionales.\footnote{Hoy, en la mayoría de los países de la UE, el gasto y la tributación de las autoridades nacionales ascienden a cerca del 50 por ciento del PIB. Las competencias de gasto y tributación de las autoridades europeas representan alrededor del 1\% del PIB. Esta situación no ha cambiado desde el inicio de la unión monetaria en 1999.
} Al modificar impuestos y gasto, las autoridades de un país individual pueden generar grandes shocks asimétricos. Por su propia naturaleza, estos shocks quedan bien contenidos dentro de las fronteras nacionales. Por ejemplo, cuando las autoridades de un país aumentan los impuestos sobre la renta salarial, esto solo afecta al trabajo en ese país e influirá sobre el gasto y los niveles salariales en ese país. Como resultado, las curvas de demanda y oferta agregadas del país afectado se desplazarán, creando perturbaciones que conducirán a desarrollos divergentes de precios y salarios.

El hecho de que los países mantengan la mayor parte de sus competencias de gasto y tributación dentro de la unión monetaria crea la posibilidad de que se produzcan grandes desarrollos asimétricos en la unión, lo que plantea la cuestión de cómo deberían conducirse las políticas presupuestarias en una unión monetaria.

Existen otros aspectos de la existencia de los Estados-nación que pueden ser fuente de perturbaciones asimétricas y que pueden inducir una dinámica nacional de auges y recesiones. Como se argumentó antes, muchas instituciones económicas son nacionales y el efecto de estas diferencias puede ser que los países experimenten condiciones económicas muy diferentes, que pueden reforzar condiciones de auge en un país, mientras otros países experimentan un crecimiento mucho más moderado de la actividad económica.\footnote{Los sistemas de negociación salarial, por ejemplo, difieren ampliamente entre países, creando la posibilidad de desarrollos asimétricos. Además, las diferencias en los sistemas legales y en las costumbres generan diferencias significativas en el funcionamiento de los mercados financieros. Por ejemplo, las regulaciones sobre las condiciones bajo las cuales los bancos conceden hipotecas difieren de un país a otro. Además, las diferencias en las normas fiscales que determinan la cantidad que los prestatarios pueden deducir de su renta pueden conducir a evoluciones muy distintas de los precios de la vivienda en distintos países. Pueden darse muchos más ejemplos. 
} Así, la persistencia de Estados-nación con sus propias políticas e instituciones también creó \textit{animal spirits} nacionales, es decir, oleadas de optimismo y pesimismo que, aunque correlacionadas en toda la Eurozona, fueron muy diferentes en intensidad y condujeron a flujos de capital desestabilizadores durante la fase descendente.

Estas diferencias en los movimientos de la actividad económica pueden dar lugar a otro fenómeno importante: la aparición de grandes divergencias en las posiciones competitivas de los miembros de una unión monetaria (exactamente lo que ocurrió en la Eurozona):\footnote{El coste laboral unitario relativo se define como el coste laboral unitario de un país (por ejemplo, Alemania) en relación con el coste laboral unitario medio de los demás países miembros de la Eurozona. 
}

\imagenfit{Fig2_5.png}{Coste relativo de trabajo unitario en los países de la Eurozona}{}

Cuando el coste laboral unitario relativo disminuye, como en el caso de Alemania, puede decirse que los costes laborales unitarios alemanes disminuyeron en relación con los de los demás países miembros. Como resultado, Alemania mejoró su posición competitiva durante el período 2000-08, principalmente porque experimentó condiciones de auge muy moderadas. Por el contrario, Irlanda, Grecia y España, que experimentaron los auges más intensos de la actividad económica, vieron aumentar sus costes laborales unitarios relativos respecto de los demás países miembros. Como resultado, perdieron competitividad durante $2000-08$ lo que, con la aparición de las crisis financieras en $2008$, condujo a movimientos inversos en varios países que habían perdido competitividad. 

Obsérvese que los costes laborales unitarios se definen como costes laborales corregidos por la productividad del trabajo.\footnote{El coste laboral unitario se define como: $CLU = W/(Q/L)$, donde $W$ es el salario, $Q$ es el valor de la producción y $L$ es la cantidad de trabajo utilizada en la producción. $Q/L$ es la productividad media del trabajo. Obsérvese que la fórmula también puede reescribirse como $CLU = WL/Q$. Esta expresión deja claro qué significa el coste laboral unitario: es el coste salarial incorporado en un euro (o dólar) de producción.
}

De ello se sigue que los costes laborales unitarios pueden aumentar por dos razones: (i) cuando los salarios aumentan; o, (ii) cuando la productividad del trabajo disminuye. Así, parece que países como Italia, España, Portugal e Irlanda han perdido una competitividad significativa desde $2000$ porque los salarios en estos países aumentaron más rápido que la productividad. Esto ha generado graves problemas de ajuste que obligaron a iniciar un proceso de reducción de sus niveles salariales relativos respecto de los demás países de la Eurozona (si no pueden aumentar la productividad): la devaluación interna, proceso lento y doloroso. En un escenario distinto, si estos países no estuvieran en una unión monetaria, habrían podido devaluar sus monedas, restaurando así su competitividad de una sola vez. 

Del análisis de esta sección y de las anteriores podemos concluir lo siguiente. Aunque la integración económica probablemente produce estructuras económicas que pueden conducir a una mayor convergencia, la persistencia de los Estados-nación en las uniones monetarias puede poner en marcha una dinámica macroeconómica que conduce a fuertes divergencias. Así, pese a que a nivel microeconómico la integración comercial puede generar menos shocks asimétricos, la ausencia de una unión política puede desencadenar una dinámica macroeconómica que conduce a grandes desarrollos asimétricos. Dado que la unión monetaria es incompleta, la dinámica de auges y recesiones sigue teniendo un fuerte carácter nacional, creando difíciles problemas de ajuste para los Estados miembros. Por tanto, existe la necesidad de insertar una unión monetaria en una unión política más fuerte. 

$$\text{\textbf{} ¿Qué tan eficaces son las políticas monetarias nacionales?}$$
El coste de renunciar a la moneda nacional reside en que un país ya no puede utilizar políticas monetarias nacionales y, en particular, el tipo de cambio, para corregir desarrollos diferenciales en la demanda, o en costes y precios. Pero ¿qué tan eficaces son estos instrumentos nacionales de política monetaria?

Esta es una cuestión crucial. Si la respuesta es que estos instrumentos son ineficaces, podría concluirse que los países, incluso si presentan diferencias importantes entre sí, no tendrían que afrontar costes adicionales al incorporarse a una unión monetaria. El instrumento que pierden no les permite realmente corregir esas diferencias.

Para analizar esta cuestión sobre la eficacia de las políticas monetarias nacionales, volvemos a las perturbaciones asimétricas analizadas.

$$\text{\textbf{} Políticas monetarias nacionales para corregir shocks de demanda asimétricos permanentes}$$

Tomemos el caso de Francia, suponiendo que se poduce un desplazamiento de la demanda desde los productos franceses hacia los productos alemanes. Para hacer frente a este problema, Francia puede utilizar políticas monetarias expansivas, que conducirán a una depreciación de la moneda, o alternativamente, si tiene un tipo de cambio fijo con Alemania, puede devaluar su moneda. Aquí nos concentramos en este último caso. Primero analizamos el caso de un shock de demanda permanente y luego pasamos a los shocks temporales.

En el caso de un shock de demanda permanente, \textbf{los precios y costes de los bienes franceses tendrán que disminuir en relación con los de Alemania} si Francia quiere volver a su nivel inicial de producción (incluidos los costes salariales). Dicho de otro modo, cualquiera que sea el régimen monetario utilizado, el precio de los bienes franceses tendrá que disminuir respecto al de los bienes alemanes para estimular la demanda de bienes franceses y restaurar así el nivel inicial de producción. ¿Puede una depreciación del franco francés lograr este resultado de forma permanente? 

\imagenfit{Fig2_6.png}{Efectos sobre los costes y los precios de una depreciación de la moneda}{}

Como resultado de la depreciación, la demanda agregada en Francia se desplaza de nuevo hacia arriba y corrige el desplazamiento desfavorable inicial de la demanda, alcanzando el punto inicial de equilibrio. Así, el nivel de precios de los bienes franceses (en francos franceses) ha sido restaurado. Dado que el franco francés se ha depreciado frente al marco alemán, el precio en marcos alemanes de los bienes franceses es ahora menor, mientras que el precio en francos franceses de los bienes alemanes es ahora mayor. Por tanto, parece que la devaluación del franco francés logra reducir el precio de los bienes franceses en relación con los bienes alemanes. Las cosas parecen ir en la dirección correcta.

Sin embargo, es improbable que este nuevo punto de equilibrio pueda sostenerse. La razón es que la depreciación eleva el precio en francos franceses de los bienes alemanes importados, lo que: (i) eleva directamente el coste de producción en Francia; (ii) reduce el salario real de los trabajadores franceses; y, (iii) ejercerá presión al alza sobre el salario nominal en Francia, ya que es probable que los trabajadores quieran compensación por la pérdida de poder adquisitivo. Todo esto significa que la curva de oferta agregada se desplazará hacia arriba. Así, el nivel de precios francés aumenta y la producción disminuye. Este aumento de precios retroalimenta nuevamente el proceso de formación salarial y conduce a nuevos desplazamientos hacia arriba de la curva de oferta agregada. 


En definitiva, el equilibrio final se situará en un punto tal que los efectos inicialmente favorables de la depreciación del franco francés sobre la producción francesa desaparezcan con el tiempo: los efectos inicialmente favorables de la depreciación del franco francés sobre el precio relativo de los bienes franceses tienden a disiparse. Al pasar de $F\to F’$, el precio relativo de los bienes franceses colverá a aumentar. De todas formas, tampoco será posible decir si los efectos inicialmente favorables sobre la producción y los precios relativos desaparecerán por completo, ya que dependerá del grado de apertura de la economía y del grado en que los asalariados ajusten sus reivindicaciones salariales para corregir la pérdida de poder adquisitivo.\footnote{Esto último también depende de rasgos institucionales de los mercados de trabajo y de productos (véase HANCKE, 2013). Existe abundante evidencia empírica, sin embargo, de que para la mayoría de los países europeos esta desaparición de los efectos inicialmente favorables de una depreciación será fuerte, aunque no siempre completa.
}

La conclusión anterior también puede formularse del siguiente modo: \textbf{las variaciones del tipo de cambio nominal tienen efectos temporales sobre los precios relativos}. Con el tiempo, las depreciaciones nominales conducen a aumentos de costes y precios internos, que tienden a restaurar los precios relativos iniciales. En otras palabras, las depreciaciones nominales conducen a depreciaciones reales temporales. A largo plazo, las variaciones del tipo de cambio nominal dejan el tipo de cambio real de un país en gran medida inalterado.

¿Implica esta conclusión sobre la limitada ineficacia a largo plazo de las variaciones del tipo de cambio que los países no pierden nada al renunciar a este instrumento? No necesariamente. Debemos comparar el proceso de ajuste obtenido aquí con el proceso de ajuste en una unión monetaria:

\imagenfit{Fig2_7.png}{Depreciación de la moneda y políticas deflacionistas comparadas: Tipo de Cambio vs. Unión Monetaria}{}

El panel izquierdo muestra cómo es el ajuste en una unión monetaria tras el shock negativo permanente de demanda en Francia.\footnote{Para que se restaure el nivel inicial de producción, los salarios nominales tienen que disminuir. Esto conduce a un desplazamiento hacia abajo de la curva de oferta. Obsérvese que, cuando se restaura la producción inicial —$Y_1$—, los salarios reales en Francia también disminuyen. Esto puede no ser evidente en en la figura porque observamos que tanto los salarios nominales franceses como el precio de los productos franceses disminuyen. Sin embargo, debe recordarse que el precio de los bienes alemanes importados no ha disminuido. Por tanto, los salarios reales en Francia disminuyen. Esto forma parte del proceso de cambios en los precios relativos.
} Por el contrario, el panel derecho muestra el proceso de ajuste cuando Francia conserva su moneda nacional y se ajusta al shock negativo de demanda permitiendo que su moneda se deprecie.\footnote{Como hemos visto, esto desplaza de nuevo hacia arriba la curva de demanda. Para restaurar permanentemente el nivel inicial de producción —$Y_1$—, debe resistirse la presión al alza sobre el salario nominal y, por tanto, sobre la curva de oferta. Esto solo puede lograrse si los trabajadores franceses están dispuestos a aceptar la disminución de su salario real implícita en situarse en el punto $F$.
} Así, observamos que la condición necesaria para restaurar el nivel inicial de producción es la misma en los dos regímenes monetarios: los trabajadores franceses deben aceptar una caída de sus salarios reales, lo que es difícil de lograr en ambos regímenes.\footnote{La cuestión que surge entonces es la siguiente: ¿en qué régimen monetario puede satisfacerse esta condición con mayor facilidad? La respuesta es que, en un mundo perfecto sin ilusión monetaria, no hay diferencia. Si los trabajadores franceses se resisten a una disminución de sus salarios reales, se resistirán tanto dentro como fuera de una unión monetaria. En ambos regímenes, el cambio en los precios relativos será difícil de lograr, y también lo será el retorno al nivel inicial de producción. Sin embargo, en un mundo menos perfecto, en el que los trabajadores tienen ilusión monetaria, estos pueden resistirse con más fuerza a las reducciones del salario real provocadas por una caída del salario nominal que a la misma reducción del salario real producida por aumentos de precios, manteniéndose constante el salario nominal. Así, en tal mundo puede decirse que será más difícil, y más costoso en términos de producción perdida, ajustarse a un shock de demanda dentro de una unión monetaria que fuera de ella.

Conclusión que se refuerza al considerar las implicaciones presupuestarias de los dos mecanismos de ajuste. Cuando, en una unión monetaria, un país debe reducir salarios y precios, la consiguiente caída de la producción también conduce a un aumento del déficit presupuestario y a un incremento de la ratio deuda pública/PIB. Como subrayamos anteriormente, los miembros de una unión monetaria no controlan el dinero en el que han emitido su deuda. Esto los hace vulnerables a ataques especulativos autocumplidos que pueden conducir a una crisis de liquidez, la cual a su vez puede degenerar en una crisis de solvencia. Así, durante el período de ajuste, estos países pueden enfrentarse a una crisis de deuda soberana, que conduce a un nuevo movimiento descendente de la actividad económica. Este movimiento descendente corre el riesgo de ser intenso y prolongado. Hubo varias devaluaciones en la UE antes del inicio de la unión monetaria que tuvieron éxito en restaurar el equilibrio interno y de la balanza comercial a un coste que probablemente fue menor que si no se hubiera utilizado el instrumento del tipo de cambio. Las devaluaciones francesas de 1982-83 —tras un período de importantes errores de política— destacan como casos exitosos (véase Sachs y Wyplosz, 1986). De forma similar, las devaluaciones belga y danesa de 1982 tuvieron bastante éxito en restablecer el equilibrio externo sin costes significativos en términos de desempleo (véase De Grauwe y Vanhaverbeke, 1990).
}

$$\text{\textbf{} Políticas monetarias nacionales para estabilizar shocks del ciclo económico}$$
Muchos shocks de demanda están relacionados con dinámicas de auge y caída que los hace endógenos y temporales. Estos shocks del ciclo económico pueden ser asimétricos en la medida en que los ciclos económicos de los distintos países que participan en la unión monetaria no están bien sincronizados o cuando las amplitudes de ciclos económicos similares son muy diferentes. La cuestión que surge aquí es la de la estabilización macroeconómica. 

¿Qué sucede cuando hay recesión en Francia y auge en Alemania? Si estos dos países forman una unión monetaria, tienen un problema. El banco central común queda paralizado: si reduce el tipo de interés para aliviar el problema francés, aumentará las presiones inflacionarias en Alemania; si eleva el tipo de interés para contrarrestar las presiones inflacionarias en Alemania, intensificará la recesión en Francia.\footnote{Además, dado que los shocks son temporales, no puede invocarse la flexibilidad salarial ni la movilidad del trabajo para resolver este problema. Por ejemplo, no tendría mucho sentido intentar inducir a los trabajadores franceses a abandonar Francia durante la recesión solo para atraerlos de vuelta durante el siguiente auge francés. De hecho, en una unión monetaria simplemente no hay solución a este problema: el banco central común no puede estabilizar la producción a nivel nacional; solo puede hacerlo a nivel de la unión.
} Sin embargo, cuando Francia y Alemania conservan sus propias monedas, disponen de los instrumentos para estabilizar la producción a nivel nacional. Así, cuando Francia sufre una recesión, el banco central francés puede estimular la demanda agregada reduciendo el tipo de interés y permitiendo que el franco francés se deprecie. De modo similar, cuando Alemania experimenta un auge, su banco central puede elevar el tipo de interés y permitir que la moneda se aprecie para moderar dicho auge.

Peor aún, en una unión monetaria, los gobiernos que no controlan la moneda en la que emiten su deuda se enfrentan a la desconfianza de los mercados financieros durante una recesión. Si esta desconfianza conduce a aumentos de los tipos de interés de los bonos públicos, los países que se enfrentan a una recesión pueden verse obligados a desactivar los estabilizadores automáticos de los presupuestos públicos y reforzando las fuerzas recesivas.

Por tanto, cabe concluir que en una unión monetaria existe poco margen para políticas de estabilización a nivel nacional. Las políticas monetarias no pueden utilizarse para estabilizar ciclos económicos impulsados nacionalmente. Además, es probable que estos movimientos cíclicos se vean exacerbados, ya que los presupuestos públicos pierden sus capacidades de estabilización automática. En esencia, los gobiernos nacionales pierden una de sus funciones esenciales: la estabilización del ciclo económico. Y, como resultado, es probable que pierdan legitimidad y, en consecuencia, que muchas personas situadas en el lado perdedor durante las recesiones económicas den la espalda a la unión monetaria y a la élite política nacional que la hizo posible. 

$$\text{\textbf{} Productividad e inflación en una unión monetaria: el efecto Balassa-Samuelson}$$
Hasta ahora hemos supuesto que las tasas nacionales de inflación se igualan en una unión monetaria. Esto no es necesariamente cierto. A menudo observamos que en las uniones monetarias existentes hay diferencias regionales en las tasas de inflación. Aunque estas diferencias son relativamente pequeñas, pueden ser significativas:

\imagenfit{Fig2_8.png}{Inflación media anual, Eurozona: 1999-2011 (\%)}{}

Puede verse que estas diferencias ciertamente no son nulas: en un extremo, Alemania e Irlanda tuvieron una tasa media anual de inflación inferior al 1,5\%, mientras que Italia, Grecia y España experimentaron una tasa anual de inflación de entre el 2,5\% y el 3\% durante 1999-2011. ¿Cómo pueden interpretarse estos fenómenos? La hipótesis de BALASSA y SAMUELSON puede ser una posible explicación.\footnote{Para consultar el desarrollo de la hipótesis [\authorfont{Ver Tema 3.B12}]
} Si existen diferenciales en el crecimiento de la productividad entre países de una unión monetaria, entonces las tasas de inflación también deben diferir: si la productividad crece más rápido en Italia que en Alemania, la inflación italiana tendrá que superar a la alemana dentro de una unión monetaria. Esta mayor inflación en Italia es entonces un fenómeno de equilibrio que refleja el mayor crecimiento de la productividad en ese país.

En definitiva, lo que se pretende es ensalzar la importancia de que estos diferenciales de inflación no deberían ser motivo de preocupación per se: precisamente porque la productividad crece más rápido en Italia que en Alemania, los salarios (y, por tanto, los precios de los bienes no comerciables) deben aumentar más rápido en Italia que en Alemania para mantener inalterada la posición competitiva de los sectores de bienes comerciables de ambos países. No obstante, no todas las diferencias observadas en las tasas de inflación dentro de una unión monetaria son resultado de este mecanismo equilibrador. De hecho, en algunos países como España, Irlanda, Grecia e Italia, los salarios aumentaron significativamente más rápido que la productividad. Como resultado, la mayor inflación registrada en estos países no formó parte del mecanismo equilibrador descrito en el modelo BALASSA-SAMUELSON, y condujo a deterioros significativos de sus posiciones competitivas.

$$\text{\textbf{}Políticas monetarias nacionales, inconsistencia temporal y credibilidad}$$
La idea de que, cuando el gobierno sigue determinadas políticas, juega un juego con el sector privado ha dominado la teoría macroeconómica desde la publicación de los artículos fundamentales de KYDLAND y PRESCOTT (1977) y BARRO y GORDON (1983). En particular, la reputación que adquieren los gobiernos al aplicar las políticas anunciadas tiene un gran impacto sobre cómo dichas políticas afectarán a la economía. Esta literatura también tiene implicaciones importantes para nuestra discusión sobre los costes de una unión monetaria. Conduce a una crítica fundamental de la visión según la cual las políticas monetarias nacionales son instrumentos de política que los gobiernos tienen a su disposición para utilizarlos de manera discrecional, en particular para estabilizar la economía ante shocks temporales. 

Supongamos que existen dos países que aún no forman una unión monetaria. Llamamos al primero Alemania y suponemos que su gobierno es duro. El segundo país se denomina Italia, donde el gobierno es blando. Utilizamos la condición de paridad del poder adquisitivo, que afirma que si la inflación en Italia excede la inflación en Alemania, la lira italiana tendrá que depreciarse frente al marco alemán para compensar este diferencial de inflación, es decir: $\dot e = p_I − p_G$.

\imagenfit{Fig2_14.png}{Equilibrio inflacionario en un modelo de dos países}{}

Como se observa, Italia tiene una tasa de inflación de equilibrio más alta que Alemania. Por tanto, su moneda tendrá que depreciarse continuamente. El problema de Italia es que podría alcanzar un equilibrio de inflación mucho menor que el punto $E$ si su gobierno pudiera convencer a sus ciudadanos de que, una vez en el punto $F$, no intentará alcanzar el punto $G$.

¿Puede Italia resolver su problema anunciando que se incorporará a una unión monetaria con Alemania? Para responder, supongamos primero que Italia anuncia que fijará su tipo de cambio con el marco alemán $\dot e = 0$. Dada la paridad del poder adquisitivo, esto fija la tasa de inflación italiana en el nivel alemán, la línea horizontal desde $C$. Italia parece ahora capaz de disfrutar de una menor tasa de inflación y las ganancias potenciales de bienestar son grandes, porque en el nuevo equilibrio la economía se sitúa sobre una curva de indiferencia más baja.

La cuestión, sin embargo, es si esta regla puede ser creíble. Una vez en el nuevo punto de equilibrio $F$, las autoridades italianas tienen un incentivo para generar una devaluación sorpresa de la lira. Esta devaluación sorpresa conduce a un aumento inesperado de la inflación y permite que la economía se desplace hacia el punto $G$. Con el tiempo, sin embargo, los agentes económicos ajustarán sus expectativas, de modo que la tasa de inflación de equilibrio terminará siendo la misma que antes de fijar el tipo de cambio. Así, fijar simplemente el tipo de cambio no resuelve el problema, porque la regla de tipo de cambio fijo no es más creíble que una regla de inflación fija.\footnote{Volveremos sobre esta cuestión cuando analicemos el funcionamiento de los sistemas de tipo de cambio fijo. Algunos economistas han sostenido que fijar el tipo de cambio puede ser una regla que posee inherentemente más credibilidad que anunciar una regla de inflación constante.
}

Existen, sin embargo, otros arreglos que potencialmente pueden resolver el problema de alta inflación en Italia. Imaginemos que Italia decide abolir su moneda y adoptar la moneda de Alemania. Si tal arreglo pudiera hacerse creíble, es decir, si los ciudadanos italianos estuvieran convencidos de que, una vez tomada esta decisión y convertido el marco en la moneda nacional, las autoridades italianas nunca rescindirían esta decisión, entonces Italia podría alcanzar el mismo equilibrio de inflación que Alemania y la línea horizontal que conecta el equilibrio de inflación alemán con Italia definiría un equilibrio creíble. Dado que Italia ya no tiene política monetaria independiente, sus autoridades monetarias, con preferencias \textit{blandas}, han dejado de existir y, por tanto, no pueden devaluar la lira. En palabras de Giavazzi y Pagano (1988), Italia ha tomado prestada credibilidad de Alemania, porque su gobierno tiene las manos monetarias firmemente atadas.

Este es ciertamente un resultado muy fuerte. Lleva a la conclusión de que existe una gran ganancia potencial para Italia al incorporarse a una unión monetaria con Alemania. Además, no hay pérdida de bienestar para Alemania. Así, una unión monetaria solo genera ganancias. Este análisis fue muy influyente antes del inicio de la Eurozona, especialmente en países de alta inflación, que percibieron la entrada en la Eurozona como un almuerzo gratis que les permitía introducir estabilidad macroeconómica a coste cero.

Existen, sin embargo, dos consideraciones que tienden a suavizar esta conclusión:
\begin{lnum}
    \item En primer lugar, debe quedar claro del análisis anterior que solo una unión monetaria plena establece la credibilidad requerida para Italia. Es decir, Italia debe estar dispuesta a eliminar su moneda nacional, de manera muy similar a como Alemania Oriental hizo el 1 de julio de 1990, cuando adoptó el marco alemán occidental. Cualquier cosa inferior a una unión monetaria plena afrontará un problema de credibilidad.\footnote{Algunos países han experimentado con cajas de conversión (por ejemplo, Argentina y Hong Kong). Estos son regímenes monetarios más restrictivos que los tipos de cambio fijos, pero no llegan a constituir una unión monetaria plena. Como resultado, estos países han estado ocasionalmente sometidos a ataques especulativos intensos. El régimen de caja de conversión argentino no sobrevivió al ataque especulativo de 2001 y fue abandonado en enero de 2000.
    } Como se señaló, cuando Italia fija su tipo de cambio respecto al marco y conserva su propia moneda, la credibilidad de este arreglo de tipo de cambio fijo estará en duda;
    \item En segundo lugar, y más importante para nuestro propósito actual, hemos supuesto que el banco central de la unión monetaria es el banco central alemán. Bajo este arreglo, Italia se beneficia de la reputación del banco central alemán en la consecución de baja inflación. Supongamos, sin embargo, que el nuevo banco central es una institución nueva, en la que tanto las autoridades alemanas como las italianas están representadas por igual. ¿Tendría ese nuevo banco central la misma reputación que el antiguo banco central alemán? Esto está lejos de estar claro. Si se percibe que el banco central de la unión es menos duro que el banco central alemán antes de la creación de la unión, entonces el nuevo equilibrio de inflación de la unión será superior al que prevalecía en Alemania antes de la unión. Italia aún podría ganar con tal arreglo. Alemania, sin embargo, perdería y no estaría muy entusiasmada por formar tal unión.
\end{lnum}

El modelo de Barro-Gordon para economías abiertas también nos permite entender por qué algunos países deciden dolarizar sus economías. La dolarización es un régimen monetario en el que un país decide adoptar unilateralmente el dólar para los pagos internos. Ecuador, El Salvador y Panamá son países que han pasado a una economía dolarizada. Invariablemente, estos han sido países que experimentaban alta inflación y enfrentaban grandes dificultades para reducirla debido a la falta de credibilidad de sus autoridades monetarias nacionales. La dolarización es entonces un instrumento para resolver el problema de credibilidad mediante la abolición de la moneda nacional y la adopción del dólar, respaldado por la credibilidad del Sistema de la Reserva Federal de Estados Unidos. También existen países de Europa Central y Oriental —Kosovo y Montenegro— que han decidido adoptar unilateralmente el euro como moneda nacional por razones muy similares a las que llevaron a pequeños países latinoamericanos a adoptar el dólar.

Concluimos del análisis precedente que los problemas de credibilidad son importantes al evaluar los costes de una unión monetaria. Primero, la opción de permitir una depreciación de la moneda es un arma de doble filo para las autoridades nacionales. El conocimiento de que puede usarse en el futuro complica considerablemente las políticas macroeconómicas. Segundo, la literatura sobre inconsistencia temporal también nos enseña lecciones importantes sobre los costes de una unión monetaria: las depreciaciones cambiarias no pueden usarse para corregir toda perturbación que ocurra en una economía. Una depreciación no es, como en el análisis de Mundell, un instrumento flexible que pueda utilizarse frecuentemente. Cuando se usa una vez, afecta a su uso futuro, porque genera fuertes efectos sobre las expectativas. Es un instrumento peligroso que puede perjudicar a quienes lo utilizan. Cada vez que los responsables de política usen este instrumento, tendrán que evaluar las ventajas obtenidas hoy frente al coste de que será más difícil utilizarlo eficazmente en el futuro.

Esto ha llevado a algunos economistas a concluir que el instrumento del tipo de cambio no debería utilizarse en absoluto y que los países incluso ganarían renunciando irrevocablemente a su uso. Esta conclusión va demasiado lejos. Hubo muchos casos observados en Europa durante las décadas de 1980 y 1990 en los que las devaluaciones monetarias se utilizaron con gran éxito. Los ingredientes de este éxito han sido típicamente que la devaluación estuvo acompañada de otros cambios drásticos de política (a veces con un cambio de gobierno, por ejemplo Bélgica y Dinamarca en 1982). Como resultado, la devaluación fue percibida como un cambio extraordinario de política que no podía repetirse fácilmente en el futuro. Bajo esas condiciones, los efectos negativos sobre la reputación pudieron mantenerse bajo control. Algunos países, en particular Dinamarca en 1982, incluso parecen haber mejorado rápidamente su reputación tras la devaluación. Renunciar a la posibilidad de usar este instrumento por tiempo indefinido implica un coste para una nación.

Esta conclusión se refuerza al tener en cuenta que, cuando los países renuncian a sus monedas nacionales, la deuda pública tendrá que emitirse en una moneda que no está bajo el control del gobierno. Como se argumentó anteriormente, esto hace a estos gobiernos más vulnerables a verse envueltos en una crisis de deuda. Así, parece existir aquí un trade-off. Los países de alta inflación pueden mejorar su reputación antiinflacionaria entrando en una unión monetaria. Sin embargo, al mismo tiempo, si su reputación en el ámbito presupuestario no es fuerte, su entrada en la unión los vuelve más vulnerables a una crisis de deuda soberana. La entrada en una unión monetaria no crea un “almuerzo gratis” —véase De Grauwe y Ji (2018), donde esta idea se desarrolla más y se muestra que los países que experimentaron crisis cambiarias a comienzos de los años noventa fueron, en términos generales, los mismos que experimentaron crisis de deuda soberana en 2010-12—.

$$\text{\textbf{} Mundell una vez más}$$
La contribución de MUNDELL a la teoría de las áreas monetarias óptimas es sustancial. Esta es también la principal razón por la que recibió el Premio Nobel en 1999. En esa ocasión fue aclamado como uno de los padres de la UEM. Esto puede parecer bastante sorprendente, dado que su artículo fundamental, que sentó las bases del análisis posterior de los costes y beneficios de una unión monetaria, era muy escéptico sobre las posibilidades de una unificación monetaria exitosa en Europa. Existe, sin embargo, otra contribución de MUNDELL, escrita en 1973, que es mucho más optimista sobre los beneficios de una unión monetaria.

La reconsideración de Mundell se basa en dos argumentos. Primero, sostuvo que una unión monetaria es una forma más eficiente de organizar un sistema de seguro para afrontar shocks asimétricos que un sistema de monedas nacionales con incertidumbre cambiaria. Para entenderlo, considérese nuevamente el shock asimétrico que afecta a Francia y Alemania. Supóngase, sin embargo, que el shock asimétrico es temporal, por ejemplo, resultado de una desincronización de los ciclos económicos o de factores estacionales. Cuando los dos países forman una unión monetaria, existen flujos automáticos de capital que suavizarán el golpe para los residentes franceses. Los consumidores franceses que saben que la pérdida de renta es temporal y desean mantener sus niveles de consumo encontrarán relativamente fácil endeudarse con consumidores alemanes, que experimentan un aumento temporal de la renta, al menos si los mercados de capitales están suficientemente integrados. Es probable que este sea el caso en una unión monetaria. En ausencia de una unión monetaria, tal mecanismo será más difícil, porque la existencia de monedas separadas y de un tipo de cambio que puede variar hará que los consumidores alemanes sean reacios a prestar a los consumidores franceses por el riesgo implicado. Así, en ausencia de una unión monetaria, los shocks asimétricos temporales no pueden asegurarse fácilmente.

Una segunda razón por la que Mundell reconsideró su pesimismo sobre la unión monetaria se relaciona con lo siguiente. En un mundo incierto, es probable que los movimientos del tipo de cambio sean una fuente de shocks asimétricos, en lugar de ser un mecanismo que permite a los países ajustarse mejor a shocks asimétricos. Existe abundante evidencia de que los movimientos de los tipos de cambio a menudo están desconectados de fundamentos subyacentes como los diferenciales de inflación o el crecimiento de la producción. Los tipos de cambio suelen estar impulsados por factores psicológicos y comportamiento gregario, generando gran volatilidad que, a su vez, se transmite a las economías nacionales (De Grauwe y Grimaldi, 2006). Sobre la base de estos argumentos, el nuevo MUNDELL pasó a considerar la unión monetaria como una forma de reducir los shocks asimétricos y mejorar el seguro frente a ellos.

Hay ciertamente un importante grano de verdad en este análisis. No obstante, debe tenerse presente que es fácil exagerar este optimismo. Primero, la conclusión de Mundell sobre el seguro frente a shocks asimétricos se sostiene cuando los shocks son temporales. Cuando son permanentes, es improbable que los consumidores alemanes estén dispuestos a conceder crédito indefinidamente a los consumidores franceses. El ajuste debe realizarse mediante cambios en precios y salarios. Volvemos al análisis de Mundell de 1961.

Segundo, incluso si los shocks son temporales, cuando los mercados financieros comienzan a desconfiar de uno o más países miembros de una unión monetaria, las propiedades estabilizadoras y aseguradoras de los mercados financieros desaparecen. El país que experimenta un shock negativo no puede contar con flujos financieros procedentes de otros países; por el contrario, en un régimen de desconfianza los inversores probablemente venderán los bonos públicos del país afectado por el shock negativo, provocando un aumento del tipo de interés interno y una salida de liquidez. Los mercados financieros dejan de ser una fuerza estabilizadora y no logran facilitar mecanismos de seguro frente a shocks asimétricos.

Tercero, aunque la volatilidad cambiaria puede ser una fuente independiente de shocks asimétricos, sigue siendo cierto que pueden ocurrir grandes shocks que pueden afrontarse más fácilmente cuando se permite que el tipo de cambio se ajuste.

$$\text{\textbf{} El coste de la unión monetaria y la apertura de los países}$$

Hemos desarrollado varias ideas relacionadas con la cuestión de cómo la apertura de un país afecta al coste de la unión monetaria. Aquí nos concentramos en dos de ellas, que, como se verá, tienen efectos opuestos. En primer lugar, está la relación entre el grado de apertura y la aparición de shocks asimétricos. Hemos presentado dos visiones: la visión de la Comisión Europea, que sostiene que una mayor integración conduce a una mayor simetría de los shocks, y la visión de KRUGMAN, que considera negativa esta relación entre integración y simetría. Según la primera visión, podemos concluir que una mayor apertura reduce el coste de una unión monetaria, ya que reduce la probabilidad de que se produzcan shocks asimétricos. Según la segunda visión, sin embargo, esta conclusión se invierte: los costes de una unión monetaria aumentan con el grado de apertura de los países.

La segunda idea relevante para nuestro análisis sobre cómo la apertura afecta al coste de la unión monetaria tiene que ver con la eficacia del tipo de cambio para hacer frente a shocks asimétricos. Consideremos dos países: uno relativamente abierto y otro relativamente cerrado

\imagenfit{Fg2_14.png}{Efectividad de una devaluación en función del grado de apertura}{}

Tanto los efectos de demanda como los de oferta de una depreciación difieren entre los dos países. En lo que respecta a los efectos por el lado de la demanda, la misma depreciación tiene un efecto más fuerte en la economía relativamente abierta que en la relativamente cerrada. Para entenderlo, considérese dos casos extremos. Supóngase que la economía relativamente abierta exporta el 99 por ciento de su PIB, mientras que la relativamente cerrada solo exporta el 1 por ciento de su PIB. La misma depreciación, por ejemplo del 10 por ciento, necesariamente elevará más la demanda agregada en el primer país que en el segundo. En el gráfico mostramos esta diferencia mediante el hecho de que la curva de demanda se desplaza más hacia afuera en el país relativamente abierto que en el país relativamente cerrado.

Los dos países también difieren respecto a los efectos por el lado de la oferta de la depreciación. Cabe esperar que, en la economía relativamente abierta, el desplazamiento hacia arriba de la curva de oferta tras la depreciación sea más pronunciado que en la economía relativamente cerrada. Esto se debe a que la economía más abierta importa más —como porcentaje del consumo total—, de modo que el índice de precios al consumo aumenta más, generando una espiral salarios-precios más intensa que en la economía relativamente cerrada.

Llegamos ahora a la siguiente conclusión. Los efectos combinados de demanda y oferta de una depreciación en los dos países son tales que no podemos decir a priori en qué país es más eficaz la depreciación para estimular la producción. Lo que sí podemos concluir, sin embargo, es que la depreciación se sentirá con mayor intensidad sobre el nivel agregado de precios en la economía más abierta que en la relativamente cerrada. Esto significa que el uso sistemático del instrumento del tipo de cambio conducirá a una mayor variabilidad de precios en la economía más abierta que en la relativamente cerrada.

Podemos generalizar este punto del siguiente modo. El uso sistemático de la política monetaria para estabilizar la producción y el empleo probablemente conducirá a una mayor variabilidad de precios en la economía más abierta que en la relativamente cerrada. La razón es que estas políticas monetarias provocarán movimientos del tipo de cambio, es decir, una política monetaria expansiva —o contractiva— conducirá a una depreciación —o apreciación— de la moneda. Estos movimientos cambiarios tendrán efectos más pronunciados sobre los precios internos en la economía relativamente abierta. En la medida en que la variabilidad de precios implique costes, el uso sistemático de políticas monetarias nacionales será más costoso en la economía más abierta. Este punto fue reconocido por primera vez por McKinnon (1963) en su importante contribución a la teoría de las áreas monetarias óptimas.

De esta discusión sobre la eficacia de una depreciación en economías abiertas y relativamente cerradas puede concluirse que, para un mismo efecto sobre la producción, la depreciación probablemente será más costosa en la economía abierta debido a la mayor variabilidad de precios que implica. Así, la pérdida de la capacidad de utilizar políticas monetarias nacionales probablemente será menos costosa para la economía relativamente abierta que para la relativamente cerrada.

Combinando el análisis de la eficacia de una depreciación con el análisis de la relación entre apertura y shocks asimétricos, puede derivarse la conclusión de que el coste de una unión monetaria muy probablemente disminuye con el grado de apertura de un país. Mostramos esto en la Figura 2.16, tomada de Krugman (1990). En el eje vertical se representa el coste de una unión monetaria —es decir, el coste de renunciar al instrumento del tipo de cambio—. Este coste se expresa como porcentaje del PIB. En el eje horizontal se representa la apertura del país en relación con los países con los que desea formar una unión monetaria. Esta apertura se representa mediante la participación del comercio en el PIB del país considerado. Observamos que, a medida que aumenta la apertura, disminuye el coste de incorporarse a una unión monetaria.

Esta conclusión, sin embargo, no se sostiene de forma general. No podemos excluir la posibilidad de que la relación tenga pendiente positiva. Esta situación puede darse cuando la probabilidad de shocks asimétricos aumenta significativamente con el grado de apertura del país —el escenario de Krugman—. La presunción, sin embargo, es que es improbable que este caso ocurra. Primero, es más probable que, con la integración comercial, los shocks asimétricos se vuelvan menos probables —véase nuestro análisis de la visión de la Comisión Europea—. Segundo, incluso si la integración comercial conduce a más shocks asimétricos, el coste de utilizar el instrumento del tipo de cambio probablemente compensará ese efecto en economías más altamente integradas. Por tanto, mantendremos nuestra presunción de que el coste de la unión monetaria disminuye con la apertura.

$$\text{Conclusión}$$

La crítica de la teoría tradicional de las áreas monetarias óptimas, desarrollada por Mundell y McKinnon, nos ha permitido añadir matices importantes a esta teoría. En particular, ha cambiado nuestra visión sobre los costes de una unión monetaria. La teoría tradicional de las áreas monetarias óptimas tiende a ser bastante pesimista sobre la posibilidad de que los países se incorporen a una unión monetaria a bajo coste. La crítica discutida en este capítulo es menos pesimista: los costes de formar una unión monetaria parecen menos prohibitivos. Las principales razones de esta conclusión son dos.

Primero, la capacidad de las variaciones del tipo de cambio para absorber shocks asimétricos es más débil de lo que la teoría tradicional de las áreas monetarias óptimas, de inspiración keynesiana, nos llevó a creer. Las variaciones del tipo de cambio normalmente no tienen efectos permanentes sobre la producción y el empleo.

Segundo, los países que mantienen políticas monetarias y cambiarias independientes a menudo descubren que los movimientos del tipo de cambio se convierten en una fuente de perturbaciones macroeconómicas, en lugar de ser instrumentos de estabilización macroeconómica. Paradójicamente, la volatilidad de los tipos de cambio puede ser una fuente significativa de shocks asimétricos para los países que conservan sus propias monedas. En contra de la visión antigua, los cambios del tipo de cambio no son instrumentos que los responsables de política puedan utilizar de manera flexible y sin coste.

¿Significa esto que las ideas de la teoría tradicional de las áreas monetarias óptimas han dejado de ser relevantes? La respuesta es no. A pesar de nuestra crítica, el núcleo duro del análisis de las áreas monetarias óptimas sigue en pie. Puede formularse del siguiente modo.

Existen diferencias importantes entre países que no van a desaparecer en una unión monetaria. Muchas, si no la mayoría, de estas diferencias tienen un origen político e institucional. Los Estados-nación que actualmente conforman la unión monetaria en Europa han mantenido muchas de sus peculiaridades nacionales. Los mercados laborales siguen teniendo rasgos institucionales nacionales específicos. Los sistemas legales no son idénticos, generando diferencias en el funcionamiento de los mercados financieros, de los mercados de vivienda y posiblemente de otros mercados. Los gobiernos de los Estados miembros de la unión monetaria utilizan sistemas fiscales distintos y aplican políticas de gasto diferentes. La UEM por sí misma elimina algunas de estas diferencias, pero ciertamente no todas. Las que permanecen conducen a movimientos divergentes en la producción y los precios nacionales, creando la necesidad de ajustes nacionales difíciles. La ausencia de políticas monetarias y cambiarias nacionales que ayuden en este proceso de ajuste se percibe entonces como un coste de la unión monetaria.

Además, y de forma más importante, parece que en una unión monetaria los gobiernos nacionales no solo pierden su capacidad de estabilizar el ciclo económico mediante políticas monetarias, sino que también es probable que pierdan los estabilizadores automáticos de los presupuestos públicos cuando, durante una recesión, los mercados financieros entran en pánico sobre la capacidad de estos gobiernos para seguir atendiendo su deuda. Esto puede conducir a crisis de liquidez autocumplidas que obliguen a los países a aplicar una austeridad severa en medio de una recesión. Volveremos sobre esta característica importante de las uniones monetarias que carecen de unión fiscal en el Capítulo 5.

Así, cuando los países de la UE se incorporaron a la UEM el 1 de enero de 1999, asumieron un riesgo. La naturaleza de este riesgo se ha hecho visible desde el inicio de la crisis de deuda soberana en la Eurozona. Todo esto debería servir como advertencia para aquellos países que contemplan incorporarse a la UEM, o para países del resto del mundo que avanzan hacia sus propias uniones monetarias.

El riesgo de altos costes de ajuste frente a perturbaciones asimétricas puede reducirse aplicando dos estrategias. Una consiste en hacer los mercados más flexibles, de modo que los shocks asimétricos puedan ajustarse mejor. La otra consiste en acelerar el proceso de unificación política. Esto reducirá las idiosincrasias nacionales y, por tanto, la aparición de perturbaciones asimétricas de origen político o institucional. Además, facilitará la estabilización del ciclo económico al nivel de la unión monetaria.


\clearpage
\anexo{Amplicación: Crítica a los argumentos sobre los beneficios de un Area Monetaria}\label{anx:critica_beneficios}


$$\textbf{Incertidumbre}: \text{Bienestar y beneficios de la empresa}$$
En el cuerpo del Tema hemos visto como los beneficios en materia de bienestar tienen también sus debilidades. En particular, hemos visto como estas anancias puedes ser menores cuando evalúamos los beneficios esperados en el mundo de la empresa.

Existe un aspecto de la incertidumbre que es bastante serio y que puede socavar la relevancia del análisis anterior. Se trata de que las variaciones del tipo de cambio no se distribuyen normalmente. La naturaleza de los movimientos del tipo de cambio es que existen períodos de relativa tranquilidad seguidos de períodos de gran turbulencia (clustering). Durante estos períodos turbulentos, los cambios del tipo de cambio pueden ser muy grandes y sostenidos en una dirección, produciendo burbujas seguidas posteriormente de colapsos:

\imagenfit{Fig3_4.png}{Burbujas y colapsos en el mercado de intercambio de divisas (FOREX)}{DE GRAUWE y GRIMALDI (2006)}

Tales movimientos crean riesgos de cola, es decir, el riesgo de un cambio muy grande que ocurre con baja probabilidad y que, cuando ocurre, puede tener un efecto devastador: la caída del tipo de cambio puede ser tan grande que el precio caiga muy por debajo de la curva de coste marginal (y de de coste medio), obligando a la empresa a cerrar. En este contexto, el coste de la quiebra y de reasignación de factores de producción empleados serán típicamente sustancial, creando grandes costes de los movimientos del tipo de cambio.

Cabe tener en cuenta que estos movimientos cambiarios son un problema recurrente con tipos de cambio libremente flotantes. Generan grandes costes de ajuste y mucha miseria económica. También ocurrieron masivamente dentro de la Unión Europea a comienzos de la década de 1990, cuando algunas monedas (por ejemplo, la lira italiana y la peseta española) se depreciaron entre un 20 y un 30 por ciento en unas pocas semanas, creando grandes costes de ajuste en países como Alemania y el Benelux. Estos movimientos del tipo de cambio entre monedas de países altamente integrados fueron un factor importante para convencer a muchos líderes de estos países de avanzar hacia una unión monetaria, principalmente por dos razones: (i) se percibía cada vez más que sería difícil, si no imposible, gestionar los tipos de cambio de forma ordenada en un mundo de libre movilidad de capitales; y, (ii) estos movimientos cambiarios se veían cada vez más como fuentes importantes de perturbaciones asimétricas (argumento de MUNDELL), en lugar de como variables que pudieran utilizarse para ajustarse a shocks asimétricos (argumento de FRIEDMAN).



\clearpage
\anexo{¿Cómo \textit{salvar} las imperfecciones de una Unión Monetaria?: Mecanismos de seguro dentro de la Unión Monetaria}\label{anx:mecanismos_seguros}
El shock asimétrico analizado en los párrafos anteriores es un acontecimiento exógeno con efectos permanentes, producido por un cambio en las preferencias de los consumidores. Sin embargo, muchos shocks asimétricos son de naturaleza distinta.

El capitalismo es una invención humana extraordinaria que logra canalizar la iniciativa individual y la creatividad hacia la acumulación de capital y un progreso material cada vez mayor. Sin embargo, también es inherentemente inestable. Períodos de optimismo y pesimismo se alternan, creando fases de auge y recesión en la actividad económica. Los auges son maravillosos; las recesiones generan grandes dificultades para muchas personas.

Los auges y caídas son endémicos en el capitalismo porque muchas decisiones económicas están orientadas hacia el futuro. Inversores y consumidores miran al futuro para decidir si invertir o consumir. Pero el futuro es incierto. Nadie lo conoce. Como resultado, al hacer previsiones, consumidores e inversores se observan entre sí. Esto hace posible que el optimismo de un individuo se transmita a otros, generando un movimiento de optimismo autocumplido. El optimismo induce a los consumidores a consumir más y a los inversores a invertir más, validando así dicho optimismo. Pero lo contrario también es cierto. Cuando aparece el pesimismo, el mismo mecanismo de comportamiento gregario conduce a una caída autocumplida de la actividad económica. Predominan los \textit{animal spirits} (KEYNES 1936; AKERLOF y SHILLER 2009; DE GRAUWE 2012).

Mientras estos movimientos en los \textit{animal spirits} estén sincronizados entre los Estados miembros de la unión monetaria, no generan problemas adicionales para la unión; es decir, el hecho de pertenecer a una unión monetaria no agrava los auges y caídas. La situación es distinta si estos movimientos no están sincronizados, es decir, cuando algunos países experimentan auges y otros recesiones. En dichos escenarios podemos experimentar dos posibles resultados: (i) un resultado benigno, en el que la unión puede convivir con el ciclo desincronizado; y, (ii) una escenario \textit{maligno}, en el que la asincronía del ciclo tensa la Unión Monetaria hasta tal punto que puede forzar la desintegración por las razones mencionadas. ¿Por qué? ¿Cuándo experimentaríamos uno u otro?

Examinemos primero un escenario benigno basado en la confianza (por ejemplo, un shock como el respresntado en la primera figura).

Dado un shock cíclico, que se presupone es temporal, se espera que, después de algún tiempo, el proceso observado sea el contrario: Francia experimentando un auge y Alemania una recesión. En este escenario, no es necesario que Francia intente ajustarse mediante reducciones salariales y de precios, ni que Alemania lo haga mediante aumentos salariales y de precios, ni tampoco mediante la emigración de trabajadores franceses hacia Alemania. Por tanto, el ajuste se estructura sobre los estabilizadores automáticos del presupuesto, que cumplen su función de estabilizar el ciclo económico: en Francia, la recesión conduce a un déficit presupuestario; en Alemania, el auge conduce a un superávit presupuestario. Este mecanismo tiende a reducir la intensidad de la recesión en Francia, porque al incurrir en un déficit el gobierno francés inyecta poder adquisitivo en la economía. También reduce la intensidad del auge en Alemania, porque los superávits presupuestarios reducen el poder adquisitivo en ese país.

Si los inversores mantienen su confianza en la capacidad del gobierno francés para hacer frente a su deuda (que inevitablemente aumenta durante una recesión), están dispuestos a comprar la deuda pública adicional sin exigir un mayor tipo de interés. Por tanto, el tipo de interés francés puede mantenerse constante. ¿Por qué? La razón es que en Alemania el gobierno tiene un superávit presupuestario y, en consecuencia, retira bonos del mercado (la oferta de bonos alemanes disminuye), mientras que en Francia, la oferta de bonos aumenta. Si los mercados confían en el gobierno francés tanto como en el alemán, estarán dispuestos a compensar la menor tenencia de bonos alemanes con una mayor tenencia de bonos franceses, considerándolos sustitutivos perfectos. De ello se sigue que el gobierno francés puede financiar fácilmente su déficit, porque los inversores (principalmente alemanes en este caso) están dispuestos a comprar esos bonos.

Así, en este escenario benigno, observamos que los mercados de capital dentro de la unión monetaria desempeñan un papel estabilizador: cuando Francia atraviesa dificultades debido a una recesión, los mercados permiten transferir recursos desde el país en auge hacia el país en recesión, aliviando así el coste de la recesión.

Consideremos ahora el otro escenario, en el que el aumento del déficit y de la deuda en Francia lleva a los inversores a perder la confianza en el gobierno francés, que puede ocurrir si la recesión es especialmente profunda y existe gran incertidumbre sobre su duración (algo no muy lejano a lo que se puede preveer sin cambios en el rumbo actual). En este caso, los inversores comenzarán a vender bonos del gobierno francés y a comprar bonos alemanes. Esto genera un flujo de liquidez desde Francia hacia Alemania (lo contrario que en el escenario anterior) y un aumento del tipo de interés a largo plazo en Francia, acompañado de una disminución en Alemania. La curva de demanda agregada en Francia se desplazará con más intensidad hacia el origen, intensificando y prolongando la recesión; mientras que en Alemania ocurrirá lo contrario. Obsérvese nuevamente la naturaleza autocumplida de las expectativas: si los inversores anticipan problemas en Francia debido a déficits y deuda elevados, sus acciones contribuyen a que dichos problemas se materialicen.

Así, en este escenario de desconfianza, los movimientos del ciclo económico se amplifican: la recesión es más profunda en Francia y el auge más intenso en Alemania. La pertenencia a una unión monetaria conduce entonces a una mayor volatilidad de la producción y el empleo, lo cual no es una característica deseable.

Si Francia y Alemania no hubieran formado una unión monetaria, podrían haber mitigado estas dinámicas desestabilizadoras. Consideremos el caso de Francia suponiendo que ha mantenido su independencia monetaria y su propia moneda. Cuando, durante una recesión, los inversores venden bonos franceses y compran bonos alemanes, deben pasar por el mercado de divisas: venden francos y compran marcos. Como resultado, el franco se deprecia y el marco se aprecia. La depreciación del franco tiende a estimular la demanda agregada en Francia, mientras que la apreciación del marco tiende a reducir la demanda agregada en Alemania. Existe, por tanto, un efecto estabilizador derivado de los movimientos del tipo de cambio, que está ausente cuando ambos países pertenecen a una unión monetaria.

Por consiguiente, en una unión monetaria, los movimientos del ciclo económico se amplifican si los mercados financieros no confían plenamente en la solvencia de uno o más gobiernos miembros. En el Recuadro 1.2 se presenta un estudio de caso de la Eurozona durante la reciente “Gran Recesión”, mostrando cómo los shocks asimétricos fueron amplificados por grandes divergencias en los tipos de interés a largo plazo, incrementando la necesidad de ajustes dolorosos en salarios y movilidad laboral.

Resulta útil contar con un mecanismo de seguro que permita transferencias de renta hacia el país que experimenta un shock negativo de demanda. Sin embargo, este mecanismo no sustituye al ajuste cuando el shock es permanente; lo que hace es otorgar más tiempo para realizar dicho ajuste. En la medida en que los países enfrentan rigideces y disponen de sistemas de seguro poco desarrollados, los costes de la unión monetaria pueden ser elevados.

Cuando los shocks asimétricos son temporales, es decir, resultado de auges y caídas desincronizados, la cuestión relevante no es tanto la flexibilidad como la estabilidad. El hecho de que los países miembros de una unión monetaria sean vulnerables a cambios en el sentimiento de los mercados puede generar mayor volatilidad cíclica. Así, un país en recesión y con aumento del déficit puede verse afectado por ventas masivas de su deuda pública, lo que conduce a una crisis de liquidez y a un aumento de los tipos de interés. Esto probablemente obligará al gobierno a aplicar políticas de austeridad (aumentar impuestos y reducir el gasto), exacerbando la recesión. Los gobiernos descubren entonces que su capacidad para estabilizar la economía está severamente limitada; peor aún, se ven obligados a aplicar políticas fiscales que desestabilizan la economía. 

$$\text{\textbf{Mecanismos de seguro (I):} Unión Presupuestaria}$$
Cuando los Estados pertenecientes a una Unión Monetaria se ven afectados por grandes shocks asimétricos enfrentan difíciles problemas de ajuste. Dado que los shocks de demanda asimétricos tienden a generar aumentos en los déficits presupuestarios en algunos países, los mercados financieros pueden forzar crisis de liquidez en estos países, amplificando así los shocks. ¿Es posible diseñar un mecanismo que alivie estos problemas y reduzca los costes de una unión monetaria?

En principio, es posible diseñar tal mecanismo en dos partes. La primera se refiere al papel del banco central común en hacer posible evitar crisis de liquidez. La segunda consiste en centralizar una parte significativa de los presupuestos nacionales en un presupuesto común de la unión. Aquí nos concentramos en esta segunda parte. 

La centralización de los presupuestos nacionales equivale a tener una unión monetaria junto con una unión presupuestaria. Tal unión presupuestaria cumple dos funciones. En primer lugar, crea un mecanismo de seguro que desencadena transferencias de renta desde el país que experimenta buenos tiempos hacia los países afectados por shocks adversos. Al hacerlo, reduce el coste de estos shocks negativos. En segundo lugar, una unión presupuestaria permite consolidar una parte significativa de la deuda pública nacional en una deuda común, protegiendo así a sus miembros frente a crisis de liquidez y suspensiones de pagos forzadas. Analicemos estos dos mecanismos.

⸻

Una unión presupuestaria como mecanismo de seguro

Volvamos al modelo de dos países (Francia y Alemania) y supongamos que una gran parte de sus presupuestos públicos se centraliza a nivel europeo. Supongamos, por tanto, que existe un gobierno europeo que recauda directamente impuestos (incluidas cotizaciones sociales) y transfiere directamente ingresos (por ejemplo, pensiones y prestaciones por desempleo) a los residentes en Francia y Alemania.

Como resultado de esta centralización presupuestaria, una caída de la producción en Francia conduce a una reducción de los ingresos fiscales del gobierno europeo procedentes de Francia, mientras que los ingresos fiscales procedentes de Alemania aumentan debido al crecimiento de su producción. Al mismo tiempo, el gobierno europeo incrementa su gasto (prestaciones por desempleo) en Francia y lo reduce en Alemania. El resultado neto es que el presupuesto central redistribuye automáticamente renta desde Alemania —donde la producción ha aumentado— hacia Francia —donde ha disminuido—.

Dicho de otro modo, esta centralización permite a los ciudadanos franceses suavizar su consumo tras un shock negativo de producción. Obsérvese que también existe suavización del consumo en Alemania, aunque en sentido inverso. Como resultado, el coste de la unión monetaria se reduce, ya que los ciudadanos pueden estabilizar su consumo a lo largo del tiempo pese a los shocks asimétricos. La razón por la que Alemania podría estar interesada en este mecanismo es que también se beneficiaría cuando sufra un shock negativo.

Como en muchos sistemas de seguro, el principal problema es el riesgo moral. Esto se observa en su funcionamiento dentro de los países. En muchos Estados (por ejemplo, Bélgica, Alemania e Italia), el presupuesto nacional transfiere automáticamente renta desde regiones con alto crecimiento hacia regiones con bajo crecimiento. Estas transferencias tienden a reducir la presión para que las regiones se ajusten, y como resultado pueden volverse permanentes. La aplicación de estos esquemas a nivel europeo sería problemática, ya que podría generar transferencias grandes y permanentes entre países, provocando resistencia en aquellos que financian dichas transferencias (véase Schelkle, 2017).

⸻

Una unión presupuestaria como mecanismo de protección

Hemos visto que una unión monetaria en la que cada país mantiene su independencia presupuestaria es muy frágil. En tal unión, los gobiernos emiten deuda en una moneda que no controlan, lo que los hace vulnerables a episodios de desconfianza que pueden provocar crisis de liquidez y suspensiones de pagos. Es evidente que, en principio, una unión presupuestaria puede resolver este problema.

La razón es que, en una unión presupuestaria, la deuda pública nacional (o al menos una parte significativa de ella) se centraliza en una deuda común de la unión. Esto tiene dos efectos. En primer lugar, la centralización de los mercados de deuda pública en un mercado común elimina los movimientos desestabilizadores de capital entre mercados nacionales. En segundo lugar, el gobierno de la unión adquiere el carácter de un gobierno “independiente”, es decir, emite deuda en una moneda sobre la cual tiene pleno control. Por tanto, no puede ser objeto de una crisis de liquidez (al menos si la unión mantiene un tipo de cambio flexible frente al resto del mundo, como en el caso del Reino Unido).

Esta unión presupuestaria implica además la existencia de un gobierno de la unión suficientemente fuerte como para obligar al banco central común a proporcionar liquidez en momentos de crisis.

[Nota al pie: Esta capacidad de forzar al banco central a actuar como prestamista de última instancia es una característica fundamental de los Estados con soberanía monetaria.]

⸻

Viabilidad en Europa

¿Existe alguna posibilidad de que Europa avance hacia una unión presupuestaria de este tipo? El presupuesto de la Unión Europea representa aproximadamente el 1% del PIB de la UE, mientras que los presupuestos nacionales suelen representar entre el 40% y el 50% del PIB. Existe muy poca probabilidad de que se produzca una centralización significativa de los presupuestos y de la deuda pública a nivel europeo en el futuro cercano.

Tal centralización requeriría un alto grado de unificación política, incluyendo una transferencia sustancial de soberanía nacional en materia fiscal hacia un gobierno y un parlamento europeos. Actualmente no existe voluntad política en Europa para avanzar en esa dirección. Como resultado, los mecanismos de seguro y de protección derivados de la centralización presupuestaria no están disponibles en la unión monetaria europea actual.

⸻

De lo anterior se deduce que una unión monetaria sin unión presupuestaria funciona de manera muy distinta a una unión monetaria acompañada de una unión presupuestaria. La primera puede denominarse una “unión monetaria incompleta”, mientras que la segunda constituye una “unión monetaria plena”. Volveremos sobre esta distinción en el Capítulo 5, donde analizamos distintos tipos de uniones monetarias incompletas, en particular la fragilidad de la Eurozona.

En el Capítulo 7 analizaremos si es posible diseñar instituciones que, aun sin alcanzar una unión presupuestaria y política completa, puedan proporcionar cierto grado de seguro y protección a los Estados miembros de una unión monetaria incompleta, y cómo diseñarlas de manera que se minimice el problema de riesgo moral.

$$\text{\textbf{Mecanismos de seguro (II):} Mercados financieros}$$
Una unión presupuestaria proporciona un mecanismo de seguro dentro de una unión monetaria. Existe otra forma de organizar un esquema de seguro en una unión monetaria. Este esquema opera a través de los mercados financieros. Supongamos, como antes, un shock asimétrico que afecta negativamente a Francia y positivamente a Alemania. Supongamos (y este es un supuesto crucial) que los mercados financieros de Francia y Alemania están completamente integrados.

Concentrémonos aquí en cómo los mercados integrados de bonos y acciones facilitan el ajuste. Como resultado del shock negativo, las empresas francesas incurren en pérdidas, lo que hace caer los precios de sus acciones. Dado que el mercado de renta variable está plenamente integrado, estas acciones francesas también están en manos de residentes alemanes. Por tanto, estos últimos soportan parte del coste de la caída de la actividad económica en Francia. De forma inversa, el auge en Alemania eleva los precios de las acciones de las empresas alemanas. Como estas también están en manos de residentes franceses, estos reciben cierta compensación por las dificultades económicas en Francia. Dicho de otro modo, un mercado de acciones integrado funciona como un sistema de seguro. El riesgo de un shock negativo en un país es compartido por todos los países. Como resultado, el impacto del shock negativo sobre la renta de los residentes del país afectado se ve mitigado.

Un mecanismo similar opera a través del mercado de bonos integrado. Como consecuencia del shock negativo, las empresas francesas sufren pérdidas y algunas quiebran, lo que reduce el valor de los bonos franceses en circulación. Parte de estos bonos está en manos de residentes alemanes, que por tanto también soportan el coste de la recesión en Francia.

La ventaja de este esquema de seguro basado en los mercados financieros privados es que reduce el problema de riesgo moral. Sin embargo, presenta también un inconveniente importante. Los desempleados pobres en Francia, que no poseen activos financieros emitidos en Alemania, obtienen poca compensación de este sistema. En cambio, los ciudadanos franceses con mayor riqueza y carteras de activos más amplias son quienes probablemente reciben la mayor parte de las transferencias. Como resultado, este esquema de seguro privado, en ausencia de un canal público, proporciona una cobertura insuficiente para la mayoría de los ciudadanos franceses.








\end{document}